\section{Zeros, Euler products and prime races under Dirichlet GRH}\label{sec:zer}
\providecommand{\zerJ}{\mathcal{J}}
\providecommand{\zerK}{\mathcal{K}}
\providecommand{\zersinc}{\operatorname{sinc}}
\providecommand{\zergE}{\gamma_{\mathrm{E}}}
\providecommand{\zerLL}{\mathcal{L}}
\providecommand{\zerconv}{\operatorname{conv}}

This section collects the consequences of \cref{src:thm:dgrh} that are statements about the
zeros themselves and about the prime sums they govern: sharp partial Euler products on the
critical line, the distribution and rigidity of signed zero multiplicities, the covariance,
oscillation and extremal values of the smoothed prime errors of the source proof at several
scales, and the sign and size of several weighted prime sums. Every result marked ``\stG'' uses
\cref{src:thm:dgrh} as its only input on the location of zeros; simplicity of zeros and linear
independence of ordinates are never assumed. Nontrivial zeros are written
$\rho=\tfrac12+i\gamma$ and are counted with multiplicity unless stated otherwise, $\Lambda$ is
the von Mangoldt function, and $\zergE$ is Euler's constant.

One elementary estimate is used several times.

\begin{lemma}[Local mean square of an exponential sum; \stU]\label{zer:lem:bins}
Let $0<h\le1$ and let $T(v)=\sum_{\lambda}d_\lambda e^{i\lambda v}$ be a finite sum with real
frequencies. Group the frequencies in the bins $I_j=[j/h,(j+1)/h)$, $j\in\Z$, and put
$b_j=\sum_{\lambda\in I_j}\abs{d_\lambda}$. Then, with an absolute implied constant,
\[
 \sup_{U\in\R}\ \frac1h\int_U^{U+h}\abs{T(v)}^2\,\dd v\ \ll\ \sum_{j\in\Z}b_j^2 .
\]
Consequently, if an infinite family of coefficients has $\sum_jb_j^2<\infty$, its partial sums
over $\abs{\lambda}\le Y$ converge, as $Y\to\infty$, in the norm
$\norm{f}_{S^2}=\sup_b\bigl(\int_b^{b+1}\abs{f(t)}^2\,\dd t\bigr)^{1/2}$.
\end{lemma}
\begin{proof}
The function $W(y)=(\sin(y/2)/(y/2))^2$ is nonnegative, is bounded below by a positive constant
on $[-1/2,1/2]$, and its Fourier transform is bounded and supported on $[-1,1]$. Integrate
$\abs{T}^2$ against $h^{-1}W((v-U-h/2)/h)$. The interaction of two frequencies
$\lambda,\lambda'$ vanishes unless $\abs{\lambda-\lambda'}\le1/h$, so only equal or adjacent
bins contribute, and $2b_jb_{j+1}\le b_j^2+b_{j+1}^2$. The second assertion follows by applying
the bound, with $h=1$, to differences of partial sums.
\end{proof}

\subsection{Paired Euler products at the critical line}\label{zer:sec:euler}
\sources{dirichlet-grh-research; exponent_seven_and_grh_transfers}

For a primitive nonprincipal Dirichlet character $\chi$ put
\[
 P_x(s,\chi)=\prod_{p\le x}\bigl(1-\chi(p)p^{-s}\bigr)^{-1},\qquad u=\log x .
\]
The deep Riemann hypothesis predicts the behavior of this ordinary partial product at a fixed
point of the critical line for every sufficiently large cutoff; for instance, for the character
$\chi_4$ modulo four, with $\beta(s)=L(s,\chi_4)$, it predicts
$P_x(1/2,\chi_4)\to\sqrt2\,\beta(1/2)$. Conrad isolates the fixed-height prime-sum obstruction
and relates it to a prime-sum estimate stronger than the pointwise estimate supplied by GRH
\cite{Conrad2005Partial,Kaneko2022Euler}, and Sheth establishes the predicted fixed-center
asymptotic outside a set of finite logarithmic measure \cite{Sheth2025Euler}. Hybrid
Euler--Hadamard products connect the scale $1/\log x$ near a zero to exponential integrals
\cite{GonekHughesKeating2007Hybrid,BuiKeating2007Mean}. The results below concern sharp
products at linked heights at that scale: two offsets whose phases at the cutoff are opposite
cancel the uncontrolled prime-sum error, for every cutoff.

Fix $s_0=\tfrac12+i\theta$ and let $m=\ord_{s=s_0}L(s,\chi)$, with $m=0$ at a nonzero value.
Define
\begin{gather*}
 A_0=e^{-m\zergE}\frac{L^{(m)}(s_0,\chi)}{m!},\qquad
 \delta=\mathbf 1_{\{\theta=0,\ \chi^2\text{ principal}\}},\\
 \zerJ(a)=\int_0^1\frac{1-e^{-iat}}{t}\,\dd t,\qquad
 \zerK(a)=\frac12\int_1^2\frac{e^{-iat}}{t}\,\dd t .
\end{gather*}
All implied constants may depend on the fixed character and on $\theta$. The logarithm
$\log P_x$ is always the one obtained by summing the absolutely convergent power series of
each local factor; this convention also defines complex linear combinations of logarithms.

\begin{theorem}[Cancellation at the cutoff; \stG]\label{zer:thm:pair}
Let $a_1,\dots,a_r$ be real and $c_1,\dots,c_r$ complex. Suppose
$\sum_jc_je^{-ia_j}=0$ and write $S=\sum_jc_j$. There is a fixed logarithm $C_0$ of $A_0$ such
that, for every real $x$ tending to infinity,
\begin{equation}\label{zer:eq:lin}
 \sum_jc_j\log P_x\bigl(s_0+ia_j/u,\chi\bigr)+mS\log u
 =SC_0+\sum_jc_j\bigl(m\zerJ(a_j)+\delta\zerK(a_j)\bigr)+O\Bigl(\frac1u\Bigr).
\end{equation}
The error is uniform for bounded offsets and bounded $\sum_j\abs{c_j}$. For integer $c_j$,
exponentiation gives an ordinary product or quotient with an unambiguous limiting constant.
\end{theorem}

The proof isolates the endpoint. Write $\Psi(x)=\sum_{n\le x}\chi(n)\Lambda(n)$ and, for $v>0$,
\[
 E(v)=e^{-s_0v}\Psi(e^v),\qquad F_0(v)=E(v)+\frac m{s_0},\qquad
 H(v)=\sum_{2\le n\le e^v}\frac{\chi(n)\Lambda(n)}{n^{s_0}\log n}.
\]
Step functions use right endpoints; changing values at jumps has no effect on their integrals.

\begin{lemma}\label{zer:lem:primitive}
The primitive $\int_1^VF_0(v)\,\dd v$ is bounded for $V\ge1$. There is a number $C_0$ with
$\exp C_0=A_0$ such that
\begin{equation}\label{zer:eq:H}
 H(v)=C_0-m\log v+\frac{F_0(v)}v+O(1/v)\qquad(v\ge1).
\end{equation}
\end{lemma}
\begin{proof}
The explicit formula for a primitive nonprincipal character
\cite{Davenport2000Multiplicative,Koukoulopoulos2019Distribution}, in the sense of
distributions on $v\ge1$, is
\begin{equation}\label{zer:eq:explicit}
 E(v)=-\sum_{\rho=1/2+i\gamma}\frac{e^{i(\gamma-\theta)v}}{\rho}+D(v),
 \qquad D(v)=O_\chi\bigl((1+v)e^{-v/2}\bigr).
\end{equation}
This is Perron's formula for $-L'/L$ after the shift past the nontrivial zeros, which lie on
the critical line by \cref{src:thm:dgrh}. The pole of the Perron kernel at zero, the possible
trivial zero there and the remaining negative trivial zeros give the stated $D$. A smooth test
in $v$ makes the zero sum absolutely convergent. The $m$ terms at $\gamma=\theta$ give
$-m/s_0$. After one integration every remaining coefficient is bounded by
$2/\abs{\rho(\gamma-\theta)}$. Their sum converges by the zero count
$N(T+1)-N(T)\ll_\chi\log(2+T)$ for $T\ge1$, applied to both signs. This proves the
bounded-primitive assertion, also when $\theta$ itself is the ordinate of a multiple zero.

Stieltjes partial summation, starting below $\log2$, gives an endpoint $E(v)/v$ and an
integral with integrand $E(w)(s_0/w+1/w^2)$. The constant mode $-m/s_0$ contributes
$-m\log v$ and a constant, because its $1/v$ terms cancel. The integral of
$F_0(w)(s_0/w+1/w^2)$ to infinity converges, with tail $O(1/v)$ by integration by parts
against its bounded primitive. This gives \eqref{zer:eq:H}, pointwise including jumps, with
some $C_0$.

For $\sigma>0$, partial summation and the GRH estimate $\Psi(x)\ll_\chi\sqrt x\log^2(2x)$ give
\[
 \sigma\int_0^\infty H(v)e^{-\sigma v}\,\dd v=\log L(s_0+\sigma,\chi).
\]
Indeed the logarithmic series agrees with the Euler product in $\Re s>1$ and defines its
continued logarithm in $\Re s>1/2$ by locally uniform convergence. The constant and logarithm
in \eqref{zer:eq:H} contribute $C_0+m\log\sigma+m\zergE$. The other terms have Abel mean
$o(1)$ as $\sigma\downarrow0$: the $O(1/v)$ remainder contributes $O(\sigma\log(2/\sigma))$
for small $\sigma$, and the bounded primitive treats $F_0(v)/v$ by integration by parts.
Comparing the exponential with $L(s_0+\sigma,\chi)\sim\sigma^mL^{(m)}(s_0,\chi)/m!$ proves
the claimed normalization.
\end{proof}

\begin{proof}[Proof of \cref{zer:thm:pair}]
Set $W(t)=\sum_jc_je^{-ia_jt}$. Truncation at prime powers $p^k\le x$ gives exactly
\[
 \int_{0-}^uW(v/u)\,\dd H(v)=W(1)H(u)-u^{-1}\int_0^uH(v)W'(v/u)\,\dd v .
\]
Put $R(v)=H(v)-(C_0-m\log v)$. Its primitive is bounded on $[0,\infty)$. At zero the logarithm
is integrable. For $v\ge1$ the partial summation in the proof of \cref{zer:lem:primitive}
gives the exact identity
\[
 R(v)=F_0(v)/v-\int_v^\infty F_0(w)k(w)\,\dd w,\qquad k(w)=s_0/w+1/w^2 .
\]
The first term has a bounded primitive by \cref{zer:lem:primitive}. Integrating the tail from
$1$ to $V$ gives
$\int_1^V(w-1)k(w)F_0(w)\,\dd w+(V-1)\int_V^\infty k(w)F_0(w)\,\dd w$. The first term is
bounded because $(w-1)k(w)$ is bounded and has integrable derivative. The second uses the
$O(1/V)$ tail bound. Hence
$\int_0^uR(v)W'(v/u)\,\dd v=O(\norm{W'}_\infty+\norm{W''}_\infty)$ by integration by parts.
Substitution of the model $C_0-m\log v$ and integration by parts at zero now give
\begin{equation}\label{zer:eq:endpoint}
 \int_{0-}^uW(v/u)\,\dd H(v)=W(0)(C_0-m\log u)+m\int_0^1\frac{W(0)-W(t)}t\,\dd t
   +W(1)R(u)+O_W(1/u).
\end{equation}
The only uncontrolled pointwise quantity is multiplied by $W(1)$.

The ordinary Euler product also contains terms with $p\le x<p^k$. Terms of order $k\ge3$ have
total absolute value $O_W(x^{-1/6})$ for these finite exponential weights, which are bounded
for all real arguments: for $p>x^{1/3}$ sum $O(p^{-3/2})$; for $p\le x^{1/3}$ sum the
geometric tail starting with its first omitted power, which is $O(x^{-1/2})$ per prime. The
square terms contribute
\[
 \frac12\sum_{\sqrt x<p\le x}\frac{\chi(p)^2}{p^{1+2i\theta}}W(2\log p/u)
 =\frac\delta2\int_1^2\frac{W(t)}t\,\dd t+O_W(1/u).
\]
In detail, fixed-modulus prime number theory gives
\[
 \sum_{p\le y}\chi^2(p)\log p=\mathbf 1_{\chi^2\text{ principal}}\,y+O_\chi(y/\log^3y).
\]
Partial summation on $[\sqrt x,x]$ makes its error $O_{W,\theta}(1/u)$ or smaller. The main
integral before setting $\theta=0$ is $\tfrac12\int_1^2W(t)e^{-i\theta ut}\,\dd t/t$;
integration by parts gives $O_{W,\theta}(1/u)$ for fixed $\theta\ne0$. Using $W(1)=0$ in
\eqref{zer:eq:endpoint} completes the proof.
\end{proof}

\begin{corollary}[\stG]\label{zer:cor:pairs}
In the setting of \cref{zer:thm:pair}, as $x\to\infty$:
\begin{enumerate}[label=\textup{(\alph*)}]
\item Taking offsets $0,\pi$ with both coefficients equal to $1$,
\[
 u^{2m}P_x(s_0,\chi)P_x(s_0+\pi i/u,\chi)\longrightarrow
 A_0^2\exp\bigl\{m\zerJ(\pi)+\delta[\zerK(0)+\zerK(\pi)]\bigr\}.
\]
This product contains the actual fixed-height sharp product.
\item For every fixed integer $k$,
\[
 \frac{P_x(s_0+2\pi ki/u,\chi)}{P_x(s_0,\chi)}\longrightarrow
 \exp\bigl\{m\zerJ(2\pi k)+\delta[\zerK(2\pi k)-\zerK(0)]\bigr\}.
\]
\item For a real character at $\theta=0$,
\[
 u^m\abs[\big]{P_x\bigl(\tfrac12+\pi i/(2u),\chi\bigr)}\longrightarrow
 \abs{A_0}\exp\bigl\{m\Re\zerJ(\pi/2)+\Re\zerK(\pi/2)\bigr\}.
\]
In particular, with no exceptional set of cutoffs,
\begin{equation}\label{zer:eq:beta}
 \lim_{x\to\infty}\abs[\Big]{P_x\Bigl(\tfrac12+\frac{\pi i}{2\log x},\chi_4\Bigr)}
 =\beta(\tfrac12)\exp\Bigl(\tfrac12\int_1^2\frac{\cos(\pi t/2)}t\,\dd t\Bigr)
 =0.547115431519\ldots .
\end{equation}
\end{enumerate}
\end{corollary}
\begin{proof}
Parts (a) and (b) are \cref{zer:thm:pair} with offsets $0,\pi$ and coefficients $1,1$, and
with offsets $2\pi k,0$ and coefficients $1,-1$. For a real character at $\theta=0$ the
products at opposite offsets are conjugates, and offsets $\pm\pi/2$ with coefficients $1,1$
give (c). The alternating series for $\beta(1/2)$ has positive sum, so $m=0$ in
\eqref{zer:eq:beta}.
\end{proof}

These assertions apply to the original unsmoothed factors at all primes $p\le x$; their
logarithmic combinations implement a weight that vanishes exactly at the endpoint. The shift
in \eqref{zer:eq:beta} tends to zero at the scale on which the largest prime factors rotate
through a fixed angle, and it changes the prime-square contribution as well as the endpoint
error: at height zero the prediction for the same sharp cutoff is
$\sqrt2\,\beta(1/2)=0.944258314238\ldots$, where $\beta(1/2)=0.667691457190\ldots$. A
numerical illustration at $130$ cutoffs up to $10^7$ is in
\artifact{dirichlet-grh-research/critical-line/euler_example.csv}; at $x=10^7$ the two finite
products are $0.94591$ and $0.54823$.

\paragraph{Individual products.}
The same calculation identifies what remains when the endpoint is retained. For real $a$ let
\[
 B(v,a)=\log P_{e^v}(s_0+ia/v,\chi)+m\log v-C_0-m\zerJ(a)-\delta\zerK(a).
\]
Equations \eqref{zer:eq:endpoint} and \eqref{zer:eq:H} give, uniformly for $\abs a\le A$,
\begin{equation}\label{zer:eq:single}
 B(v,a)=e^{-ia}\frac{F_0(v)}v+O_A(1/v).
\end{equation}
Thus the pointwise convergence of the canonical logarithm at the fixed height is equivalent to
$\Psi(x)=o(\sqrt x\log x)$ for this character. The GRH bound $O(\sqrt x\log^2x)$ leaves one
logarithm to save and requires a little-oh at the target scale; the cancellation in
\cref{zer:thm:pair} removes this entire term for the paired observable.

\begin{proposition}[Prescribed subsequences at fixed height; \stG]\label{zer:prop:subseq}
Along the following explicit sequences of cutoffs the endpoint term is absent.
\begin{enumerate}[label=\textup{(\alph*)}]
\item Let $\chi$ be real and $\theta\ne0$. Then, also when $s_0$ is a zero,
\[
 (\log X_n)^m\abs{P_{X_n}(s_0,\chi)}=\abs{A_0}\bigl(1+O(1/\log X_n)\bigr),\qquad
 X_n=\exp\bigl((\pi/2+n\pi)/\abs\theta\bigr)\quad(n\ge0).
\]
\item Let $\chi$ be arbitrary, let $\theta_1\ne\theta_2$, write $s_j=\tfrac12+i\theta_j$ and
$m_j,A_j,\delta_j$ for the corresponding constants, and put
$v_n=2\pi n/\abs{\theta_1-\theta_2}$, $x_n=\exp(v_n)$. Then
\[
 v_n^{m_1-m_2}\frac{P_{x_n}(s_1,\chi)}{P_{x_n}(s_2,\chi)}
 =\frac{A_1}{A_2}\,2^{(\delta_1-\delta_2)/2}\bigl(1+O(1/v_n)\bigr).
\]
\end{enumerate}
\end{proposition}
\begin{proof}
(a) Use \eqref{zer:eq:single} at $a=0$. Since $\Psi(X_n)$ is real and
$\exp(-i\theta\log X_n)$ is purely imaginary, taking real parts removes the term
$E(v)/v$; the remaining constant $m/(s_0v)$ is $O(1/v)$.
(b) Use \eqref{zer:eq:single} at $a=0$ separately at the two heights. The two uncontrolled
terms contain the same $\Psi(x_n)$ with identical phase $\exp(-i\theta_jv_n)$, so they cancel
on subtraction; $\zerK(0)=\tfrac12\log2$ gives the power of two.
\end{proof}

\begin{theorem}[Local mean square; \stG]\label{zer:thm:ms}
For $U\ge2$, $0<h\le1$ and fixed $A<\infty$,
\[
 \Bigl(\frac1h\int_U^{U+h}\sup_{\abs a\le A}\abs{B(v,a)}^2\,\dd v\Bigr)^{1/2}
 \ll_{\chi,\theta,A}\frac{\log^2(2/h)}U .
\]
In particular, if $\log(1/h)=o(\sqrt U)$, the normalized individual products converge to
their respective profiles $A_0\exp\{m\zerJ(a)+\delta\zerK(a)\}$ in relative logarithmic
measure in every such interval, uniformly in bounded $a$.
\end{theorem}
\begin{proof}
Apply \cref{zer:lem:bins} to \eqref{zer:eq:explicit}, dropping the common unit phase
$e^{-i\theta v}$. The sum of absolute coefficients for $\abs\gamma\le1/h$ is
$O_\chi(\log^2(2/h))$ by summing unit bins. For the positive tail, each bin $j\ge1$ of length
$1/h$ contains $O_\chi(h^{-1}\log(2+(j+1)/h))$ ordinates, so its total absolute coefficient
is at most $C_\chi\log(2+(j+1)/h)/j$. The sum of the squares of these bin bounds is
$O_\chi(\log^2(2/h))$, and the negative tail obeys the same estimate. This also gives local
$L^2$ convergence of the high-frequency series, identifying its limit with the distribution
in \eqref{zer:eq:explicit}. The normalized $L^2$ norm of $F_0$ on the interval is
consequently $O_{\chi,\theta}(\log^2(2/h))$, and \eqref{zer:eq:single} proves the stated
bound. For each fixed tolerance, Chebyshev's inequality bounds the exceptional fraction by
the square of the right-hand side divided by the tolerance squared; continuity of
exponentiation near zero transfers this to the relative product error.
\end{proof}

For instance $h=\exp(-U^{1/3})$ is admissible: the corresponding interval of prime cutoffs
starts at $X=e^U$ and has length asymptotic to $X\exp[-(\log X)^{1/3}]$, and the theorem
supplies a vanishing exceptional fraction inside every such interval. A qualitative statement
that the global exceptional set has finite logarithmic measure does not give this local
conclusion. At the boundary scale $h=\exp(-a\sqrt U)$ the mean-square bound with Chebyshev's
inequality gives only $O_\chi(a^4/\eps^2)$ for the exceptional fraction at tolerance $\eps$,
without decay in $U$. At the central point a frequency split matched to the threshold gives
decay there.

\begin{theorem}[Exceptional cutoffs for the central product; \stG]\label{zer:thm:exceptional}
Let $\theta=0$, so that $m=\ord_{1/2}L(s,\chi)$ and $\delta=1$ exactly when $\chi^2$ is
principal, and write
\[
 B(v)=B(v,0)=\log P_{e^v}(\tfrac12,\chi)+m\log v-C_\chi,\qquad
 \exp C_\chi=2^{\delta/2}e^{-m\zergE}\frac{L^{(m)}(1/2,\chi)}{m!},
\]
where $C_\chi=C_0+\tfrac\delta2\log2$. For each fixed $\eps>0$ and every $U\ge2$,
\[
 \abs[\big]{\set{v\in[U,U+1]:\abs{B(v)}>\eps}}\ll_{\chi,\eps}\exp\bigl(-\sqrt{2\pi\eps U}\bigr).
\]
Thus the set of exceptional $x\ge e^2$ has finite measure for
$e^{b\sqrt{\log x}}\,\dd x/x$ whenever $b<\sqrt{2\pi\eps}$. For $a>0$ and
$h=\exp(-a\sqrt U)$, the fraction of $v\in[U,U+h]$ with $\abs{B(v)}>\eps$ is
\[
 O_{\chi,\eps}\bigl(\exp[-(\sqrt{2\pi\eps}-a)\sqrt U]\bigr).
\]
\end{theorem}

At the boundary window scale $h=\exp(-a\sqrt U)$ the exceptional fraction therefore vanishes
for every fixed tolerance $\eps>a^2/(2\pi)$; for example $a=1$ allows every tolerance greater
than $1/(2\pi)$ in the canonical logarithm. This does not give convergence for every
tolerance in each such window at fixed $a$, or a pointwise statement. The estimates also
bound relative product discrepancies: $\abs B\le\eps$ implies
$\abs{(\log x)^mP_xe^{-C_\chi}-1}\le e^\eps-1$ at $x=e^v$.

\begin{lemma}\label{zer:lem:tail}
Let $\theta=0$, so that $E(v)=e^{-v/2}\Psi(e^v)$. For every real $\lambda\ge1$ and every
$U\ge2$,
\[
 \abs[\big]{\set{v\in[U,U+1]:\abs{E(v)}>\lambda}}\ll_\chi e^{-\sqrt{2\pi\lambda}} .
\]
\end{lemma}
\begin{proof}
Count both signs of the ordinate with multiplicity. The Dirichlet zero count
\cite{BennettMartinOBryantRechnitzer2021Counting} gives $N(T)=(T/\pi)\log T+O_\chi(T)$ and
$N(t+1)-N(t)\ll_\chi\log(2+t)$. Partial summation gives, for $T\ge2$,
\[
 \sum_{\abs\gamma\le T}\abs\rho^{-1}\le\frac{(\log T)^2}{2\pi}+O_\chi(\log(2T)).
\]
Indeed one can replace $1/\abs\rho$ by $1/\abs\gamma$ above height one at bounded total cost
and integrate $N(t)/t^2$. The smooth remainder $D$ in \eqref{zer:eq:explicit} is uniformly
bounded on $v\ge2$.

For the tail $S_T(v)=\sum_{\abs\gamma>T}e^{i\gamma v}/\rho$,
\[
 \int_U^{U+1}\abs{S_T(v)}^2\,\dd v\ll_\chi\frac{\log^2(2T)}T .
\]
This is \cref{zer:lem:bins} with $h=1$: the sum $b_j$ of the absolute coefficients in the bin
$[j,j+1)$ satisfies $b_j\ll_\chi\log(2+\abs j)/(1+\abs j)$ in the tail, which yields the
estimate by summing squares. The same bound makes the infinite series converge in local
$L^2$, and its distributional limit is the tail of the explicit formula.

For large $\lambda$ choose $\log T=\sqrt{2\pi\lambda}-D_\chi$, where the fixed $D_\chi$ is
large enough that the low-ordinate sum plus the smooth remainder is at most
$\lambda-\sqrt\lambda$; expanding the square in the low-ordinate bound shows that such a
choice exists. Chebyshev's inequality then bounds the exceptional measure by
$O_\chi(\log^2(2T)/(\lambda T))=O_\chi(e^{-\sqrt{2\pi\lambda}})$. Enlarge the constant for
bounded $\lambda$.
\end{proof}

\begin{proof}[Proof of \cref{zer:thm:exceptional}]
At $\theta=0$ and $a=0$ equation \eqref{zer:eq:single} reads $B(v)=E(v)/v+O_\chi(1/v)$ for
$v\ge2$, since $F_0=E+2m$, $\zerJ(0)=0$ and $\zerK(0)=\tfrac12\log2$. Apply
\cref{zer:lem:tail} at $\lambda=\eps U-O_\chi(1)$. This proves the first assertion. Sum over
integer $U$ with weight $e^{b\sqrt{U+1}}$ for the second, and divide the unit-interval bound
by $h$ for the third.
\end{proof}

No optimality is claimed for the constant $2\pi$; it comes from the two-sided zero count and
the cutoff used in the proof.

\paragraph{The signed zero spectrum.}
The Mellin residues of the smoothed prime remainder in the source proof are signed
differences of zero multiplicities. How these signed data determine the modulus and the
twist, in the limit and from a finite band or a finite initial segment of zeros, and the
resulting criterion for four central values at once, are treated in \cref{sec:rec}.

\subsection{Signed multiplicity differences in short intervals and in phase boxes}\label{zer:sec:short}
{\emergencystretch=6em\sources{grh-methods-polynomial-density; grh_cm_four_terms_and_prime_bias; grh_spectral_rigidity; grh_filtered_prime_races; grh_prime_error_balancing}}

The next two theorems locate both signs of differences of zero multiplicities on the
critical line: in every interval of triple-logarithmic length for an arbitrary fixed finite
family, and in every prescribed box of prime phases for the twist-family difference of the
source proof. A third statement records that the multiplicity vectors of all characters to a
fixed modulus have full rank in windows of logarithmic length, and a fourth that no nonzero
polynomially weighted combination of the multiplicity divisors of distinct primitive
characters has finite support.

\paragraph{Intervals of triple-logarithmic length.}
Let $\chi_1,\dots,\chi_r$ be distinct primitive Dirichlet characters, of conductors
$q_1,\dots,q_r$, and let $m_j(t)$ be the order of $L(s,\chi_j)$ at $s=\tfrac12+it$. The
conductor-one character gives the Riemann zeta function. Fix $r\ge2$. For $c\in\R^r$ put
$z_c(t)=\sum_jc_jm_j(t)$ and $z_{c,\pm}=\max(\pm z_c,0)$. Write
$\log_3T=\log\log\log T$. Sums below count each ordinate once.

\begin{theorem}[\stG]\label{zer:thm:short}
Fix $\eps>0$ and the characters. For all sufficiently large $T$, every
\[
 (2/\pi+\eps)\log_3T\le H\le T
\]
satisfies, simultaneously for both signs and every $c$ with $\sum_jc_j=0$ and
$\norm c_2=1$,
\[
 \sum_{T<t\le T+H}z_{c,\pm}(t)\ \gg_{\eps,\chi_1,\dots,\chi_r}\ H\log\log H .
\]
\end{theorem}

For two characters take $c=(1,-1)/\sqrt2$: each such interval contains both an ordinate
where the first multiplicity is greater and one where the second is greater. If the
multiplicities at a common zero are $2$ and $1$, the sum counts one unit there. Srinivas
proved that for distinct functions in the Selberg class of equal degree each direction of
excess occurs within distance $O(\log\log T)$ of every large $T$ \cite{Srinivas2002Distinct};
the use of real critical-line ordinates gives a smaller interval for fixed Dirichlet
families and a lower bound for the mass in it. Banks gives bounds uniform in the character
modulus \cite{Banks2023Comparing}; Fujii and Bombieri--Perelli obtain much larger total
multiplicity counts in global ranges \cite{Fujii1976Zeros,BombieriPerelli1998Distinct}, and
Banks and Loftus detect genuinely unshared zeros in longer windows
\cite{BanksLoftus2026Unshared}.

\begin{proof}
\emph{A positive weight and its prime frequencies.} Let $Q$ be the least common multiple of
the conductors and $C(n)=\sum_jc_j\chi_j(n)$. The induced characters of $(\Z/Q\Z)^\times$
are distinct by uniqueness of the primitive inducer, and character orthogonality gives the
mean square of $C(v)$ over the group as $1$. Choose $v$ with $\abs{C(v)}\ge1$ and $\theta$
with $e^{i\theta}C(v)=\abs{C(v)}$. Put $Y=\lambda\log H$ for a sufficiently small fixed
$\lambda>0$. Take all primes $p\le Y$ congruent to $v$ modulo $Q$ and write
\[
 w_p=\frac{\log p}{\sqrt p},\qquad V^2=\sum_pw_p^2\sim\frac{(\log Y)^2}{2\varphi(Q)},
\]
the asymptotic following by partial summation from the prime number theorem in the finitely
many fixed progressions. Form the real trigonometric sums
\[
 A(t)=\Im\prod_p\bigl(1+i(w_p/V)\cos(t\log p+\theta)\bigr),\qquad
 P_\pm(t)=1\pm e^{-1/2}A(t).
\]
On the real axis $\abs A\le\prod_p\exp(w_p^2/(2V^2))=\sqrt e$, so $0\le P_\pm\le2$. Extend
their finite exponential expansions to complex $t$, and write either of the two as
$P(t)=\sum_Rb_Re^{it\log R}$, where the positive rationals $R$ are in lowest terms. Unique
factorization gives $b_1=1$ and $b_p=\pm e^{-1/2}e^{i\theta}w_p/(2V)$ at the chosen primes.
Every other nonconstant frequency, apart from their reciprocals, has at least two different
primes with nonzero exponents and cannot be a prime power or its reciprocal. Also
$b_{1/R}=\overline{b_R}$. Both the numerator and the denominator of $R$ are at most
$D=\prod_{p\le Y}p\le e^{2Y}$ for large $Y$. Expanding the product and using
Cauchy--Schwarz gives $\sum_R\abs{b_R}\le1+\prod_p(1+w_p/V)\le1+\exp(\sqrt{\pi(Y)})$.

\emph{An explicit-formula test localized in height.} Choose $0<a<\pi$ and $0<\kappa<1/2$
such that $a(1/2-\kappa)(2/\pi+\eps)>1$. Set $B=\kappa H$ and $X=T+H/2$. Let $k(t)$ be
$e^{-\cosh(at)}$ divided by its real integral, and put
\[
 K_B(t)=\frac1{2B}\int_{-B}^Bk(t-y)\,\dd y,\qquad h(t)=P(t)K_B(t-X).
\]
Thus $0\le h(t)\le1/B$ on $\R$. Use the convention
$\widehat h(u)=\int_\R h(t)e^{-itu}\,\dd t$. For fixed $b$ with $1/2<b<\pi/(2a)$, shifting
the contour to $\Im t=\pm b$ gives $\abs{\widehat k(u)}\ll_{a,b}e^{-b\abs u}$, since
$\abs{e^{-\cosh(a(t+ib))}}=e^{-\cos(ab)\cosh(at)}$. Consequently
\[
 \widehat h(u)=\sum_Rb_Re^{-iX(u-\log R)}\,\widehat k(u-\log R)\,
 \zersinc\bigl(B(u-\log R)\bigr),
\]
with value $1$ of the last two factors at $u=\log R$: resonant frequencies $u=\log R$
contribute exactly $b_R$, independently of the center $X$.

For a primitive $\chi_j$ write $a_j=(1-\chi_j(-1))/2$ and $\psi=\Gamma'/\Gamma$. Under
Dirichlet GRH the Guinand--Weil formula in this normalization reads
\begin{align*}
 2\pi\sum_tm_j(t)h(t)={}&2\pi\,\mathbf 1_{q_j=1}\set{h(i/2)+h(-i/2)}\\
 &+\int_\R h(t)\Bigl[\log(q_j/\pi)+\Re\psi\bigl((1/2+a_j+it)/2\bigr)\Bigr]\dd t\\
 &-2\Re\sum_{n\ge2}\frac{\Lambda(n)\chi_j(n)}{\sqrt n}\,\widehat h(\log n).
\end{align*}
The real-valued test need not be even. The formula follows by integrating the completed
logarithmic derivative on the two sides of the critical strip, using its functional
equation on the left; the strip of decay just proved justifies both contour shifts and
absolute summation on the prime side. For the convention see
\cite[\S5]{FourierOptimizationQuadratic}.

After multiplication by $c_j$ and summation, the archimedean contribution is $O(1)$
uniformly: the common $\log(\abs t/2)$ from Stirling's formula cancels since
$\sum_jc_j=0$, the combined density is bounded near zero, and $\int\abs h\le2$. The pole
contribution is $o(1)$: on $\Im t=\pm1/2$ one has $\abs{P(t)}\le D^{1/2}\sum_R\abs{b_R}$,
whereas each translated kernel in $h(\pm i/2)$ has real displacement at least $X-B\ge T$.

On the prime side the total of the off-diagonal terms is $O(e^{CY}/B)$ for a fixed $C$. For
$n\le2D$ and $n\ne R$, the numerator of $n-R$ in the common denominator is a nonzero
integer, hence $\abs{\log n-\log R}\ge1/(2D^2)$. The sinc factor costs at most $2D^2/B$, and
$\sum_{n\le2D}(\log n)/\sqrt n\ll D^{1/2}\log(2D)$. For $n>2D$ the logarithmic difference
exceeds $\log2$ and the Fourier bound gives $O(B^{-1}D^bn^{-b})$; the resulting sum
converges since $b>1/2$. Summing these bounds with $\abs{b_R}$ proves the assertion,
uniformly in the phases and in the vector $c$. Choose $\lambda$ so that $C\lambda<1/2$. The
prime diagonal now gives
\[
 2\pi\sum_tz_c(t)P_\pm(t)K_B(t-X)=\mp e^{-1/2}\abs{C(v)}\,V+O(1).
\]
In particular the two positive test functions have signed sums of opposite signs and size
$\gg\log\log H$.

\emph{Removing the tails.} The number of zeros of any fixed primitive function in a unit
interval at height $y$, with multiplicity, is $O(\log(2+\abs y))$
\cite{BennettMartinOBryantRechnitzer2021Counting}. Put $L=H/2-B$. In the $j$th unit
interval outside either end of $(T,T+H]$, the maximal kernel value is
$O(\exp(-c_ae^{a(L+j)}))$, $j\ge0$; this bound also covers a zero exactly at the excluded
endpoint $T$. Summing the unit-interval bounds, using $H\le T$, gives a tail of size
\[
 O\bigl(\log(2+T)\exp(-c_a'e^{aL})\bigr)=o(1),
\]
because $e^{aL}\ge(\log\log T)^{1+\eta}$ for some $\eta>0$. The same estimates apply to
negative ordinates. Since $0\le h\le1/B$, each sign inside the interval has mass $\gg BV$.
\end{proof}

The constant $2/\pi$ is a sufficient limiting coefficient in this construction; its
optimality is open.

\begin{example}\label{zer:ex:seven}
For $\zeta$ and the Dirichlet beta function $L(s,\chi_{-4})$, use the unnormalized
difference of their multiplicities. The unamplified conductor term is $-\log4$. Using
$P=1\pm\cos(t\log7)$, the two limiting signed readings after multiplication by $2\pi$ are
$-\log4\mp2\log7/\sqrt7$. The absolute prime term exceeds $\log4$, so this one prime
already certifies both signs for sufficiently large windows of the stated length. The
growing family of primes in the proof supplies the factor $\log\log H$.
\end{example}

\paragraph{Prescribed phase boxes.}
Fix $q\ge1$ and an odd prime $d\nmid q$. Put $h=\varphi(q)$ and $\eta=(\cdot/d)$. For each
character $\chi\bmod q$ use its primitive inducing character $\chi^*$, of conductor
$f_\chi$; the character $\chi^*\eta$ is primitive of conductor $df_\chi$. At positive height
put
\[
 z(\gamma)=\sum_{\chi\bmod q}\Bigl(\ord_{1/2+i\gamma}L(s,\chi^*)
   -\ord_{1/2+i\gamma}L(s,\chi^*\eta)\Bigr),\qquad z_\pm=\max(\pm z,0).
\]
Omitted Euler factors have their zeros on $\Re s=0$, so the same integer is obtained with
the imprimitive functions; for $8\mid q$ and $d\equiv7\pmod8$ it is the residue numerator of
the source, studied as a whole in \cref{sec:rec}. Write $\mathbb T=\R/(2\pi\Z)$. For a
finite set $\mathcal B$ of primes put
$\theta_{\mathcal B}(t)=(t\log b\bmod2\pi)_{b\in\mathcal B}$.

\begin{theorem}[\stG]\label{zer:thm:phasebox}
Let $U$ be a nonempty open box of $\mathbb T^{\mathcal B}$, and fix $0<\alpha\le1$. There
are positive constants $c,T_0$ depending only on $q,d,\mathcal B,U,\alpha$ such that,
simultaneously for both signs, for $T\ge T_0$ and $T^\alpha\le H\le T$,
\begin{equation}\label{zer:eq:phase}
 \sum_{\substack{T<\gamma\le T+H\\\theta_{\mathcal B}(\gamma)\in U}}z_\pm(\gamma)\ \ge\
 cH\log\log T .
\end{equation}
Each distinct ordinate is summed once.
\end{theorem}

The left side counts excess multiplicity after all coincidences, including coincidences
across different characters, have been canceled. The case $\mathcal B=\varnothing$ gives
the unrestricted statement. For $q=1$ these are respectively the zeros and the poles, with
multiplicities, of $\zeta(s)/L(s,\eta)$.

\begin{lemma}[A uniform rational-frequency estimate; \stG]\label{zer:lem:rational}
For a primitive character $\xi$ of fixed conductor (including the zeta case), put
$D_\xi(x)=\sum_{T<\gamma_\xi\le T+H}x^{i\gamma_\xi}$. If $x=a/b>1$ is in lowest terms with
$a,b\le Q$, $Q\ge2$, then
\begin{equation}\label{zer:eq:rationallandau}
 D_\xi(x)=-\frac H{2\pi}\frac{\Lambda(x)\xi(x)}{\sqrt x}+O_\xi\bigl(Q^4\log^2(2T)\bigr),
\end{equation}
where the numerator in the main term is zero unless $x$ is an integral prime power. The
value at $x^{-1}$ is the conjugate. Moreover, in the notation above,
\begin{equation}\label{zer:eq:conductordifference}
 \sum_{T<\gamma\le T+H}z(\gamma)=-\frac{hH}{2\pi}\log d+O_{q,d}(\log(2T)).
\end{equation}
\end{lemma}
\begin{proof}
Banks proves a sharper uniform estimate for $x\ge2$ \cite[Theorem~1.2]{Banks2023Comparing};
the coarser estimate is given for the full range $x>1$ because some ratios of products of
primes approach $1$. Integrate $x^sL'(s,\xi)/L(s,\xi)$, divided by $2\pi i$, on the
rectangle with vertical sides $-1,2$. Move each horizontal side by at most $1$ to a height
at distance $\gg_\xi1/\log(2T)$ from all nearby zero ordinates; the standard local zero
count $N_\xi(t+1)-N_\xi(t-1)\ll_\xi\log(2t)$ permits this choice. The local
partial-fraction expansion for $L'/L$ then gives $O_\xi(\log^2(2T))$ on the two horizontal
sides \cite[Lemma~2.2]{Banks2023Comparing}. The horizontal integrals cost
$O_\xi(Q^2\log^2(2T))$. Restoring the endpoints costs $O_\xi(\sqrt x\log(2T))$ in the sum
of $x^\rho$ by GRH. On the right side expand $L'/L=-\sum_{n\ge2}\Lambda(n)\xi(n)n^{-s}$.
The diagonal $n=x$ gives the displayed main term before division by $\sqrt x$. For $n\ne x$
the oscillatory integral of $(x/n)^{it}$ is $O(1/\abs{\log(x/n)})$. For $x/2<n<2x$ we have
$\abs{x-n}\ge1/Q$ and $\abs{\log(x/n)}\gg Q^{-2}$; elsewhere the reciprocal is bounded.
Since $\sum_{n\ge2}\Lambda(n)/n^2<\infty$, the total off-diagonal cost is $O(Q^4)$. On the
left side the functional equation and Stirling's formula give a logarithmic term
$-\log(f_\xi t/(2\pi))$, a uniformly $O_\xi(1/t)$ remainder, and a logarithmic derivative
at $2-it$. Integration by parts bounds the logarithmic term by $O_\xi(\log(2T)/\log x)$,
with $1/\log x\ll Q$. The absolutely convergent Dirichlet series at $2-it$ has frequencies
$\log(nx)\ge\log2$, so costs $O(1)$. These estimates prove \eqref{zer:eq:rationallandau}
after division by $\sqrt x$. At frequency $1$ use the Riemann--von Mangoldt formula:
subtracting the equal numbers of degree-one functions cancels the common $t\log t$ term and
gives \eqref{zer:eq:conductordifference}.
\end{proof}

\begin{proof}[Proof of \cref{zer:thm:phasebox}]
\emph{A bounded test with a growing prime contribution.} Let $Y=c_0\log T$ for a
sufficiently small fixed $c_0>0$ to be specified by the error bound below. Let $S$ consist
of the primes $p\le Y$ outside $\mathcal B$ satisfying $p\equiv1\pmod q$ and $\eta(p)=-1$.
Define $w_p=(\log p)/\sqrt p$ and $\sigma^2=\sum_{p\in S}w_p^2$. This union of reduced
classes modulo $qd$ has density $1/(2h)$ among primes, so the prime number theorem in fixed
progressions and partial summation give $\sigma^2\sim(\log Y)^2/(4h)$. For large $T$ define
the real trigonometric polynomial
\[
 A(t)=\Im\prod_{p\in S}\Bigl(1+i\frac{w_p}\sigma\cos(t\log p)\Bigr).
\]
Pointwise $\abs A\le\exp(\sum_pw_p^2/(2\sigma^2))=\sqrt e$. Its constant term is zero. The
coefficient of $\cos(t\log p)$ arising from the singleton $p$ is $w_p/\sigma$. All
remaining nonzero terms involve at least two distinct primes, each with exponent $1$ or
$-1$ in the frequency ratio, so neither this ratio nor its inverse is an integral prime
power. At a singleton $p$ we have $\sum_\chi\chi^*(p)(1-\eta(p))=2h$. Taking the real part
in \eqref{zer:eq:rationallandau} consequently gives the weighted singleton contribution
$-hH\sigma/\pi$.

\emph{Restricting the phases.} Take a smooth function $0\le\Phi\le1$ supported inside $U$
with normalized Haar integral $\mu>0$ (for empty $\mathcal B$ take $\Phi=1$). Write
$\Phi_M$ for its rectangular Fourier truncation to exponents $\abs{k_b}\le M$, where
$M=\floor{\sqrt{\log T}}$. Smoothness gives
$\norm{\Phi-\Phi_M}_\infty\ll_{\Phi,N}M^{-N}$ for every fixed $N$ and a uniformly bounded
sum of absolute Fourier coefficients. The sum of $\abs z$ in the window is
$O_{q,d}(H\log T+\log T)$ by the zero count. Thus $\Phi_M$ can replace $\Phi$ in sums
against $z$ and $zA$ with an error $o(H)$; here $\Phi$ in a sum over ordinates stands for
$\Phi(\theta_{\mathcal B}(\gamma))$.

In $\Phi_MA$ every frequency with a nonzero exponent in $\mathcal B$ also has a nonzero
exponent in $S$. Unique factorization makes each such frequency, and its reciprocal,
ineligible for a prime-power main term. Only the constant Fourier coefficient $\mu$ of
$\Phi_M$ couples to a singleton in $A$. For $\Phi_M$ alone the constant frequency
contributes $O_{q,d}(H+\log T)$ by \eqref{zer:eq:conductordifference}. Its nonconstant
prime-power terms contribute $O_{\Phi,q,d,\mathcal B}(H)$, since the factor
$\log b/b^{\abs k/2}$ is bounded for the fixed primes $b\in\mathcal B$ and the Fourier
coefficients are absolutely summable.

All the rational frequencies occurring here, in lowest terms, have numerator and
denominator at most
\[
 Q=\exp\Bigl(\sum_{p\le Y}\log p+M\sum_{b\in\mathcal B}\log b\Bigr)
 \le\exp\bigl(2Y+O_{\mathcal B}(M)\bigr)
\]
for large $T$. The sum of the absolute Fourier coefficients of $A$ is at most
$\prod_{p\in S}(1+w_p/\sigma)\le\exp(\sqrt{\abs S})$ by Cauchy--Schwarz. Using
\eqref{zer:eq:rationallandau} termwise therefore costs
\[
 O_{\Phi,q,d}\bigl(Q^4\log^2(2T)\exp(\sqrt{\abs S})\bigr).
\]
Choose, for example, $c_0=\alpha/32$; this error is at most $T^{\alpha/4+o(1)}=o(H)$. We have proved the two estimates
\[
 \sum_{T<\gamma\le T+H}z(\gamma)\Phi=O(H),\qquad
 \sum_{T<\gamma\le T+H}z(\gamma)\Phi A(\gamma)=-\frac{\mu h}\pi H\sigma+o(H).
\]
The functions $\Phi(\sqrt e-A)$ and $\Phi(\sqrt e+A)$ take their values in
$[0,2\sqrt e]$ and vanish outside the box. Their signed sums are, respectively,
$+(\mu h/\pi)H\sigma+O(H)$ and $-(\mu h/\pi)H\sigma+O(H)$. Dividing the first by $2\sqrt e$
bounds the positive mass from below, and negating the second does the same for the
negative mass. Since $\sigma\asymp\log\log T$, \eqref{zer:eq:phase} follows.
\end{proof}

Liu and Zaharescu study smooth tests on multidimensional tori in races between the zeros of
two Dirichlet $L$-functions \cite{LiuZaharescu2021Races}. In \cref{zer:thm:phasebox} the
cancellation is between two whole products, the angular box is fixed arbitrarily before the
amplifier is chosen, and the same height window works in both directions. The count allows
different multiplicities at a common zero; simplicity, disjointness of zero sets, and the
conjectural $H\log T$ scale for almost all excess multiplicity are stronger assertions.

\paragraph{Full rank of the multiplicity vectors.}
The next statement is used in \cref{zer:sec:balancing} and \cref{zer:sec:cov}.

\begin{corollary}[Multiplicity rank; \stU]\label{zer:thm:rank}
Fix $q\ge1$ and put $n_q=\varphi(q)$. There is $K_q$ such that if $T$ is sufficiently large
and $K_q\log T\le H\le T$, the vectors
\[
 M(\rho)=\bigl(\ord_\rho L(s,\chi^*)\bigr)_{\chi\bmod q}\qquad
 (0<\Re\rho<1,\ T<\Im\rho<T+H)
\]
span $\R^{n_q}$. Distinct zeros index the vectors; a zero absent from one $L$-function
contributes a zero entry. Under \cref{src:thm:dgrh} the same vectors are indexed by
distinct positive ordinates.
\end{corollary}
\begin{proof}
For large $T$ the window $(T,T+H)$ contains a window of the much shorter length allowed in
\cref{rec:thm:localrank}, around its midpoint, so the vectors $M(\rho)$ span over $\C$, and
being real they span $\R^{n_q}$.
\end{proof}

This yields a nonsingular integer multiplicity matrix of size $\varphi(q)$ in every such
height window, including when different characters share zeros. In particular one can
choose $\gg_qT/\log T$ pairwise disjoint bases with ordinates in $[T,2T]$, by dividing it
into intervals of length $2K_q\log(2T)$. The rank argument keeps $q$ fixed; nonsingularity
places no condition on rational linear relations among the heights themselves.

\paragraph{Polynomially weighted multiplicity divisors.}
The rank theorem localizes in height but uses constant weights. The next statement allows
weights that are polynomials in the zero, characters of different conductors, and an
arbitrary finite exceptional set; it is the algebraic input of the covariance theorems of
\cref{zer:sec:cov}. Let $\chi_1,\dots,\chi_r$ be distinct primitive Dirichlet characters, of
conductors $q_j$, including the conductor-one character if desired. Write $a_j\in\set{0,1}$ for
the parity, $\chi_j(-1)=(-1)^{a_j}$, and $\delta_j=1$ for the conductor-one character and $0$
otherwise. The completed functions
\[
 X_j(s)=[s(s-1)]^{\delta_j}(q_j/\pi)^{(s+a_j)/2}\Gamma\bigl((s+a_j)/2\bigr)L(s,\chi_j)
\]
are entire of order one, do not vanish at $0$ and $1$, and their zeros are precisely the
nontrivial zeros. Put $m_j(\rho)=\ord_{s=\rho}X_j(s)$, with $m_j(\rho)=0$ when $X_j(\rho)\ne0$.
By Hadamard factorization,
\begin{equation}\label{zer:eq:hadamard}
 \frac{X_j'}{X_j}(s)=B_j+\sum_{\rho}m_j(\rho)\Bigl(\frac1{s-\rho}+\frac1\rho\Bigr),\qquad
 \sum_{\rho}\frac{m_j(\rho)}{\abs\rho^2}<\infty,
\end{equation}
with locally normal convergence away from the zeros and no assumption of simple zeros.

\begin{theorem}[Polynomial zero-divisor rigidity; \stU]\label{zer:thm:polyrigid}
If polynomials $P_1,\dots,P_r\in\C[s]$ satisfy
\[
 \sum_{j=1}^rP_j(\rho)m_j(\rho)=0
\]
at every point $\rho$ of the union of the zero sets of the $X_j$ except finitely many, then
$P_1=\dots=P_r=0$.
\end{theorem}

\begin{proof}
Let $d\ge0$ bound the degrees and put
\[
 F(s)=\sum_jP_j(s)\frac{X_j'}{X_j}(s),\qquad A(s)=\sum_jP_j(s)\frac{L'}{L}(s,\chi_j).
\]
First, $F$ is a rational function. Truncate the sums in \eqref{zer:eq:hadamard} to a common
disk $\abs\rho\le R$ containing the exceptional points. For each $\rho$,
$P_j(s)/(s-\rho)=P_j(\rho)/(s-\rho)+(P_j(s)-P_j(\rho))/(s-\rho)$ splits a simple pole and a
polynomial of degree at most $d-1$; the terms $P_j(s)/\rho$ and $B_jP_j(s)$ have degree at most
$d$. By hypothesis the pole part of the truncated $F$ does not change once $R$ is large, and the
remaining polynomial has degree at most $d$. Convergence at $d+1$ distinct regular points forces
convergence of its coefficients, by invertibility of the Vandermonde matrix. So $F$ is a fixed
finite sum of simple fractions plus a polynomial.

Second, write $g_a(s)=\tfrac12\psi((s+a)/2)$ with $\psi=\Gamma'/\Gamma$. Logarithmic
differentiation gives
\[
 F(s)=A(s)+\sum_jP_j(s)\Bigl[\tfrac12\log(q_j/\pi)+g_{a_j}(s)
 +\delta_j\Bigl(\frac1s+\frac1{s-1}\Bigr)\Bigr].
\]
Let $\Delta f(s)=f(s+2)-f(s)$. The recurrence of the Gamma function gives
$g_a(s+2k)=g_a(s)+\sum_{h=0}^{k-1}(s+a+2h)^{-1}$. In $\Delta^{d+1}(P_jg_{a_j})$ the
coefficient of $g_{a_j}(s)$ is $\Delta^{d+1}P_j=0$, and every other term is rational. Hence
$\Delta^{d+1}A$ is rational. For real $\sigma\to+\infty$ the Euler product gives
\begin{equation}\label{zer:eq:polydir}
 A(\sigma)=-\sum_{n\ge2}\Lambda(n)n^{-\sigma}\sum_jP_j(\sigma)\chi_j(n)
 =O\bigl((1+\sigma)^d2^{-\sigma}\bigr),
\end{equation}
for instance because $\sum_{n\ge2}(\log n)n^{-\sigma}\le2^{2-\sigma}\sum_{n\ge2}(\log n)n^{-2}$
for $\sigma\ge2$, and the same bound holds for the fixed difference $\Delta^{d+1}A$. A nonzero
rational function does not decay exponentially along the positive axis, so $\Delta^{d+1}A=0$ on
$\Re s>1$. For each real $\sigma>1$ the sequence $u_n=A(\sigma+2n)$ then has vanishing
difference of order $d+1$, so it is a polynomial in $n$; it tends to zero by
\eqref{zer:eq:polydir}, so it vanishes, and $A(\sigma)=0$.

Third, all coefficients $C_n(s)=\Lambda(n)\sum_jP_j(s)\chi_j(n)$ vanish identically: if $n_0$
were the least index with $C_{n_0}\ne0$, multiplying \eqref{zer:eq:polydir} by $n_0^\sigma$
would give $\abs{C_{n_0}(\sigma)}\ll(1+\sigma)^d(n_0/(n_0+1))^{\sigma-2}$, which is impossible
for a nonzero polynomial. Put $Q=\lcm(q_1,\dots,q_r)$. Dirichlet's theorem gives a prime
$p\equiv b\pmod Q$ for every unit $b$, and $\Lambda(p)>0$, so $\sum_jP_j(s)\chi_j(b)=0$ for
every $b\in(\Z/Q\Z)^\times$. The characters induced by the $\chi_j$ modulo $Q$ are distinct, by
uniqueness of the primitive inducing character, and orthogonality gives $P_j=0$ for all $j$.
\end{proof}

\begin{corollary}[Tail spanning; \stG]\label{zer:cor:tail}
Write $m_j(\gamma)=m_j(\tfrac12+i\gamma)$. For every $H\ge0$ and every integer $d\ge0$ the
vectors
\[
 v(\gamma)=\bigl(m_j(\gamma)(i\gamma)^\ell\bigr)_{1\le j\le r,\ 0\le\ell\le d},\qquad
 \abs\gamma>H,
\]
span $\C^{r(d+1)}$.
\end{corollary}
\begin{proof}
A linear functional with coefficients $c_{j\ell}$ annihilating these vectors gives polynomials
$P_j(s)=\sum_\ell c_{j\ell}(s-\tfrac12)^\ell$ that satisfy the hypothesis of
\cref{zer:thm:polyrigid}: by \cref{src:thm:dgrh} every zero is $\tfrac12+i\gamma$, and only
finitely many distinct ordinates have $\abs\gamma\le H$. Hence all $c_{j\ell}=0$.
\end{proof}

Common zeros, arbitrary multiplicities and both signs of the ordinates are allowed. For $d=0$
and characters to one modulus the statement also follows from \cref{zer:thm:rank} (applied
modulo the least common multiple of the conductors), which in addition localizes the spanning
zeros in every window of logarithmic length; the polynomial weights $(i\gamma)^\ell$ are what
the derivatives in \cref{zer:thm:jets} require.

\subsection{Suzuki's prime-power questions}\label{zer:sec:suzuki}
\sources{grh_cm_four_terms_and_prime_bias; grh_method_transfers_2026-11-26}

Suzuki's study of variants of Chebyshev's conjecture \cite{Suzuki2025Variants} leaves two
questions about weighted sums over prime powers: how often an unsmoothed error is negative,
and what happens to an average over moduli when the range of moduli is fixed. Both are
governed by Mellin residues at zeros, as in the source proof; here the residue at the
center computes a sign.

\subsubsection{The unsmoothed prime-power race}
For $x>1$ let
\[
 E(x)=\sideset{}{'}\sum_{n\le x}\frac{\Lambda(n)}{\sqrt n}-2\sqrt x ,
\]
where the prime on the sum gives half weight to its last term when $x$ is a prime power.
Suzuki asks for the logarithmic density of $E(x)\le0$
\cite[\S1.1, after (6)]{Suzuki2025Variants}.

\begin{theorem}[\stG]\label{zer:thm:suzukidensity}
Let $p(a,L)=L^{-1}\abs{\set{t\in[a,a+L]:E(e^t)>0}}$ for $a\ge1$, where $\abs{\cdot}$ is
Lebesgue measure. Then
\[
 10^{-937}<\liminf_{L\to\infty}\ \inf_{a\ge1}p(a,L)\le
 \limsup_{L\to\infty}\ \sup_{a\ge1}p(a,L)<0.00007943 .
\]
In particular, the lower logarithmic density of the set $E(x)<0$ exceeds $99.992057\%$, and
its upper logarithmic density is at most $1-10^{-937}$. The placement of the logarithmic
interval is arbitrary.
\end{theorem}

The only global zero-location input is the zeta case of \cref{src:thm:dgrh}. All the
independence information used in the proof is finite and verified; global simplicity and
linear independence of the zero ordinates are not assumed. The certificates are in the
following files.
\begin{center}
\artifact{grh_method_transfers_2026-11-26/checks/suzuki/verified.txt}\\
\artifact{grh_method_transfers_2026-11-26/checks/suzuki/zero_cells.txt}
\end{center}

\begin{proof}
\emph{The mean and its Fourier series.} Write $\gamma$ for a distinct positive ordinate of
a zeta zero and $m_\gamma$ for its multiplicity. Suzuki's explicit formula (36) and central
value (47) give, under RH,
\[
 E(e^t)=-c+F(t)+r(t),\qquad
 c=\tfrac12\bigl(\pi/2+\log(8\pi)+\zergE\bigr)>2.686,\qquad
 F(t)=-2\sum_{\gamma>0}\frac{m_\gamma}\gamma\sin(\gamma t),
\]
with $r(t)=\sum_{j\ge1}e^{-(2j+1/2)t}/(2j+1/2)>0$, which tends exponentially to zero. The
zero count in a unit block at height $j$ is $O(\log(2+\abs j))$ with multiplicities, so by
\cref{zer:lem:bins} the symmetric partial sums of $F$ converge in the norm
$\norm\cdot_{S^2}$, with their limit agreeing almost everywhere with the pointwise explicit
formula; this is an instance of the Stepanov approach of
\cite{KaczorowskiRamare2003Almost}. Finite-polynomial means on long intervals converge
uniformly in the starting point, by integrating each exponential. Equivalently, these means
are Haar integrals on the closure $G$ of the phase flow
$t\mapsto(e^{i\gamma t})_{\gamma>0}$ in the product of circles. Uniform mean-square
approximation extends this interpretation to the positive part of any sine-series tail,
since $y\mapsto y_+$ is $1$-Lipschitz. Write $\norm\cdot_p$ for the Haar norm. Inversion in
the phase group sends every sine sum $Y$ to $-Y$, giving
$\norm{(Y)_+}_p=2^{-1/p}\norm Y_p$ whenever the moment exists.

\emph{Moment estimates from finite certificates.} Number the zeros by increasing positive
ordinate. They are simple over the entire finite range used below. Decompose
$F=P+M_1+M_2+M_3+R$: the indices for $P$ are $1$ through $500$, each $M_j$ uses the range
in the table, and $R$ is the remaining tail. For a finite block $I$ put
$v_I=2\sum_{i\in I}\gamma_i^{-2}$.
\begin{center}\small
\begin{tabular}{rrrr}\toprule
Indices & $v_I$ (rounded) & Integers sorted & Minimum gap\\\midrule
1--500 & 0.04391240805 & \multicolumn{2}{l}{\cite[Table~2]{BestTrudgian2015Linear}}\\
501--6,000 & 0.00190153507 & 30,250,000 & 2,892\\
6,001--12,000 & 0.00016163380 & 17,997,000 & 236\\
12,001--18,000 & 0.00006286009 & 17,997,000 & 2,478\\\bottomrule
\end{tabular}

\smallskip
Gaps are in units of $10^{-14}$. The arithmetic uses unrounded enclosures.
\end{center}
For the first $500$ ordinates, Table~2 of \cite{BestTrudgian2015Linear} excludes every
nontrivial integer relation with coefficients of absolute value at most $554{,}000$. In
particular the $18$th moment of $P$ equals the moment for independent uniform phases:
a monomial in the expanded sine polynomial has coefficient vector with $\ell^1$ norm at
most $18$, and only the constant character integrates to one. For a single independent
phase with amplitude $2/\gamma$, the exponential moment series has coefficients
$1/(\gamma^{2k}(k!)^2)$ in degree $2k$ and zero in odd degree; it is coefficientwise
bounded by $\exp(z^2/\gamma^2)$ as a series in $z$. Multiplication and degree extraction
therefore give
\[
 \mathbb E[P^{2r}]\le(2r-1)!!\,v_P^r,\qquad r=9 .
\]

For each middle block only pair cancellation contributes to the fourth moment. All
$18{,}000$ ordinates were computed as certified balls with FLINT, retaining the enclosing decimal cell
of width $10^{-22}$. Put $u_i=\floor{10^{14}\gamma_i}$, computed from that cell. The table
records the minimum separation among positive differences $u_j-u_i$ with $j>i$; in the
first middle block all sums $u_j+u_i$ with $j\ge i$ are included in the same sort.
Different positive differences separated by more than three units exclude a nontrivial
two-against-two equality between ordinates, since such an equality would give two equal
positive differences after canceling common terms. A three-against-one equality would give
a sum equal to a positive difference. Their scaled rounding errors are respectively in
$[0,2]$ and $[-1,1]$, so the same gap threshold suffices. In each of the last two blocks
$3\min\gamma_i>\max\gamma_i$ already excludes three-against-one. Four terms with a common
sign have positive absolute sum. Thus the fourth moment of each $M_j$ equals its
independent-phase value and is at most $3v_{M_j}^2$. This certificate imposes no condition
on relations between the blocks.

\emph{The part beyond the certificate.} Equation (20) of
\cite{BrentPlattTrudgian2022Mean} gives $\sum_{\gamma>0}m_\gamma/\gamma^2\le0.023105$.
Simplicity is verified to $H=3\cdot10^{12}$ \cite{PlattTrudgian2021Riemann}. Above $H$ a
jump of the counting function is bounded by $m_\gamma\le0.56\log\gamma$, by Corollary~1 of
\cite{BrentPlattTrudgian2022Mean}, whose Lemmas 5 and 6 bound the sums of $1/\gamma^2$ and
$\log(\gamma/2\pi)/\gamma^2$ beyond $H$, counting multiplicity. With $D=\log H$ they give
\[
 \sum_{\gamma>H}\frac{m_\gamma(m_\gamma-1)}{\gamma^2}\le
 \frac{0.56}{2\pi H}\bigl\{D^2+(\log(2\pi)-1)D\bigr\}<2.524\cdot10^{-11}.
\]
The full variance is $2\sum m_\gamma^2/\gamma^2$ by orthogonality. Subtracting the four
finite block variances, with their certified endpoints, yields
$v_R=\norm R_2^2<0.000171563042$.

\emph{The upper bound.} Put $a_0=0.79$, $b_0=0.37$ and $M=\sum_jM_j$. Pointwise,
\[
 (F-a_0-b_0)_+\le(P-a_0)_++(M-b_0)_++R_+ .
\]
For $\tau>0$ and $r\ge2$, maximizing $(y-\tau)/y^r$ for $y\ge\tau$ gives
$(y-\tau)_+\le(r-1)^{r-1}(y_+)^r/(r^r\tau^{r-1})$. With $r=9$ for $P$ and $r=2$ for $M$,
symmetry and Minkowski's inequality give
\begin{align*}
 \norm{(P-a_0)_+}_2&\le\frac{(r-1)^{r-1}}{r^ra_0^{r-1}}\sqrt{(2r-1)!!\,v_P^r/2}
   <0.000923214,\\
 \norm{(M-b_0)_+}_2&\le\frac{\sqrt{3/2}}{4b_0}\Bigl(\sum_{j=1}^3\sqrt{v_{M_j}}\Bigr)^2
   <0.003415944,\\
 \norm{R_+}_2&=\sqrt{v_R/2}<0.009261832 .
\end{align*}
Chebyshev's inequality bounds the upper frequency of $F>c$ by the square of the sum of
these norms divided by $(c-a_0-b_0)^2$. To absorb $r(t)$ use the threshold $c-\eta$,
discard a fixed initial interval at cost $O_\eta(1/L)$, then let $\eta\downarrow0$. The
unrounded computation gives an upper bound below $0.00007942911$. Equality $E(x)=0$ occurs
on a countable set, since $E$ decreases strictly between consecutive prime powers.

\emph{The lower bound.} Let $K_T$ be convolution along the one-parameter flow on $G$ with
the nonnegative real-line Fej\'er kernel of integral one and Fourier multiplier
$(1-\abs\lambda/T)_+$. Then $K_T$ preserves nonnegativity and constants, is self-adjoint,
and contracts the supremum norm. For $T=1420$ the first $1{,}000$ ordinates lie below the
cutoff, with the next above. Put $d_i=2(1/\gamma_i-1/T)$ for $1\le i\le1{,}000$. The
smoothed sum $K_TF$ is the finite polynomial $-\sum_id_i\sin\theta_i$, where
$e^{i\theta_i}$ is coordinate $i$ on $G$.

Set $k=23$, and normalize $w(\theta)=(1-\sin\theta)^k$ to have mean one on the circle.
Direct integration gives
\[
 \max w=4^k\Big/\binom{2k}k=:B,\qquad
 \int w(\theta)(-\sin\theta)\,\frac{\dd\theta}{2\pi}=\frac k{k+1}:
\]
the circle mean before normalization is $2^{-k}\binom{2k}k$, and the second identity
follows by taking the ratio of the successive means and subtracting one. Set
$W=\prod_{i=1}^{500}w(\theta_i)$. The finite certificate of \cite{BestTrudgian2015Linear}
applies coordinate by coordinate to every monomial of $W$ and of $W\sin\theta_i$ for
$i\le500$, since no exponent exceeds $24$. Thus $\int_GW=1$, while $0\le W\le B^{500}$. For
the remaining $500$ modes use only $\abs{\sin\theta_i}\le1$. Therefore
\[
 \int_GW\,(K_TF-c)\ge\Delta:=\frac k{k+1}\sum_{i=1}^{500}d_i-\sum_{i=501}^{1000}d_i-c
 >0.00091075 .
\]
Moving the convolution onto $W$ and using $0\le K_TW\le B^{500}$ gives
\[
 \Delta\le B^{500}\int_{F>c}(F-c)\le B^{500}\,\norm F_2\,\mathbb P(F>c)^{1/2}.
\]
The arithmetic verifies $\Delta^2/(B^{1000}\norm F_2^2)>10^{-937}$, using the preceding
upper bound for $\norm F_2^2$. To pass to a uniform lower frequency, bounded Lipschitz test
functions of $F$ have Haar means uniformly on long intervals, by the local mean-square
approximation above. Approximate the indicator of the open half-line $(c,\infty)$ from
below by such functions. The decaying remainder $r(t)$ also tends to zero in mean,
uniformly on long intervals starting at $a\ge1$.
\end{proof}

The existence and the value of the ordinary logarithmic density at this threshold are not
determined by the argument: a finite list of ruled-out resonances controls specified
moments and leaves possible relations with arbitrarily high ordinates.

\subsubsection{The average over moduli at fixed range}
Suzuki's equation (29) \cite{Suzuki2025Variants} concerns
\[
 A_Q(x)=\sum_{q=3}^Q\Bigl(\varphi(q)\sum_{\substack{n\le x\\n\equiv1\bmod q}}
   \frac{\Lambda(n)}{\sqrt n}\log\frac xn-4\sqrt x\Bigr).
\]
His Theorem~7 gives negativity when $Q\ge x$, and he asks about a large fixed $Q$. The
order of limits changes the answer.

\begin{theorem}[\stG]\label{zer:thm:suzukiQ}
For each fixed integer $3\le Q\le138$, $A_Q(x)$ is eventually negative. For each fixed
$139\le Q\le2{,}000{,}000$, it is eventually positive. More generally, for any fixed
$Q\ge139$ it is eventually negative if and only if some primitive Dirichlet $L$-function of
conductor at most $Q$ vanishes at $1/2$; in the other case it is eventually positive. The
starting value of $x$ may depend on $Q$.
\end{theorem}

Thus eventual negativity for arbitrarily large fixed $Q$ in (29) is equivalent to the
existence of a single central Dirichlet zero.

\begin{proof}
\emph{The central residue.} Set $Z_q(s)=\prod_{\chi\bmod q}L(s,\chi)$,
$m_q=\ord_{1/2}Z_q$, and define $C(q)$ by
$Z_q'/Z_q(1/2+z)=m_q/z+C(q)+O_q(z)$. Put $M(Q)=\sum_{q=3}^Qm_q$ and
$B(Q)=\sum_{q=3}^QC(q)$. Character orthogonality and the twice-integrated explicit formula
give
\begin{equation}\label{zer:eq:AQ}
 A_Q(x)=-\tfrac12M(Q)(\log x)^2-B(Q)\log x+O_Q(1).
\end{equation}
In detail, the Mellin integrand for the $q$th inner weighted sum, including its factor
$\varphi(q)$, is $-(Z_q'/Z_q)(1/2+z)x^z/z^2$. The pole at $z=1/2$ contributes $4\sqrt x$.
The residue at zero yields the two displayed powers and a constant. The noncentral critical
zeros yield a bounded oscillating sum because their denominators are $\abs{\Im\rho}^2$ and
the modulus is fixed. For primitive factors this is the integrated explicit formula of
\cite[\S\S3.1, 4.1, 5.1]{Suzuki2025Variants}. For each Euler factor deleted on induction,
expand its logarithmic derivative in powers of $p^{-1/2-z}$; the coefficients at $z=0$ are
geometrically decreasing, so termwise summation with $\log(x/p^j)$ adds its linear term and
a bounded remainder. This also justifies the formula for the principal character.

\emph{An elementary coefficient valid at a central zero.} Let $h=\varphi(q)$,
$K=\Q(\zeta_q)$ and $D=\abs{\disc K}$. For $p^a\Vert q$ put $q'=q/p^a$, $e=\varphi(p^a)$,
$f=\ord_{q'}(p)$ with $f=1$ when $q'=1$, and $g=h/(ef)$. Euler products give
$Z_q=\zeta_K\prod_{p\mid q}(1-p^{-fs})^g$. The completed Dedekind zeta function, up to a
fixed positive multiple, is $D^{s/2}(2\pi)^{-hs/2}\Gamma(s)^{h/2}\zeta_K(s)$ and is
symmetric about $1/2$. Its logarithmic derivative therefore has zero constant Laurent
coefficient at $1/2$, including at a zero. The cyclotomic discriminant formula
$\log D=h(\log q-\sum_{p\mid q}\log p/(p-1))$ and
$\Gamma'(1/2)/\Gamma(1/2)=-\zergE-2\log2$ now give
\begin{equation}\label{zer:eq:Cq}
 C(q)=\frac h2\Bigl(\log(8\pi)+\zergE-\log q+\sum_{p\mid q}\frac{\log p}{p-1}\Bigr)
 +\sum_{p^a\Vert q}\frac{\varphi(q/p^a)\log p}{p^{f/2}-1}.
\end{equation}
Thus both the cutoff computation and the large-$q$ sign estimate apply irrespective of
uncomputed central values.

\emph{Certification of the crossover.} A verifier encloses $B(Q)$ for every
$3\le Q<10^4$; it is the file
\begin{center}
\artifact{grh_cm_four_terms_and_prime_bias/checks/suzuki_fixed_q.py}
\end{center}
It finds
\[
 45.431<B(138)<45.490,\qquad-29.828<B(139)<-29.767,
\]
with $B(Q)>0$ throughout $3\le Q\le138$ and $B(Q)<0$ throughout $139\le Q<10^4$. It uses
rational remainder bounds for the logarithm series after binary range reduction and integer
square roots for the ramified Euler factors. Machin's identity, the harmonic sum through
$1024$ with a trapezoid-rule remainder, and the same logarithm series certify
$3.80138<\log(8\pi)+\zergE<3.80140$.

For $q\ge10^4$, each $C(q)$ is negative. Indeed
\[
 \sum_{p\mid q}\frac{\log p}{p-1}\le\tfrac14\log q+\tfrac34\log2+\tfrac14\log3
 <\tfrac14\log q+0.8,
\]
where only the primes $2,3$ require the two corrections. As $p^f\ge q'+1$, we have
$p^{f/2}-1\ge(\sqrt2-1)\sqrt{q'}$, also for $q'=1$. Since $e\ge p^a/2$ and $p^a\ge2$, the
last sum in \eqref{zer:eq:Cq} divided by $h$ is at most $(2+\sqrt2)\log q/\sqrt q$.
Consequently $C(q)/h<2.31-(3/8-0.035)\log q<0$ in this range. It follows that $B(Q)<0$
for every $Q\ge139$.

The central orders are nonnegative. Imprimitive Euler factors do not vanish at $1/2$, so
$M(Q)>0$ exactly when a primitive character of conductor at most $Q$ gives a central zero
(the zeta value at $1/2$ is nonzero). Platt's computation verifies central nonvanishing
through modulus $2{,}000{,}000$ \cite[Theorem~9.2 of the author manuscript]{Platt2016Numerical};
this central-value bound is distinct from the finite-height GRH bound in the same paper. Substitution in
\eqref{zer:eq:AQ} proves the theorem.
\end{proof}

\begin{proposition}[The second cutoff; \stG]\label{zer:prop:UQ}
Suzuki's equation (30) uses the sum
\[
 U_Q(x)=\sum_{q=3}^Q\varphi(q)\sum_{\substack{n\le xe^2\\n\equiv1\bmod q}}
   \frac{\Lambda(n)}{\sqrt n}\log\frac xn .
\]
For every fixed $Q\ge3$ its sign agrees with the eventual sign of $A_Q$ outside a set
$\mathcal E_Q\subset[2,\infty)$ with $\int_{\mathcal E_Q}\dd x/x<\infty$.
\end{proposition}
\begin{proof}
Put $y=xe^2$ and
$D_Q(y)=\sum_{q=3}^Q\varphi(q)\sum_{n\le y,\,n\equiv1\bmod q}\Lambda(n)/\sqrt n$. The
unsmoothed explicit formulas (36), (43), (44) of \cite{Suzuki2025Variants} give
\[
 D_Q(e^t)=2(Q-2)e^{t/2}-M(Q)\,t+R_Q(t),\qquad
 \sup_{b\ge1}\int_b^{b+1}\abs{R_Q(t)}^2\,\dd t<\infty .
\]
Indeed each noncentral zero contributes a frequency with coefficient of absolute value
$1/\abs{\Im\rho}$ per unit of multiplicity, the fixed-conductor zero count in a unit block
is $O_Q(\log(2+\abs k))$, and \cref{zer:lem:bins} applies to the finite family. The
finitely many deleted Euler factors contribute bounded terms through geometrically
decreasing prime-power sums. Endpoints at prime powers form a countable exceptional set.
From $U_Q(x)=A_Q(y)+4(Q-2)\sqrt y-2D_Q(y)$ and $\log y=\log x+2$ we obtain, writing
$u=\log x$,
\[
 U_Q(e^u)=-\tfrac12M(Q)u^2-B(Q)u+H_Q(u),\qquad
 \sup_{b\ge1}\int_b^{b+1}\abs{H_Q(u)}^2\,\dd u<\infty .
\]
Since $B(Q)\ne0$ by the proof of \cref{zer:thm:suzukiQ}, the absolute value of the
polynomial part exceeds $c_Qk$ on $[k,k+1]$ for all sufficiently large integers $k$, for
some $c_Q>0$. The set of sign failures in this interval has measure $O_Q(k^{-2})$. Summing
proves the assertion.
\end{proof}

\begin{proposition}[A growing range; \stU]\label{zer:prop:growing}
For every fixed $\eps>0$, $A_Q(x)<0$ for all sufficiently large $x$ and all integers
$Q\ge(1/9+\eps)x$.
\end{proposition}
\begin{proof}
The prime weights in $A_Q$ are nonnegative and $\sum_{q\mid n-1}\varphi(q)=n-1$ for
$n\ge2$. Thus, uniformly in integers $Q\ge3$,
\[
 A_Q(x)\le\sum_{n\le x}\Lambda(n)\sqrt n\,\log(x/n)-4(Q-2)\sqrt x
 =(4/9+o(1))x^{3/2}-4(Q-2)\sqrt x .
\]
Partial summation with the prime number theorem gives the last equality, since
\[
 \int_0^1y^{1/2}\log(1/y)\,\dd y=\frac49 .
\]
\end{proof}

\subsection{Covariance and oscillation of smoothed prime errors}\label{zer:sec:cov}
{\emergencystretch=6em\sources{multiscale_prime_fluctuations; grh_multiscale_covariance; grh_spectral_rigidity; grh_filtered_prime_races; grh_prime_error_balancing; prime_peak_rigidity}}

The kernel $(1/n-1/x)_+$ of the prime remainder in the source proof has Mellin denominator
$s(s-1)$, and $\rho(1-\rho)=\tfrac14+\gamma^2$ at $\rho=\tfrac12+i\gamma$. Under
\cref{src:thm:dgrh} the normalized smoothed errors are therefore absolutely convergent
almost periodic functions of $\log x$. This subsection proves that their joint covariances,
at several multiplicative scales, in all residue classes, with derivatives, or after a
compact filter, are positive definite even after any finite set of zero frequencies is
deleted, and derives uniform two-sided oscillation in every direction. Simplicity of zeros,
disjointness of zero sets and rational independence of ordinates are never assumed:
coincident ordinates are grouped before squares are formed, and positivity comes from the
arithmetic of prime powers.

\paragraph{Setting.}
Fix $q\ge1$, put $G=(\Z/q\Z)^\times$ (the trivial group for $q=1$) and $n_q=\varphi(q)$, and let
$\chi$ run over the characters modulo $q$, extended by zero to nonunits, with primitive
inducing character $\chi^*$ and $\delta_\chi=1$ for the principal character, $0$ otherwise. For
$x\ge1$ and $a\in G$ put
\[
 S_\chi(x)=\sum_{n\le x}\Lambda(n)\chi(n)\Bigl(\frac1n-\frac1x\Bigr),\qquad
 S_a(x)=\sum_{\substack{n\le x\\n\equiv a\bmod q}}\Lambda(n)\Bigl(\frac1n-\frac1x\Bigr)
 =\frac1{n_q}\sum_\chi\overline{\chi(a)}S_\chi(x),
\]
and
\[
 c_\chi=\lim_{s\to1}\Bigl(-\frac{L'}{L}(s,\chi)-\frac{\delta_\chi}{s-1}\Bigr)-\delta_\chi,\qquad
 c_a=\frac1{n_q}\sum_\chi\overline{\chi(a)}c_\chi .
\]
For the principal character $c_\chi=-1-\zergE-\sum_{p\mid q}\log p/(p-1)$. The weight vanishes at
$n=x$, so $S_\chi$ and $S_a$ are continuous. Under \cref{src:thm:dgrh} let $m_\chi(\gamma)$ be the
order of $L(s,\chi^*)$ at $\tfrac12+i\gamma$, let $\Gamma_q$ be the set of distinct ordinates of
all the $L(s,\chi^*)$, and put
\begin{equation}\label{zer:eq:va}
 v_a(\gamma)=\frac1{n_q}\sum_\chi\overline{\chi(a)}\,m_\chi(\gamma),\qquad
 \beta_a=4v_a(0),\qquad
 Z_a(t)=e^{t/2}\Bigl(S_a(e^t)-\frac t{n_q}-c_a\Bigr)-\beta_a\quad(t\ge0).
\end{equation}
Since $m_{\overline\chi}(-\gamma)=m_\chi(\gamma)$, one has $v_a(-\gamma)=\overline{v_a(\gamma)}$, and
$c_a,\beta_a$ are real. The constant $\beta_a$ records possible central zeros; their absence is
not assumed. For a character put
\[
 K_\chi(t)=\sum_{\gamma}\frac{m_\chi(\gamma)e^{i\gamma t}}{1/4+\gamma^2},
\]
the sum over all ordinates of $L(s,\chi^*)$, the central one included. The local zero count
$O_q(\log(2+\abs y))$ in $[y-1,y+1]$ makes this series absolutely and uniformly convergent.

\begin{lemma}[Ramp explicit formula; \stG]\label{zer:lem:ramp}
For every $\chi$ modulo $q$ and $x\ge2$,
\[
 S_\chi(x)=\delta_\chi\log x+c_\chi+x^{-1/2}K_\chi(\log x)+O_q\bigl((1+\log x)/x\bigr).
\]
Consequently, uniformly in $a\in G$, for $t\ge\log2$,
\begin{equation}\label{zer:eq:Zexp}
 Z_a(t)=\sum_{\substack{\gamma\in\Gamma_q\\\gamma\ne0}}\frac{v_a(\gamma)}{1/4+\gamma^2}\,e^{i\gamma t}
 +O_q\bigl((1+t)e^{-t/2}\bigr),
\end{equation}
and the series converges absolutely and uniformly on $\R$.
\end{lemma}
\begin{proof}
First let $\chi$ be primitive and $A(s)=-L'(s,\chi)/L(s,\chi)$. For $c>1$, termwise Mellin
inversion of $y^s/(s(s-1))$, whose inverse transform is $(y-1)_+$, gives
\[
 S_\chi(x)=\frac1{2\pi i}\int_{c-i\infty}^{c+i\infty}A(s)\frac{x^{s-1}}{s(s-1)}\,\dd s .
\]
Move the line to $\Re s=-\tfrac12$. The functional equation and Stirling's formula give
$A(-\tfrac12+iy)=O_q(\log(2+\abs y))$. The local zero count lets the horizontal edges pass at
heights tending to $\pm\infty$ at distance $\gg_q1/\log(2+\abs y)$ from every ordinate; there
the local partial-fraction expansion bounds $A$ by $O_q(\log^2(2+\abs y))$, so these edges
vanish for fixed $x$, and the new vertical integral is $O_q(x^{-3/2})$. The residue at $s=1$
is $\delta_\chi\log x+c_\chi$. At a zero $\rho=\tfrac12+i\gamma$ of multiplicity $m$ it is
$mx^{\rho-1}/(\rho(1-\rho))=mx^{-1/2+i\gamma}/(\tfrac14+\gamma^2)$. At $s=0$ the function $A$
has at most a simple pole, so the residue is $x^{-1}$ times a polynomial of degree at most one
in $\log x$; no other pole is crossed. For imprimitive $\chi$, each prime $p\mid q$ with
$p\nmid f_\chi$ removes
\[
 \sum_{p^k\le x}\chi^*(p)^k\log p\Bigl(\frac1{p^k}-\frac1x\Bigr)
 =\frac{\chi^*(p)\log p}{p-\chi^*(p)}+O_p\Bigl(\frac{1+\log x}x\Bigr),
\]
by the geometric tail and the at most $\log x/\log p$ terms of the $1/x$ sum. The constant is
the change of $c_\chi$ caused by the removed Euler factor, and the nontrivial zeros are those
of $\chi^*$. Character orthogonality gives \eqref{zer:eq:Zexp}; the frequency $\gamma=0$
contributes $4v_a(0)=\beta_a$.
\end{proof}

All the statements below rest on one elementary principle.

\begin{lemma}[Almost periodic vectors; \stU]\label{zer:lem:ap}
Let $F\subset\R\setminus\set0$ be countable with $F=-F$, let $z_\gamma\in\C^N$ satisfy
$z_{-\gamma}=\overline{z_\gamma}$ and $B=\sum_{\gamma\in F}\norm{z_\gamma}<\infty$, and put
$V(t)=\sum_\gamma z_\gamma e^{i\gamma t}$. Let $X:[t_0,\infty)\to\R^N$ be locally bounded and
measurable.
\begin{enumerate}[label=\textup{(\roman*)}]
\item If $\norm{X(t)-V(t)}\to0$, or if $X-V$ is bounded and $\int_{t_0}^\infty\norm{X-V}\,\dd t<\infty$,
then, uniformly in $u\ge t_0$ as $L\to\infty$,
\[
 \frac1L\int_u^{u+L}X\,\dd t\to0,\qquad
 \frac1L\int_u^{u+L}XX^{\mathsf T}\dd t\to\mathcal G:=\sum_\gamma z_\gamma z_\gamma^*,
\]
and $\mathcal G$ is real symmetric positive semidefinite. For every $w\in\R^N$ and $c\in\R$ the
mean of $(w\cdot X-c)^2$ over $[u,u+L]$ tends to $w^{\mathsf T}\mathcal Gw+c^2$.
\item Suppose $\lambda=\lambda_{\min}(\mathcal G)>0$ and $\norm{X-V}\to0$. Put
$\eps=\lambda/(32B)$ and $\eta=\lambda/(16B^2)$. There are $L_0,T_0$ such that for every
unit $w\in\R^N$, every $T\ge T_0$ and every $L\ge L_0$,
\[
 \abs{\set{t\in[T,T+L]:w\cdot X(t)\ge\eps}}\ge\eta L,\qquad
 \abs{\set{t\in[T,T+L]:w\cdot X(t)\le-\eps}}\ge\eta L .
\]
If $X$ is continuous, the convex hull of $\set{X(t):T\le t\le T+L}$ contains the closed ball
of radius $\eps$ about $0$. In any case the closed convex hull of the cluster set of $X(t)$ as
$t\to\infty$ contains that ball.
\item Suppose $\lambda>0$, $X$ is continuous, $X-V$ is bounded and integrable on
$[t_0,\infty)$, and put $B'=\max(1,\sup\norm X)$. There is $L_1$ such that every interval
$[u,u+L]$ with $u\ge t_0$ and $L\ge L_1$ contains points $t_\pm$ with
$\pm w\cdot X(t_\pm)\ge\lambda/(8B')$, for every unit $w$.
\end{enumerate}
\end{lemma}
\begin{proof}
(i) The mean of $e^{i\omega t}$ over $[u,u+L]$ has modulus at most $2/(L\abs\omega)$ for
$\omega\ne0$. This proves the limits for finite truncations of $V$, the second by retaining
only the pairs $(\gamma,-\gamma)$. Absolute summability bounds the truncation error uniformly,
and the hypothesis on $X-V$ makes its contribution to both means $o(1)$ uniformly in $u$. The
matrix is real because $z_{-\gamma}=\overline{z_\gamma}$. The last assertion follows by
expanding the square.

(ii) Let $g=w\cdot V$, so $\abs g\le B$. By (i), uniformly in $u$ and in unit $w$ (the
convergence is that of finitely many vector and matrix entries), there is $L_0$ with
$\abs{\mathbb E g}\le\lambda/(4B)$ and $\mathbb Eg^2\ge\lambda/2$ for windows of length
$L\ge L_0$, where $\mathbb E$ is the window average. Since $g^2\le B\abs g$, we get
$\mathbb E\abs g\ge\lambda/(2B)$ and $\mathbb Eg_\pm=(\mathbb E\abs g\pm\mathbb Eg)/2\ge\lambda/(8B)$.
With $\theta=2\eps=\lambda/(16B)$, the inequality $g_+\le\theta+B\mathbf 1_{g\ge\theta}$ gives
$\mathbb P(g\ge\theta)\ge\eta$, and likewise for $-g$. Choose $T_0$ with $\norm{X-V}\le\eps$
beyond $T_0$. For continuous $X$ the hull of a compact image is compact; if a point of the
ball were outside it, strict separation would give a unit $w$ whose maximum on the hull is
less than $\eps$, contradicting what was just proved. For the cluster set, each unit $w$ has
arbitrarily late $t$ with $w\cdot X(t)\ge\eps$; a convergent subsequence gives a cluster point
with the same inequality, and separation again gives the inclusion.

(iii) Let $f=w\cdot X$, so $\abs f\le B'$. By (i) there is $L_1$ with $\mathbb Ef^2\ge\lambda/2$
and $\abs{\mathbb Ef}\le\lambda/(4B')$ on every window $[u,u+L]$, $u\ge t_0$, $L\ge L_1$. Then
$\mathbb E\abs f\ge\lambda/(2B')$ and $\mathbb Ef_\pm\ge\lambda/(8B')$, and a continuous function
attains a maximum at least its mean.
\end{proof}

\paragraph{Several scales.}
Fix $r\ge1$ distinct real numbers $h_1,\dots,h_r$; the observed multiplicative scales are
$e^{h_j}$. Put $D=rn_q$. For $w=(w_{j,a})\in\C^D$, real $\gamma$, and $1<T<T+H$, define
\[
 A_w(\gamma)=\sum_{j=1}^r\sum_{a\in G}w_{j,a}e^{i\gamma h_j}v_a(\gamma),\qquad
 V_{T,H}(w)=2\sum_{\substack{\gamma\in\Gamma_q\\T<\gamma\le T+H}}
 \frac{\abs{A_w(\gamma)}^2}{(1/4+\gamma^2)^2}.
\]
For real $w$, $V_{T,H}(w)$ is the mean square of the real signal
$2\Re\sum_{T<\gamma\le T+H}A_w(\gamma)(\tfrac14+\gamma^2)^{-1}e^{i\gamma t}$ carried by the band.

\begin{theorem}[Multiscale spectral coercivity; \stG]\label{zer:thm:coercive}
For fixed $q$ and distinct $h_1,\dots,h_r$ there are positive constants $c,C,K,T_0$,
depending only on $q$ and the shifts, such that for every $T\ge T_0$, every $H$ with
$K\log(2T)\le H\le T$ and every $w\in\C^D$,
\begin{equation}\label{zer:eq:coercive}
 \frac{cH}{T^4\log(2T)}\norm w_2^2\ \le\ V_{T,H}(w)\ \le\ \frac{CH\log^2(2T)}{T^4}\norm w_2^2 .
\end{equation}
In particular each such band contains at least $D$ distinct ordinates, and the vectors
$(e^{i\gamma h_j}v_a(\gamma))_{j,a}$ for the ordinates in the band span $\C^D$.
\end{theorem}

No uniformity in $q$, in the number of scales, or as two shifts approach each other is
asserted; the dependence on merging scales is the subject of \cref{zer:sec:coalesce}. The proof
uses two ingredients. The first is Landau's formula at a fixed argument: for real $u\ge2$,
uniformly for $1<T<T+H\le2T$,
\begin{equation}\label{zer:eq:landaufixed}
 \sum_{T<\gamma\le T+H}m_\chi(\gamma)e^{i\gamma\log u}
 =-\frac{H}{2\pi\sqrt u}\Lambda(u)\chi^*(u)+O_{q,u}(\log(2T)),
\end{equation}
where the main term is zero unless $u$ is an integral prime power. This is Banks's extension of
the Landau--Gonek formula \cite[Theorem~1.2]{Banks2023Comparing}, divided by $\sqrt u$ using
\cref{src:thm:dgrh}; for fixed $u$ the distance to the nearest other prime power is positive,
and half-open instead of open endpoints change the sum by $O_{q,u}(\log(2T))$.

\begin{lemma}[Prime probes; \stU]\label{zer:lem:probes}
For a real $R>0$ with $R\ne1$, only finitely many primes $p$ make $Rp$ an integral prime power.
Consequently, for each $a\in G$ there is a prime $p_a\equiv a\pmod q$ such that
$u_{a,j,k}=p_ae^{h_k-h_j}\ge2$ for all $j,k$, and $u_{a,j,k}$ is an integral prime power only for
$j=k$.
\end{lemma}
\begin{proof}
If $Rp$ is an integer, $R=u/v$ is rational in lowest terms. If $v>1$ then $v\mid p$, which
allows one prime at most; if $v=1$ then $u\ge2$ and $up$ is a prime power only if $p\mid u$.
For the probes, exclude the finitely many primes dividing a numerator or denominator of a
rational ratio $e^{h_k-h_j}\ne1$, require $p_a\min_{j,k}e^{h_k-h_j}\ge2$, and use Dirichlet's
theorem.
\end{proof}

\begin{proof}[Proof of \cref{zer:thm:coercive}]
Let $B_{T,H}$ be the set of ordinates in $(T,T+H]$ and $L=\log(2T)$. For fixed $j,a$,
\[
 \sum_{\gamma\in B_{T,H}}A_w(\gamma)e^{i\gamma(\log p_a-h_j)}
 =\frac1{n_q}\sum_{k,b,\chi}w_{k,b}\overline{\chi(b)}\sum_{\gamma\in B_{T,H}}
 m_\chi(\gamma)e^{i\gamma\log u_{a,j,k}} .
\]
By \eqref{zer:eq:landaufixed} and \cref{zer:lem:probes} only $k=j$ gives a main term. There
$u=p_a$ is prime to $q$, so $\chi^*(p_a)=\chi(a)$, and orthogonality isolates $w_{j,a}$:
\begin{equation}\label{zer:eq:recover}
 \sum_{\gamma\in B_{T,H}}A_w(\gamma)e^{i\gamma(\log p_a-h_j)}
 =-\frac{H\log p_a}{2\pi\sqrt{p_a}}\,w_{j,a}+O_{q,h}\bigl(\norm w_2L\bigr),
\end{equation}
uniformly in $j,a,w$, since the errors total at most a constant times
$\sum_{k,b}\abs{w_{k,b}}L\le\sqrt D\norm w_2L$. Put $\alpha=\min_a\log p_a/(2\pi\sqrt{p_a})$,
choose $(j,a)$ with $\abs{w_{j,a}}\ge\norm w_2/\sqrt D$, and take $K$ so large that for
$H\ge KL$ the left side of \eqref{zer:eq:recover} has modulus at least
$H\alpha\norm w_2/(2\sqrt D)$. With $M_\gamma=\sum_\chi m_\chi(\gamma)$ the zero count gives
$\abs{B_{T,H}}\le\sum_{\gamma\in B_{T,H}}M_\gamma\le C_0HL$ and $\max M_\gamma\le C_1L$.
Cauchy--Schwarz over the distinct ordinates gives
\[
 \sum_{\gamma\in B_{T,H}}\abs{A_w(\gamma)}^2\ge\frac{\alpha^2}{4DC_0}\frac HL\norm w_2^2 ,
\]
and $\tfrac14+\gamma^2\le5T^2$ in the band gives the lower bound with $c=\alpha^2/(50DC_0)$. For
the upper bound, $\abs{v_a(\gamma)}\le M_\gamma/n_q$ gives
$\abs{A_w(\gamma)}\le(\sqrt D/n_q)M_\gamma\norm w_2$, so the sum of squares is at most
$(D/n_q^2)C_0C_1HL^2\norm w_2^2$, and $(\tfrac14+\gamma^2)^2\ge T^4$. The lower bound makes
$w\mapsto(A_w(\gamma))_{\gamma\in B_{T,H}}$ injective, which gives the last assertion.
\end{proof}

Thus any coefficient array is recovered from aggregated zero data of one band by
\eqref{zer:eq:recover}, with error $O(\norm w_2\log(2T)/H)$ per entry. For real arrays the
matrix $Q_{T,H}$ of the form $V_{T,H}$ has all its eigenvalues between the two bounds of
\eqref{zer:eq:coercive}, so $\det Q_{T,H}\ge(cH/(T^4\log(2T)))^D$ and its condition number is
at most $(C/c)\log^3(2T)$.

\begin{corollary}[All-lag strict positivity; \stG]\label{zer:cor:multiscale}
Let $J$ be a finite subset of $\Gamma_q$ with $J=-J$, and define the real vector
$U_J(t)\in\R^D$ by
\[
 U_J(t)_{j,a}=Z_a(t+h_j)-\sum_{\gamma\in J\setminus\set0}\frac{v_a(\gamma)}{1/4+\gamma^2}e^{i\gamma(t+h_j)} .
\]
Its long-time covariance in the sense of \cref{zer:lem:ap}\textup{(i)} is
\[
 \mathcal G_J=\sum_{\gamma\in\Gamma_q\setminus(J\cup\set0)}z(\gamma)z(\gamma)^*,\qquad
 z(\gamma)=\Bigl(\frac{e^{i\gamma h_j}v_a(\gamma)}{1/4+\gamma^2}\Bigr)_{j,a},
\]
and it is positive definite: $\lambda_{\min}(\mathcal G_J)\ge cH/(T^4\log(2T))$ for every band
$(T,T+H]$ of \cref{zer:thm:coercive} lying above $\max J$. Consequently
\cref{zer:lem:ap}\textup{(ii)} applies to $U_J$ with constants depending only on $q$, the
shifts and $J$: for every real unit array $w$, each of $w\cdot U_J(\log x)\ge\eps$ and
$w\cdot U_J(\log x)\le-\eps$ holds on a set of logarithmic measure at least $\eta\log K$ in
every interval $[X,KX]$ with $X$ and $K$ sufficiently large; the convex hull of the values on
each such interval contains the ball of radius $\eps$. Moreover, for one scale $j$, a nonzero
$w\in\R^{n_q}$, arbitrary real predictor coefficients $b_{k,a}$ $(k\ne j)$ and an arbitrary
intercept $c$,
\[
 \lim_{L\to\infty}\frac1L\int_u^{u+L}\Bigl(\sum_aw_aZ_a(t+h_j)-\sum_{k\ne j,\,a}b_{k,a}Z_a(t+h_k)-c\Bigr)^2\dd t
 \ \ge\ \lambda_{\min}(\mathcal G_\varnothing)\norm w_2^2 .
\]
\end{corollary}
\begin{proof}
By \eqref{zer:eq:Zexp}, $U_J(t)=\sum_{\gamma\notin J\cup\set0}z(\gamma)e^{i\gamma t}+O((1+t)e^{-t/2})$,
with $z(-\gamma)=\overline{z(\gamma)}$ and $\sum\norm{z(\gamma)}<\infty$, so \cref{zer:lem:ap}(i)
gives $\mathcal G_J$. For real $w$, $w\cdot z(\gamma)=A_w(\gamma)/(\tfrac14+\gamma^2)$; pairing $\gamma$
with $-\gamma$ gives $w^{\mathsf T}\mathcal G_Jw\ge V_{T,H}(w)$ for every band above $\max J$, and
\eqref{zer:eq:coercive} gives the bound. The ramp sums are continuous, so the hull statement
applies. For prediction, collect the coefficients in an array whose block at scale $j$ is
$w$; its squared norm is at least $\norm w_2^2$, and \cref{zer:lem:ap}(i) gives the limit
$v^{\mathsf T}\mathcal G_\varnothing v+c^2$.
\end{proof}

The theorem says more than distinctness of zero sets: common ordinates contribute to the
same vector $z(\gamma)$, and positivity survives the deletion of any bounded part of the
spectrum. Arbitrary real shifts are allowed, including logarithms of rational numbers
arbitrarily close to $1$, because the probes separate scaled prime powers rather than zeros.
The characters attached to a modulus give an invertible change of coordinates of the vector
$(Z_a)_a$, so the same conclusions hold for the errors of any family of distinct primitive
characters whose conductors divide $q$ (their imprimitive Euler factors only add terms
$O_q((1+t)e^{-t/2})$). If $\sum_jw_{j,a}=0$ for every $a$, the central constants cancel, so
with $R_a(x)=S_a(x)-\log x/n_q-c_a$ the combination
$\sum_{j,a}w_{j,a}\sqrt{e^{h_j}x}\,R_a(e^{h_j}x)$ oscillates in both directions without any
information on central zeros; for instance $\sqrt{cx}\,R_a(cx)-\sqrt x\,R_a(x)$, for fixed
$c>1$, exceeds a positive constant and falls below its negative on sets of positive lower
logarithmic density.

\begin{example}\label{zer:ex:probes}
For $q=12$ and the four shifts $0,\log2,\log(3/2),\log5$ (scales $1,2,3/2,5$), the primes
$13,17,19,11$ in the classes $1,5,7,11$ are probes in the sense of \cref{zer:lem:probes}: each
product with an off-diagonal scale ratio is either not an integer or has two distinct prime
factors. \Cref{zer:thm:coercive} therefore gives rank $16$ in every band of the stated sizes.
For $q=3$ and scales $1,2,3/2$ the primes $7$ and $5$ serve the classes $1$ and $2$. These
exact rational checks are in the files
\begin{center}
\artifact{multiscale_prime_fluctuations/verify_examples.py}\\
\artifact{grh_multiscale_covariance/numerical_check.py}
\end{center}
\end{example}

\paragraph{A certificate from one isolated prime atom.}
The probe argument also has a distributional form. For a primitive character $\eta$ of conductor
$Q$, parity $\kappa$ and $\delta=1$ for conductor one, the distribution
$U_\eta(t)=\sum_\gamma m_\eta(\gamma)e^{i\gamma t}$ (pairing with test functions converges by the
zero count) satisfies, on $(0,\infty)$,
\begin{equation}\label{zer:eq:atoms}
 U_\eta=-\sum_{n\ge2}\frac{\Lambda(n)\eta(n)}{\sqrt n}\,\delta_{\log n}+g_\eta,\qquad
 g_\eta(t)=-\frac{e^{-(\kappa+1/2)t}}{1-e^{-2t}}+\delta\bigl(e^{-t/2}+e^{t/2}\bigr).
\end{equation}
Indeed the second logarithmic derivative of the completed function $\Xi_\eta$ is
$\sum_\gamma m_\eta(\gamma)(s-\tfrac12-i\gamma)^{-2}$, and for $\Re s>1$ it equals
$-\sum_n\Lambda(n)\eta(n)(\log n)n^{-s}-\sum_{k\ge0}(s+\kappa+2k)^{-2}+\delta(s^{-2}+(s-1)^{-2})$.
At $s=\tfrac12+c+iy$, $c>\tfrac12$, the function $(c+i(y-\gamma))^{-2}$ is the Fourier transform of
$\mathbf 1_{t>0}te^{-ct}e^{i\gamma t}$; testing against $e^{ct}\phi(t)/t$ with
$\phi\in C_c^\infty(0,\infty)$, and using
$\sum_\gamma m_\eta(\gamma)/(c^2+(y-\gamma)^2)=O(\log(2+\abs y))$, the exchanges are justified, and
dividing by $te^{-ct}$ gives \eqref{zer:eq:atoms}.

\begin{proposition}[Atom certificate; \stG]\label{zer:prop:atoms}
Let $J\subset\Gamma_q$ be finite. For each coordinate $l=(j,a)$ take an isolated probe
$p_l=p_a$ as in \cref{zer:lem:probes}, a bounded interval $I_l$ about
$t_l=\log p_l-h_j$ containing no other atom location of the translates
$U_{\chi^*}(\cdot+h_k)$ and with $t+h_k>0$ on it, and $\phi_l\in C_c^\infty(I_l)$ with
$\phi_l(t_l)=1$ that annihilates the functions $g_{\chi^*}(t+h_k)$ and $e^{i\gamma t}$,
$\gamma\in J\cup\set0$. Such tests exist, the numbers
\[
 K_l=\frac{p_l}{(\log p_l)^2}\sum_{\gamma\in\Gamma_q\setminus J}(1/4+\gamma^2)^2
 \Bigl|\int\phi_l(t)e^{i\gamma t}\dd t\Bigr|^2
\]
are finite, and
\[
 \lambda_{\min}\Bigl(\sum_{\gamma\in\Gamma_q\setminus(J\cup\set0)}z(\gamma)z(\gamma)^*\Bigr)\ge\Bigl(\sum_lK_l\Bigr)^{-1}.
\]
\end{proposition}
\begin{proof}
Existence: if evaluation at $t_l$ vanished on the common kernel of the finitely many
functionals $\phi\mapsto\int\phi f$, it would be a combination of them, and a point mass would be
a smooth function on $I_l$. For $v\in\C^D$ put $d_\gamma=\sum_lv_lz(\gamma)_l$ and
$c_{j,\chi}=n_q^{-1}\sum_av_{j,a}\overline{\chi(a)}$. The distribution
$\sum_{j,\chi}c_{j,\chi}U_{\chi^*}(t+h_j)$ has coefficient $(\tfrac14+\gamma^2)d_\gamma$ at each
$\gamma\ne0$. Pairing it with $\phi_l$ removes the frequencies in $J\cup\set0$ and the smooth part
of \eqref{zer:eq:atoms}, and leaves only the isolated atom, whose coefficient is
$-(\log p_l)p_l^{-1/2}\sum_\chi c_{j,\chi}\chi^*(p_l)=-(\log p_l)p_l^{-1/2}v_l$. Cauchy--Schwarz gives
$\abs{v_l}^2\le K_l\sum_{\gamma\notin J}\abs{d_\gamma}^2$; summing over $l$ bounds the Rayleigh
quotient. Finiteness of $K_l$ follows from the zero count and the rapid decay of
$\int\phi_le^{i\gamma t}\dd t$.
\end{proof}

\paragraph{Derivatives.}
Let $\chi_1,\dots,\chi_r$ be distinct primitive characters as in \cref{zer:thm:polyrigid}, with
$\delta_j=1$ for conductor one. For $t\ge0$ and integers $k\ge1$ define the logarithmic Riesz
means and their errors
\[
 S_{j,k}(t)=\frac1{k!}\sum_{n\le e^t}\Lambda(n)\chi_j(n)(t-\log n)^k,\qquad
 E_{j,k}(t)=e^{-t/2}\bigl(S_{j,k}(t)-\delta_je^t\bigr).
\]
Let $\mathcal G$ be the set of distinct ordinates of the family and fix $d\ge0$, $k\ge d+1$.

\begin{lemma}[Riesz-mean explicit formula; \stG]\label{zer:lem:riesz}
For $k\ge1$ and $t\ge0$,
$S_{j,k}(t)=\delta_je^t-\sum_\rho m_j(\rho)e^{\rho t}\rho^{-k-1}+O_{j,k}((1+t)^{k+1})$. Consequently,
for $0\le\ell\le d$,
\begin{equation}\label{zer:eq:jetexp}
 E_{j,k}^{(\ell)}(t)=-\sum_{\gamma\in\mathcal G}\frac{m_j(\gamma)(i\gamma)^\ell}{(1/2+i\gamma)^{k+1}}e^{i\gamma t}
 +O_{j,k,d}\bigl((1+t)^{k+1}e^{-t/2}\bigr),
\end{equation}
with absolutely and uniformly convergent series ($0^0=1$).
\end{lemma}
\begin{proof}
The Laplace transform of $S_{j,k}$ is $-L'(s,\chi_j)/(s^{k+1}L(s,\chi_j))$ for $\Re s>1$. Write
$-L'/L$ through \eqref{zer:eq:hadamard} and the digamma series: it is a constant, plus
$\delta/(s-1)$, plus a multiple of $1/s$, minus $\sum_\rho m(\rho)(1/(s-\rho)+1/\rho)$, minus
$\sum_{b_n>0}(1/(s+b_n)-1/b_n)$ with $b_n=a+2n$. For $z\ne0$ the inverse Laplace transform of
$(1/(s-z)+1/z)s^{-k-1}$ is $\bigl(e^{zt}-\sum_{h<k}(zt)^h/h!\bigr)z^{-k-1}$. Since
$\sum\abs z^{-2}$ converges over both families of $z$, and $\Re z\le\tfrac12$, termwise inversion
is justified by uniqueness of Laplace transforms of continuous functions of exponential
order. The pole terms give $\delta(e^t-\sum_{h\le k}t^h/h!)$ and polynomials of degree at most
$k+1$, the polynomial parts are $O((1+t)^{k-1})$ and the trivial zeros are bounded. Since
$S_{j,k}'=S_{j,k-1}$ for $k\ge2$, one has
$E_{j,k}^{(\ell)}=\sum_{h\le\ell}\binom\ell h(-\tfrac12)^{\ell-h}E_{j,k-h}$; inserting the formula
at each order $k-h\ge1$ turns the zero numerator into $(\rho-\tfrac12)^\ell=(i\gamma)^\ell$ by
\cref{src:thm:dgrh}, and $k+1-\ell\ge2$ gives absolute convergence.
\end{proof}

For a cut $H\ge0$ remove the low frequencies completely:
\[
 Y_{j\ell,H}(t)=E_{j,k}^{(\ell)}(t)+\sum_{\abs\gamma\le H}
 \frac{m_j(\gamma)(i\gamma)^\ell}{(1/2+i\gamma)^{k+1}}e^{i\gamma t},
\]
and stack these coordinates in $Y_H(t)\in\C^N$, $N=r(d+1)$.

\begin{theorem}[High-pass jet covariance; \stG]\label{zer:thm:jets}
For every finite family of distinct primitive characters, every $d\ge0$, $k\ge d+1$ and
$H\ge0$, uniformly in $u\ge0$,
\[
 \lim_{L\to\infty}\frac1L\int_u^{u+L}Y_H(t)Y_H(t)^*\dd t
 =\mathcal G_H:=\sum_{\abs\gamma>H}\frac{v(\gamma)v(\gamma)^*}{\abs{1/2+i\gamma}^{2k+2}},
\]
with $v(\gamma)$ as in \cref{zer:cor:tail}, and $\mathcal G_H$ is positive definite. The mean of
$Y_H$ tends to zero in the same sense. The same holds if any finite set of ordinates, zero
included, is removed instead of a full band.
\end{theorem}
\begin{proof}
By \eqref{zer:eq:jetexp}, $Y_H$ is an absolutely convergent series with frequencies $\abs\gamma>H$
plus a bounded integrable remainder, so the complex analogue of \cref{zer:lem:ap}(i) gives the
limits. For $c\ne0$, $c^*\mathcal G_Hc=\sum_{\abs\gamma>H}\abs{c^*v(\gamma)}^2\abs{1/2+i\gamma}^{-2k-2}>0$ by
\cref{zer:cor:tail}. A finite deletion leaves all sufficiently high frequencies.
\end{proof}

In particular, for coefficients $c_{j\ell}$ not all zero and any finite trigonometric polynomial
$P(t)$ with real frequencies, the mean square of $\sum c_{j\ell}E_{j,k}^{(\ell)}-P$ has a positive
limit: choose $H$ above the frequencies of $P$. No finite set of frequencies explains a
nonzero combination of the smoothed errors and their derivatives. If $\gamma_1,\dots,\gamma_N$
are tail ordinates whose vectors form a basis and $W$ is the matrix with columns
$v(\gamma_h)(1/2+i\gamma_h)^{-k-1}$, then $\mathcal G_H\succeq WW^*$ and
\begin{equation}\label{zer:eq:jetcert}
 \lambda_{\min}(\mathcal G_H)\ge\frac{\abs{\det W}^2}{(\tr WW^*)^{N-1}}>0,
\end{equation}
since the product of the $N-1$ largest eigenvalues of $WW^*$ is at most $(\tr WW^*)^{N-1}$;
computing this bound needs certified multiplicities.

For residue classes take the characters modulo $q$ with their distinct primitive inducing
characters, and for $b\in G$ put
\[
 E^{[q]}_{b,k}(t)=e^{-t/2}\Bigl(\frac1{k!}\sum_{n\le e^t,\ n\equiv b}\Lambda(n)(t-\log n)^k-\frac{e^t}{n_q}\Bigr).
\]
Passing to primitive characters changes $S_k$ by $O_q((1+t)^{k+1})$, since each prime $p\mid q$
has at most $t/\log p$ powers of weight at most $t^k\log p/k!$. Let $Z_H(t)\in\R^{n_q(d+1)}$ be the
vector of $(E^{[q]}_{b,k})^{(\ell)}$ with the low-frequency terms
$n_q^{-1}\sum_\chi\overline{\chi(b)}\sum_{\abs\gamma\le H}m_{\chi^*}(\gamma)(i\gamma)^\ell(1/2+i\gamma)^{-k-1}e^{i\gamma t}$
added; it is real, because conjugate characters have reflected conjugate terms and the cut is
symmetric. Up to a decaying error it is $MY_H$ with $M=(\overline{\chi(b)}/n_q)\otimes I_{d+1}$ and
$MM^*=I/n_q$, so its covariance $M\mathcal G_HM^*$ is real symmetric positive definite.

\begin{corollary}[Uniform directional oscillation; \stG]\label{zer:cor:jetosc}
For fixed $q,k,d,H$ as above there are $\eta>0$ and $T_0$ such that for every $u\ge0$, every
$T\ge T_0$ and every real unit vector $c$, there are $t_\pm\in[u,u+T]$ with
$\pm c^{\mathsf T}Z_H(t_\pm)\ge\eta$. Equivalently, the convex hull of
$\set{Z_H(t):u\le t\le u+T}$ contains the closed ball of radius $\eta$ about $0$.
\end{corollary}
\begin{proof}
Apply \cref{zer:lem:ap}(iii); the remainder in \eqref{zer:eq:jetexp} is bounded and
integrable on $[0,\infty)$. The converse direction of the equivalence follows by maximizing a
linear functional over the hull.
\end{proof}

For $d=0$ and $H=0$ every nonzero real contrast of centred smoothed residue-class counts
changes sign with bounded waiting time on the logarithmic scale. A covariance statement does
not say that every ordering of the unsmoothed counts occurs, nor that the limiting law has full
support: a distribution on a curve can have positive definite covariance and miss open sets.

\paragraph{A filter computed from primes alone.}
The constants $c_a$ and $\beta_a$ require $L'/L(1,\chi)$ and the central orders. A compact filter
removes them. Fix $h>0$ and put
\[
 (A_hf)(t)=\frac1h\int_0^hf(t-u)\,\dd u,\qquad d_h=\frac{2(1-e^{-h/2})}h\in(0,1),
\]
\[
 \mathcal D_a(t)=(I-A_h)(d_hI-A_h)^2\bigl[e^{t/2}S_a(e^t)\bigr].
\]
Only prime powers up to $e^t$ enter $\mathcal D_a(t)$, and only the values on $[t-3h,t]$. The
operator sends $e^{zt}$ to $M_h(z)e^{zt}$ with
\[
 a_h(z)=\frac{1-e^{-hz}}{hz},\quad a_h(0)=1,\qquad M_h(z)=(1-a_h(z))(d_h-a_h(z))^2 .
\]
Since $M_h(0)=M_h(\tfrac12)=M_h'(\tfrac12)=0$ (note $a_h(\tfrac12)=d_h$), it annihilates $1$,
$e^{t/2}$ and $te^{t/2}$.

\begin{lemma}\label{zer:lem:multiplier}
For $h>0$ and real $\gamma\ne0$, $M_h(i\gamma)\ne0$.
\end{lemma}
\begin{proof}
$a_h(i\gamma)$ is the mean of $e^{-i\gamma u}$ over $[0,h]$, of modulus less than $1$, so
$a_h(i\gamma)\ne1$. If it is real, its imaginary part $-(1-\cos\gamma h)/(\gamma h)$ vanishes, so
$e^{i\gamma h}=1$ and $a_h(i\gamma)=0\ne d_h$.
\end{proof}

\begin{theorem}[Filtered prime races; \stG]\label{zer:thm:filtered}
For every $q\ge1$ and $h>0$ there are $\eps,\eta,L_0>0$ and $T_0$, depending only on $q,h$, such
that for every real unit $w\in\R^G$, every $T\ge T_0$ and every $L\ge L_0$,
\[
 \abs{\set{t\in[T,T+L]:w\cdot\mathcal D(t)\ge\eps}}\ge\eta L,\qquad
 \abs{\set{t\in[T,T+L]:w\cdot\mathcal D(t)\le-\eps}}\ge\eta L,
\]
and the convex hull of $\set{\mathcal D(t):T\le t\le T+L}$ contains the ball of radius $\eps$. The
same holds, with constants depending also on the deleted set, after subtracting from $\mathcal D$
the terms $M_h(i\gamma)v_a(\gamma)(\tfrac14+\gamma^2)^{-1}e^{i\gamma t}$ of any finite symmetric set of
nonzero ordinates.
\end{theorem}
\begin{proof}
By \cref{zer:lem:ramp},
$e^{t/2}S_a(e^t)=te^{t/2}/n_q+c_ae^{t/2}+\beta_a+\sum_{\gamma\ne0}v_a(\gamma)(\tfrac14+\gamma^2)^{-1}e^{i\gamma t}+O((1+t)e^{-t/2})$.
The filter has operator norm at most $2(1+d_h)^2$ on bounded functions and memory $3h$; it
removes the first three terms and multiplies each frequency by $M_h(i\gamma)$. So
$\mathcal D(t)=\sum_{\gamma\ne0}M_h(i\gamma)u_\gamma e^{i\gamma t}+O((1+t)e^{-t/2})$ with
$u_\gamma=(v_a(\gamma)/(\tfrac14+\gamma^2))_a$. If
$\sum_{\abs\gamma>H}\abs{M_h(i\gamma)}^2\abs{w\cdot u_\gamma}^2=0$ for a real $w$, then by
\cref{zer:lem:multiplier}
$\sum_\chi\bigl(\sum_aw_a\overline{\chi(a)}\bigr)m_\chi(\gamma)=0$ for $\abs\gamma>H$, and
\cref{zer:thm:polyrigid} with constant polynomials gives $w=0$. \Cref{zer:lem:ap}(ii) completes
the proof.
\end{proof}

For evaluation, $e^{t/2}S_a(e^t)=\sum_{n\equiv a}\Lambda(n)n^{-1/2}b(t-\log n)$ with
$b(v)=\mathbf 1_{v\ge0}(e^{v/2}-e^{-v/2})$, and the filter expands as
$d_h^2I-(d_h^2+2d_h)A_h+(1+2d_h)A_h^2-A_h^3$ with
$A_h^jb(v)=h^{-j}\sum_{k\le j}(-1)^k\binom jkb_j(v-kh)$, where $b_j$ is the $j$-fold primitive
of $b$ vanishing on $v<0$,
\[
 b_j(v)=2^j\Bigl(e^{v/2}-\sum_{k<j}\frac{(v/2)^k}{k!}\Bigr)
 -(-2)^j\Bigl(e^{-v/2}-\sum_{k<j}\frac{(-v/2)^k}{k!}\Bigr)\quad(v\ge0).
\]
For $v\ge3h$ the filtered kernel is exactly $-M_h(-\tfrac12)e^{-v/2}$, so each prime power far
below the cutoff contributes a decaying term and no growing terms cancel. Floating-point
samples for $q=1,4,5,8$ are in \artifact{grh_filtered_prime_races/experiment_output.txt}.

\subsection{Balancing all prime-power races in one multiplicative interval}\label{zer:sec:balancing}
\sources{grh_cm_four_terms_and_prime_bias; grh_prime_error_balancing; multiscale_prime_fluctuations}

Fix a positive integer $q$, put $G=(\Z/q\Z)^\times$ and $n_q=\abs G$, and let $\chi$ run
over the characters of $G$. Write $\chi^*$ for the primitive character inducing $\chi$ and
$m_\chi$ for the order of vanishing of $L(s,\chi^*)$ at $1/2$. For $a\in G$ define
\begin{equation}\label{zer:eq:racevec}
\begin{aligned}
 \psi(x;q,a)&=\sum_{n\le x,\ n\equiv a\bmod q}\Lambda(n),&
 b_a&=\frac2{n_q}\sum_\chi\overline{\chi(a)}\,m_\chi,\\
 D_a(x)&=\frac{\psi(x;q,a)-x/n_q}{\sqrt x}+b_a,&
 D(x)&=(D_a(x))_{a\in G}.
\end{aligned}
\end{equation}
The constants $b_a$ are real by conjugation. They incorporate possible central zeros, so
even characters with unknown central nonvanishing are included. The von Mangoldt weight
counts prime powers as well as primes.

\begin{theorem}[Fixed-interval balancing; \stG]\label{zer:thm:balancing}
There exist $C_q>1$ and $r_q>0$ such that, for every sufficiently large real $X$,
\begin{equation}\label{zer:eq:ballrace}
 \set{y\in\R^{n_q}:\norm y_2\le r_q}\ \subseteq\
 \zerconv\set{D(n):\ n\in\Z,\ X\le n\le C_qX}.
\end{equation}
Consequently every such $y$, including zero, is exactly a convex combination of at most
$n_q+1$ integer samples from that interval, with the same sample points and weights for
every residue class, and every unit linear combination of the $n_q$ coordinates attains
values at least $r_q$ and at most $-r_q$ in the interval.
\end{theorem}

The theorem gives existence with constants depending on the fixed modulus. It requires
neither simplicity of zeros nor rational independence of their ordinates. A ball in the
convex hull allows thresholds to be mixed; a prescribed ranking of raw prime counts asks
for all of its inequalities to hold at one threshold, which is a stronger demand (even a
circle has a convex hull containing points that the circle never visits). For the
single-threshold problem under hypotheses on the zero ordinates see
\cite{MartinNg2020Inclusive}.

The first step is a vector version of the local averaging in the source proof, where high
Laplace moments find a one-sided small value on a window (\code{IsPolar.eventually_moment_le}
and \code{IsPolar.exists_window_sub_mul_le} in
\code{New/Endgame/OneSidedCriterion.lean}). A common choice of positive weights replaces a
separate choice of a point for each direction.

\begin{lemma}[Laplace averaging; \stU]\label{zer:lem:laplace}
Let $E:[0,\infty)\to\R^n$ be measurable, with $\norm{E(v)}\le Me^{Av}$, $A\ge0$. Suppose
its coordinate Laplace transforms $\Phi(z)=\int_0^\infty E(v)e^{-zv}\,\dd v$, initially
defined for $\Re z>A$, extend holomorphically to $\Re z>0$ and to a neighborhood of $0$.
For every $\tau>0$ there is $\delta>0$ such that, for all large $T$, some probability
density supported on $[T,(1+\tau)T]$ has $E$-average of norm $O(e^{-\delta T})$. The same
estimate holds for some convex combination of at most $n+1$ values of $E$ in the interval.
\end{lemma}
\begin{proof}
Take $c=1+\tau/2$. Choose $\lambda>A$ sufficiently large as specified below, set
$k=\floor{\lambda cT}$, and let $V$ have density $\lambda^{k+1}v^ke^{-\lambda v}/k!$ on
$v>0$. The disk $\abs{z-\lambda}\le\lambda$ meets the imaginary axis only at $0$.
Compactness gives $R>\lambda$ with $\Phi$ holomorphic on a neighborhood of the closed
concentric disk of radius $R$. Cauchy's estimate yields
\[
 \norm{\mathbb E\,E(V)}=\lambda^{k+1}\norm{\Phi^{(k)}(\lambda)}/k!\ll\lambda(\lambda/R)^k .
\]
For the tails, exponential tilting gives
\[
 \mathbb E\bigl[e^{AV};\ V\notin I\bigr]=\Bigl(\frac\lambda{\lambda-A}\Bigr)^{k+1}
 \Pr(W\notin I),\qquad I=[T,(1+\tau)T],
\]
where $W$ has the same shape $k+1$ and rate $\lambda-A$. Its mean divided by $T$ tends to
$c\lambda/(\lambda-A)$. For all sufficiently large $\lambda$ this limit stays a fixed
positive distance from both endpoints $1,1+\tau$. The Chernoff bound for a gamma variable
at $y$ times its mean is $\exp[-(k+1)(y-1-\log y)]$ for the corresponding tail, by
optimizing its moment generating function. The two endpoint ratios stay bounded away from
$1$, so the probability is at most $2e^{-\kappa_\tau(k+1)}$ with $\kappa_\tau>0$
independent of all sufficiently large $\lambda,T$. Choose $\lambda$ large enough that
$\log(\lambda/(\lambda-A))<\kappa_\tau/2$. Truncating the density to $I$ thus changes its
vector integral by $O(e^{-\delta T})$, and its mass is $1-O(e^{-\delta T})$, after
decreasing $\delta$ if necessary. The normalized integral belongs to the closed convex
hull. Approximate it to within $e^{-\delta T}$ by a finite convex combination and apply
Carath\'eodory's theorem.
\end{proof}

\begin{corollary}[Exponentially accurate balance in log-windows; \stG]\label{zer:cor:expbalance}
For every $\tau>0$ there is $\delta>0$ such that for all large $T$ some convex combination
of at most $n_q+1$ values $D(e^v)$, $v\in[T,(1+\tau)T]$, has norm $O_q(e^{-\delta T})$.
\end{corollary}
\begin{proof}
For the character coordinate of \eqref{zer:eq:racevec}, write
$\psi(x,\chi)=\sum_{n\le x}\chi(n)\Lambda(n)$, $e_\chi=1$ for the principal character and
$0$ otherwise, and $s=z+1/2$. The Laplace transform of
$e^{-v/2}(\psi(e^v,\chi)-e_\chi e^v)+2m_\chi$ is
\[
 -\frac1s\frac{L'}L(s,\chi)-\frac{e_\chi}{s-1}+\frac{2m_\chi}z .
\]
The pole at $s=1$ cancels for the principal character. The residue at $z=0$ before adding
the final term is $-2m_\chi$. Finite Euler factors removed from primitive $L$-functions
have zeros only on $\Re s=0$. \Cref{src:thm:dgrh} therefore gives the continuation required
by \cref{zer:lem:laplace}, and the finite Fourier transform on $G$ gives all the coordinates
in \eqref{zer:eq:racevec}.
\end{proof}

For exact balance in a log-window of fixed length the individual nonzero frequencies are
used, through \cref{zer:thm:rank}.

\begin{proof}[Proof of \cref{zer:thm:balancing}]
For a nonzero ordinate $\gamma$ in the union of the spectra put
$\rho_\gamma=\tfrac12+i\gamma$ and
\[
 c_{\gamma,a}=-\frac1{n_q\rho_\gamma}\sum_\chi\overline{\chi(a)}\,M(\rho_\gamma)_\chi,
 \qquad c_\gamma=(c_{\gamma,a})_{a\in G},
\]
with $M(\rho)$ as in \cref{zer:thm:rank}. The finite Fourier transform and
\cref{zer:thm:rank} show that the vectors $c_\gamma$ for positive $\gamma$ span
$\C^{n_q}$. Thus the real and imaginary parts of finitely many of them span $\R^{n_q}$.

On $v>1$, the explicit formula, as an identity of distributions, reads
\begin{equation}\label{zer:eq:explicitvec}
 D(e^v)=\sum_{\gamma\ne0}c_\gamma e^{i\gamma v}+O_q\bigl((1+v)e^{-v/2}\bigr).
\end{equation}
To see the error term, first use a primitive character. The main pole gives $e_\chi e^v$,
each nontrivial zero contributes $-e^{\rho v}/\rho$ with multiplicity, and the central
contribution is canceled by the constants in \eqref{zer:eq:racevec}. The pole of the
Perron kernel at zero and a possible trivial zero there contribute $O_q(1+v)$ before
normalization. The other trivial-zero terms form a geometrically decreasing series for
$v>1$. Replacing a primitive character by the one modulo $q$ changes $\psi$ by
$O_q(v+1)$, since only powers of primes dividing $q$ are affected. The choice at a jump is
immaterial to integration. For the boundary-value formula and its convergence for smooth
averages see \cite[Theorem~1(d),(e)]{KaczorowskiRamare2003Almost}.

Take a smooth probability density $w$, positive on $(0,1)$ and supported on $[0,1]$, and
let $w_n$ be its $n$-fold convolution, supported on $[0,n]$. Put
$\widehat w(t)=\int w(y)e^{ity}\,\dd y$. For each nonzero real $t$,
$\abs{\widehat w(t)}<1$, and $\widehat w$ decreases rapidly. For a fixed $\omega$,
integration of \eqref{zer:eq:explicitvec} shows
\begin{multline*}
 \int_0^nw_n(y)\,D(e^{V+y})\,e^{-i\omega(V+y)}\,\dd y\\
 =\sum_{\gamma\ne0}c_\gamma e^{i(\gamma-\omega)V}\,\widehat w(\gamma-\omega)^n
 +O_q\bigl((1+V)e^{-V/2}\bigr)\qquad(V\ge2).
\end{multline*}
The error is uniform in $n$ since $(1+v)e^{-v/2}$ decreases for $v\ge2$. The sum converges
absolutely for $n\ge1$ by rapid decay and the classical $O_q(Y\log(2Y))$ zero count up to
height $Y$. The union of the spectra is locally finite. Hence for each fixed $\omega$ the
supremum of $\abs{\widehat w(\gamma-\omega)}$ over $\gamma\ne\omega$ is strictly below $1$.
All off-frequency terms tend uniformly in $V$ to zero as $n\to\infty$. At $\omega=0$ the
sum therefore tends to zero. At a selected $\omega=\gamma$ it tends to $c_\gamma$. For the
arithmetic remainder take first a sufficiently large fixed $n$, then $V$ sufficiently
large.

Average $D(e^v)$ against the nonnegative densities proportional to
$w_n(v-V)(1\pm\cos(\gamma v))$ and $w_n(v-V)(1\pm\sin(\gamma v))$. Their masses tend
uniformly to $1$ because $\abs{\widehat w(\gamma)}<1$. The averages approach, to any
prescribed accuracy, the vertices $\pm\Re c_\gamma$ and $\pm\Im c_\gamma$ for the selected
finite set of positive ordinates (the opposite signs for sine merely interchange the two
vertices). The convex hull of these ideal vertices contains a ball about the origin because
their real span has full dimension. Concretely, if its inradius is $r>0$, approximating
every vertex within $r/4$ preserves a ball of radius $3r/4$, by the support-function
characterization of convex inclusion.

Finally replace $D(e^v)$ in the averages by $D(\floor{e^v})$. Between consecutive integers
the arithmetic sum in the numerator is constant, and $\psi(x;q,a)\le x\log(2x)$. The mean
value theorem bounds the change by $O_q((1+v)e^{-v/2})$, uniformly for large $v$. The
averages now belong to the convex hull of a finite set of integer samples. Set
$V=\log(2X)$, so those integers lie in $[X,2e^nX]$. Taking $C_q=2e^n$ and $r_q=r/2$ proves
\eqref{zer:eq:ballrace}, and Carath\'eodory's theorem gives the sample bound.
\end{proof}

\begin{remark}\label{zer:rem:stepanov}
The conditional almost-periodicity theory of \cite[Theorem~1(e)]{KaczorowskiRamare2003Almost}
gives a second perspective. The corrected vector has Stepanov almost-periodic coordinates,
zero mean after removing the central terms, and the nonzero Fourier coefficients displayed
in \eqref{zer:eq:explicitvec}, up to an overall sign. \Cref{zer:thm:rank} makes every
nonzero real projection nonconstant. Positive and negative values of a nonzero mean-zero
projection then recur on sets of positive measure in intervals of bounded log-length: an
interval on which it exceeds a positive level on a positive-measure subset can be repeated
using relatively dense Stepanov almost-periods, and the same applies to its negative. The
positive compact averages in the proof give a direct argument uniform across directions.
\end{remark}

\begin{example}[\stU]\label{zer:ex:triangle}
Take $q=4$. Its two central orders vanish: the nonprincipal central value is a positive
alternating series, and for the principal character $\zeta(1/2)=\eta(1/2)/(1-\sqrt2)$,
where the alternating eta series is positive. The three integers $293,283,228$ lie in
$[200,400]$. For $v(n)=(\psi(n;4,1)-n/2,\ \psi(n;4,3)-n/2)$ the certified lower bounds for
the cyclic determinants are
\[
 \det(v(293),v(283))\ge5.625,\quad\det(v(283),v(228))\ge41.693,\quad
 \det(v(228),v(293))\ge14.050 .
\]
All three are positive, and the identity
$\det(v_2,v_3)v_1+\det(v_3,v_1)v_2+\det(v_1,v_2)v_3=0$ therefore gives strictly positive
balancing weights. On replacing $v(n)$ by $D(n)=v(n)/\sqrt n$, multiply the associated
weight by $\sqrt n$ and renormalize. This example is independent of GRH. The three bounds
are verified with integer sieving and rational estimates for logarithms in
\begin{center}
\artifact{grh_cm_four_terms_and_prime_bias/checks/prime_race_triangle.py}
\end{center}
\end{example}

\paragraph{Smoothed sums in bounded logarithmic windows.}
For the ramp sums of \cref{zer:sec:cov} the balance holds at real scales in windows of bounded
logarithmic length centred anywhere. With the notation \eqref{zer:eq:va} put
$R_a(x)=S_a(x)-\log x/n_q-c_a$, $F(v)=(e^{v/2}R_a(e^v))_{a\in G}$ and $\beta=(\beta_a)_{a\in G}$.

\begin{corollary}[Uniform balancing of smoothed errors; \stG]\label{zer:cor:smoothbalance}
For each $q$ there are $H>0$, $r>0$ and $T_0$ such that for every real $T\ge T_0$
\[
 \beta+\set{y\in\R^{n_q}:\norm y_2\le r}\ \subseteq\ \zerconv\set{F(v):T-H\le v\le T+H}.
\]
Hence every target $y$ with $\norm{y-\beta}\le r$, in particular $y=\beta$, equals
$\sum_{j\le k}w_j\sqrt{x_j}R_a(x_j)$ simultaneously for all $a\in G$, for some $k\le n_q+1$,
scales $x_j\in[e^{T-H},e^{T+H}]$ and convex weights $w_j$.
\end{corollary}
\begin{proof}
$F(v)-\beta=Z(v)$. Apply \cref{zer:cor:multiscale} with $r=1$, $h_1=0$ and $J=\varnothing$,
and \cref{zer:lem:ap}(ii) with windows of length $2H\ge L_0$; finish with Carath\'eodory's
theorem.
\end{proof}

The scales and weights depend on the centre and the target, and no single scale balances all
classes; for $q=1$ the intermediate value theorem gives one scale for every target in the
ball. This ball concerns the smoothed sums at real scales and \cref{zer:thm:balancing} the
unsmoothed $\psi(x;q,a)$ at integers; neither statement implies the other. The radius can be
certified from finite data: for a finite symmetric set $E$ of nonzero ordinates with tail
$\tau=\sum_{\gamma\notin E}\norm{c(\gamma)}$, where $c(\gamma)=(v_a(\gamma)/(\tfrac14+\gamma^2))_a$ and
$K=\sum_\gamma\norm{c(\gamma)}$, the uniform error of the window mean is at most
$\tau+H^{-1}\sum_{\gamma\in E}\norm{c(\gamma)}/\abs\gamma$, and that of the second moment is at most
$3K\tau+H^{-1}\sum_{\gamma,\eta\in E,\ \gamma+\eta\ne0}\norm{c(\gamma)}\norm{c(\eta)}/\abs{\gamma+\eta}$; any
certified bound for $K$ and positive lower bound for the least eigenvalue may be inserted in
the proof of \cref{zer:lem:ap}.

A different averaging rule is independent of the centre and of the target.

\begin{theorem}[Common compact kernel; \stG]\label{zer:thm:kernel}
For $B\ge2$ let $\mathcal Z_B$ be the multiset of the absolute values $0<\abs\gamma\le B$ of the
ordinates of all $L(s,\chi^*)$, $\chi\bmod q$, with multiplicity and both signs, $N_B=\abs{\mathcal Z_B}$,
\[
 H_B=\pi\sum_{g\in\mathcal Z_B}\frac1g,\qquad
 \mu_B=\mathop{\ast}_{g\in\mathcal Z_B}\operatorname{Unif}[-\pi/g,\pi/g].
\]
Then $\mu_B$ is a probability measure on $[-H_B,H_B]$, and for $T\ge H_B$ and fixed
$0<\delta<\tfrac12$,
\[
 \Bigl\|\int F(T+w)\,\dd\mu_B(w)-\beta\Bigr\|_2\le K\pi^{-N_B}+C_{q,\delta}e^{-(1/2-\delta)(T-H_B)}.
\]
Moreover $H_B=O_q(1+\log^2B)$ and $N_B\gg B\log B$. With $B=T$, for every fixed
$0<\eps<\tfrac12$,
\[
 \Bigl\|\int F(T+w)\,\dd\mu_T(w)-\beta\Bigr\|_2=O_{q,\eps}\bigl(e^{-(1/2-\eps)T}\bigr),\qquad
 H_T=O_q(\log^2T).
\]
\end{theorem}
\begin{proof}
Convolution adds support radii and preserves mass, and
\[
 \widehat\mu_B(\omega)=\prod_{g\in\mathcal Z_B}\frac{\sin(\pi\omega/g)}{\pi\omega/g}.
\]
This vanishes at every
nonzero ordinate of modulus at most $B$, and for $\abs\omega>B$ each factor has modulus at most
$g/(\pi\abs\omega)\le1/\pi$. Integrating \eqref{zer:eq:Zexp} termwise, the high frequencies
contribute at most $K\pi^{-N_B}$, and the remainder is bounded on the support. The zero count
$O_q(B\log(B+2))$ and partial summation give $\sum_{g\le B}1/g=O_q(1+\log^2B)$, with a fixed
positive minimum among the finitely many small ordinates; the zeta factor and the
Riemann--von Mangoldt formula give $N_B\gg B\log B$. For $B=T$ the first term is
$\exp(-cT\log T)$; take $\delta=\eps/2$ and absorb $\exp((\tfrac12-\delta)H_T)$ into
$\exp(\eps T/2)$.
\end{proof}

Writing $X=e^T$ and $R=(R_a)_a$, the same kernel gives
\[
 \Bigl\|\int\sqrt{e^w}\,R(Xe^w)\,\dd\mu_T(w)-X^{-1/2}\beta\Bigr\|_2=O_{q,\eps}(X^{-1+\eps}),
\]
with all scales in $[X\exp(-O_q((\log\log X)^2)),X\exp(O_q((\log\log X)^2))]$. The weight
$\sqrt{e^w}$ matters, and this is a positive averaging statement, not a pointwise bound of size
$X^{-1+\eps}$. For a spectral accuracy $\eps$, ignoring the decaying remainder, $B\asymp\log(1/\eps)$
suffices, with logarithmic support radius $O_q((\log\log(1/\eps))^2)$.

\begin{remark}[The remainder of the source proof]\label{zer:rem:sourceremainder}
Let $d\equiv7\pmod8$ be prime, $(M,d)=1$, and let $R_{d,M}$ be the centred remainder of
$\sum_{n\le x,\ n\equiv1\bmod M}\Lambda(n)(1-\chi_d(n))(1/n-1/x)$, whose main term is
$\log x/\varphi(M)$; its constant is the regular part at $1$ of the weighted logarithmic
derivative minus the pole coefficient, as in the source's \code{logDerivMean}. Take $q=Md$ and
$w(a)=\mathbf 1_{a\equiv1\bmod M}(1-\chi_d(a))$ on $(\Z/q\Z)^\times$; then $\sum_aw(a)=d-1$, so
$\sum_aw(a)/n_q=1/\varphi(M)$. The powers of $d$ add
$\log d\sum_{d^k\le x,\ d^k\equiv1\bmod M}(d^{-k}-x^{-1})$, a constant plus $O_{d,M}((1+\log x)/x)$, so
\[
 \sqrt xR_{d,M}(x)=\sum_aw(a)F_a(\log x)+O_{d,M}\bigl((1+\log x)/\sqrt x\bigr).
\]
Since $w\ne0$, \cref{zer:cor:multiscale} applies to the shifted copies
$(\sqrt{e^{h_j}x}\,R_{d,M}(e^{h_j}x)-\sum_aw(a)\beta_a)_j$: their covariance is positive definite,
every nonzero real combination has both signs on sets of positive lower logarithmic density,
and the centred remainder of the source has infinitely many spectral frequencies and is not a
linear combination of finitely many other scales. If $\varphi(M)>1$, the weights of $R_{d,M}$ and
of $R_{d,1}$ (computed on the same modulus, with the finitely many powers of primes dividing $M$
added back) are linearly independent, so the pair $(\sqrt xR_{d,M},\sqrt xR_{d,1})$ has a
two-dimensional ball about its central biases in every sufficiently late window of the fixed
logarithmic length of \cref{zer:cor:smoothbalance}, and both biases are attained by one convex
average of at most three scales.
\end{remark}

\subsection{Positive spectral weights: peaks, sharp windows and resonance}\label{zer:sec:peaks}
{\emergencystretch=6em\sources{prime_peak_rigidity; sparse_grh_balanced_windows; sharp_multiscale_prime_windows; grh_prime_resonance_2026-11-26}}

On the critical line $1/(\rho(1-\rho))=1/(\tfrac14+\gamma^2)$ is positive. In the identity class,
and for suitably balanced or autocorrelated windows, all spectral coefficients of a smoothed
prime error are then nonnegative. Such a function is bounded by its value at $0$, which is an
arithmetic constant, and returns arbitrarily close to it. This gives exact extremal values,
rigidity near the extremum, sharp inequalities between windows, and a covariance kernel whose
local structure is arithmetic. We keep the notation of \cref{zer:sec:cov}; $K_\chi$ depends only
on $\chi^*$.

\begin{lemma}[Positive spectra; \stU]\label{zer:lem:posspec}
\begin{enumerate}[label=\textup{(\roman*)}]
\item Let $F_1,\dots,F_k:\R\to\C^N$ be absolutely convergent trigonometric series and $\eps>0$. The
set of $\tau$ with $\sup_{u\in\R}\norm{F_i(u+\tau)-F_i(u)}<\eps$ for all $i$ is relatively dense
and unbounded above.
\item Let $f(t)=\sum_\gamma c_\gamma e^{i\gamma t}$ with $c_\gamma\ge0$ and $\sum c_\gamma<\infty$.
Then $\abs{f(t)}\le f(0)$ for all $t$. If $t_j\to\infty$ and $t_{j+1}-t_j\to0$, then
$\limsup_j\Re f(t_j)=f(0)$.
\item Let $\widehat G$ be the characters of a finite abelian group $G$, let
$w_{\chi,\gamma}\ge0$ be summable, and put
\[
 \mathcal K(a,t)=\sum_{\chi\in\widehat G}\sum_\gamma w_{\chi,\gamma}\overline{\chi(a)}e^{i\gamma t}.
\]
For $(a_i,t_i)\in G\times\R$ and $c_i\in\C$ put
\[
 Q=\sum_{i,j}c_i\overline{c_j}\mathcal K(a_ia_j^{-1},t_i-t_j),\qquad
 \mathcal S(T)=\sum_{i,j}c_i\overline{c_j}\mathcal K(a_ia_j^{-1},T+t_i-t_j).
\]
Then $Q\ge0$, $\abs{\mathcal S(T)}\le Q$, and $\limsup_{T\to\infty}\abs{\mathcal S(T)}=Q$.
\end{enumerate}
\end{lemma}
\begin{proof}
(i) Discard tails of total coefficient norm below $\eps/4$. For the finitely many remaining
frequencies, the closure of $\set{(e^{i\gamma t})_\gamma:t\in\R}$ is a compact subgroup $\mathcal T$
of a torus. For a symmetric neighbourhood $U$ of the identity, the translates by the dense
orbit cover $\mathcal T$, so finitely many, by $s_1,\dots,s_m\in[A,B]$, already do; for every $t$
some $t-s_j$ has its phases in $U$, so returns occur in every interval of length $B-A+1$, and
arbitrarily far to the right. No independence of the frequencies is used. (ii) The bound is
the triangle inequality. By (i) there are arbitrarily large $\tau$ with $\abs{f(\tau)-f(0)}$
small, and uniform continuity transfers this to a mesh with gaps tending to zero.
(iii) $\mathcal S(T)=\sum_{\chi,\gamma}w_{\chi,\gamma}e^{i\gamma T}\abs{\sum_ic_i\overline{\chi(a_i)}e^{i\gamma t_i}}^2$
has nonnegative coefficients and $\mathcal S(0)=Q$; apply (ii).
\end{proof}

\begin{lemma}[Mass of the zero spectrum; \stG]\label{zer:lem:mass}
For a primitive character $\chi$ of conductor $f$ and parity $b$,
\[
 M_\chi:=K_\chi(0)=\sum_\gamma\frac{m_\chi(\gamma)}{1/4+\gamma^2}=
 \begin{cases}2+\zergE-\log(4\pi),&f=1,\\[2pt]
 \log(f/\pi)+\psi\bigl((1+b)/2\bigr)+2\Re\dfrac{L'}{L}(1,\chi),&f>1,\end{cases}
\]
and $M_\chi>0$. For $f=1$ and $t\ge0$ one has exactly
\begin{equation}\label{zer:eq:zetaexact}
 K_\zeta(t)=e^{t/2}\Bigl[S(e^t)-t+1+\zergE-e^{-t}\log(2\pi)+\mathcal T(e^{-t})\Bigr],\qquad
 \mathcal T(z)=\sum_{k\ge1}\frac{z^{2k+1}}{2k(2k+1)},
\end{equation}
where $S(x)=\sum_{n\le x}\Lambda(n)(1/n-1/x)$ and
$\mathcal T(z)=z+\tfrac12(1-z)\log(1-z)-\tfrac12(1+z)\log(1+z)$.
\end{lemma}
\begin{proof}
For the completed entire function $\Xi_\chi$ the Hadamard product gives
$\Xi_\chi'/\Xi_\chi(1)-\Xi_\chi'/\Xi_\chi(0)=\sum_\rho(1/(1-\rho)+1/\rho)=M_\chi$, an absolutely convergent
sum of positive terms by \cref{src:thm:dgrh}, and the functional equation with conjugation makes
the left side $2\Re\Xi_\chi'/\Xi_\chi(1)$. Inserting the Gamma factors, with
$\xi'/\xi(1)=1+\zergE-\tfrac12\log\pi+\tfrac12\psi(\tfrac12)$ and $\psi(\tfrac12)=-\zergE-2\log2$, gives the
values (the same identity appears in the proof of \cref{ana:thm:excursion}). For
\eqref{zer:eq:zetaexact}, the contour argument of \cref{zer:lem:ramp} for $-\zeta'/\zeta$, carried to
the left past all trivial zeros, gives residues $\log x-1-\zergE$ at $1$,
$x^{\rho-1}/(\rho(1-\rho))$ at each zero, $x^{-1}\log(2\pi)$ at $0$ and $-x^{-2k-1}/(2k(2k+1))$ at $-2k$;
multiply by $x^{1/2}$. At $t=0$, $S(1)=0$ and $\mathcal T(1)=1-\log2$.
\end{proof}

\paragraph{Peak rigidity.}
Put $E_a(x)=\sqrt x\bigl(S_a(x)-\log x/n_q-c_a\bigr)$, so $E_a(e^t)=Z_a(t)+\beta_a$, and
\[
 B_a(t)=\frac1{n_q}\sum_\chi\overline{\chi(a)}K_\chi(t),\qquad v_a=B_a(0),\qquad D(t)=v_1-B_1(t),
\]
where $B_1$ is the identity class. By \cref{zer:lem:ramp},
$E_a(x)=B_a(\log x)+O_q((1+\log x)/\sqrt x)$. Conjugation shows that $B_a$ and $v_a$ are real and
$v_a=v_{a^{-1}}$; $v_1=n_q^{-1}\sum_\chi M_{\chi^*}>0$, and every $v_a$ is a finite Fourier transform
of the masses of \cref{zer:lem:mass}.

\begin{theorem}[Peak rigidity; \stG]\label{zer:thm:peak}
The following hold for every $q\ge1$.
\begin{enumerate}[label=\textup{(\roman*)}]
\item $\limsup_{x\to\infty}E_1(x)=v_1$, and $\limsup_{x\to\infty}\abs{E_a(x)}\le v_1$ for every
$a$, with $\abs{v_a}<v_1$ for $a\ne1$.
\item $D(t)\ge0$, and for all $a$ and all real $t,u$,
\[
 \abs{B_a(t+u)-B_a(u)}^2\le2v_1D(t).
\]
\item For all $a$ and $t$,
\[
 \Bigl|\frac{B_a(t)+B_{a^{-1}}(t)}2-v_a\Bigr|\le D(t),\qquad
 \frac1{n_q}\sum_\chi\frac{\abs{K_\chi(t)-M_{\chi^*}}^2}{M_{\chi^*}}\le2D(t).
\]
The matrices $V_{a,b}=v_{ab^{-1}}$ and $H(t)_{a,b}=\tfrac12(B_{ab^{-1}}(t)+B_{ba^{-1}}(t))$ satisfy
$V\succ0$ and $V-H(t)\succeq0$, and $\tr(V-H(t))=n_qD(t)$.
\item For every $\eps>0$ the set of $t$ with $\sup_{a,u}\abs{B_a(t+u)-B_a(u)}<\eps$ is relatively
dense. Hence the near-peak scales of $E_1$ are relatively dense on the logarithmic axis.
\end{enumerate}
\end{theorem}
\begin{proof}
With $w_\gamma=(\tfrac14+\gamma^2)^{-1}$ and multiplicities retained, exactly
\[
 D(t)=\frac1{n_q}\sum_\chi\sum_\gamma w_\gamma(1-\cos\gamma t)=\frac1{2n_q}\sum_\chi\sum_\gamma w_\gamma\abs{e^{i\gamma t}-1}^2\ge0 .
\]
The translate difference is $n_q^{-1}\sum_{\chi,\gamma}w_\gamma\overline{\chi(a)}e^{i\gamma u}(e^{i\gamma t}-1)$, and
Cauchy--Schwarz for the measure of mass $w_\gamma/n_q$, of total mass $v_1$, gives (ii). Put
$D_\chi(t)=M_{\chi^*}-\Re K_\chi(t)\ge0$, whose average is $D(t)$. Then
$\abs{M_{\chi^*}-K_\chi(t)}^2\le M_{\chi^*}\sum_\gamma w_\gamma\abs{1-e^{i\gamma t}}^2=2M_{\chi^*}D_\chi(t)$, and
$v_a-\tfrac12(B_a(t)+B_{a^{-1}}(t))=n_q^{-1}\sum_\chi\Re\chi(a)D_\chi(t)$, which give (iii); the
vectors $(\overline{\chi(a)})_a$ diagonalize $V$ and $H(t)$ with eigenvalues $M_{\chi^*}$ and
$\Re K_\chi(t)$. \Cref{zer:lem:posspec}(i) makes $D(t)$ small on a relatively dense set, which with
(ii) gives (iv), and then $\limsup E_1=v_1$. Since all $M_{\chi^*}>0$ and some character separates
$a\ne1$ from $1$, the triangle inequality is strict in $\abs{v_a}<v_1$.
\end{proof}

\begin{corollary}[A unique exposed peak; \stG]\label{zer:cor:peak}
For $x_j\to\infty$ the following are equivalent: $E_1(x_j)\to v_1$;
$(E_a(x_j))_a\to(v_a)_a$; $E_a(x_je^u)\to B_a(u)$ for every $a$, uniformly for
$u\ge-\tfrac12\log x_j$. The hyperplane where the identity coordinate equals $v_1$ meets the
cluster set of $(E_a(x))_a$ in one point. The same holds along integers $x$.
\end{corollary}
\begin{proof}
If $E_1(x_j)\to v_1$ then $D(\log x_j)\to0$, and (ii) gives uniformity in $u$; the error in
$E_a=B_a+O(\cdot)$ tends to zero uniformly for $x_je^u\ge\sqrt{x_j}$. Rounding $e^t$ to an integer
changes $t$ by $o(1)$, and the $B_a$ are uniformly continuous.
\end{proof}

\begin{theorem}[Universal simultaneous peaks; \stG]\label{zer:thm:univpeak}
There are integers $n_j\to\infty$ such that for every $q\ge1$ and $a\in(\Z/q\Z)^\times$,
$E_{q,a}(n_je^u)-B_{q,a}(u)\to0$ uniformly for $u\ge-\tfrac12\log n_j$. In particular
$E_{q,1}(n_j)\to v_{q,1}$ for every $q$. Let $F:\Z\to\C$ have period $Q$, put
$\mu_F=\varphi(Q)^{-1}\sum_{a\in G_Q}F(a)$,
$c_F=\sum_{a\in G_Q}F(a)c_{Q,a}+\sum_{p\mid Q}\log p\sum_{k\ge1}F(p^k)p^{-k}$, and
$E_F(x)=\sqrt x\bigl(\sum_{n\le x}\Lambda(n)F(n)(1/n-1/x)-\mu_F\log x-c_F\bigr)$. Then
$E_F(x)=B_F(\log x)+O_{F}((1+\log x)/\sqrt x)$ with $B_F=\sum_{a\in G_Q}F(a)B_{Q,a}$, and
$E_F(n_je^u)\to B_F(u)$ uniformly for $u\ge-\tfrac12\log n_j$.
\end{theorem}
\begin{proof}
For moduli $1,\dots,j$ take, by \cref{zer:lem:posspec}(i) applied to the finitely many series
$B_{q,a}$, a return $t_j$ with uniform translate errors below $1/(2j)$, large enough that the
arithmetic errors are below $1/j$ for $x\ge\tfrac12e^{t_j/2}$ and that rounding $e^{t_j}$ to an
integer $n_j>n_{j-1}$ moves all these series by less than $1/(2j)$. For a periodic weight,
nonunits $n$ with $\Lambda(n)\ne0$ are powers of primes dividing $Q$, which give the additional
constant and an error $O((1+\log x)/x)$; the unit classes follow by linearity.
\end{proof}

This includes the weight $\mathbf 1_{n\equiv1\bmod q}(1-\chi_d(n))$ of the source proof. For a signed
combination the recurring value $B_F(0)$ need not be the maximum, but the recurrence of the
whole profile holds. Conversely, an eventual bound $E_1(x)\le C$ already excludes zeros of the
$L(s,\chi)$, $\chi\bmod q$, from $\Re s>\tfrac12$, by \cref{ana:thm:zerotest} with all weights
$1/n_q$. Thus GRH for the characters modulo $q$ holds if and only if $\limsup E_1(x)<\infty$, and
then the limsup is $v_1$. For example $v_1\approx0.0461914179$ for $q=1$; for $q=3$,
$v_1\approx0.0797106939$ and $v_2\approx-0.0335192760$; and for $q=5$ the classes $1$, $2$,
$3$, $4$ give approximately $0.1522978093$, $-0.0275913840$, $-0.0275913840$ and
$-0.0509236233$. These decimals are
evaluated in floating point from \cref{zer:lem:mass} in \artifact{prime_peak_rigidity/check.py}.

\paragraph{A balanced window with a sharp normalized law.}
Let $b_\chi\ge0$ $(\chi\bmod q)$ be real, not all zero, with $b_{\overline\chi}=b_\chi$, and put
$w(n)=\sum_\chi b_\chi\chi(n)$, $\mu=b_{\chi_0}$, $S_w(x)=\sum_{n\le x}\Lambda(n)w(n)(1/n-1/x)$. For
$c>1$ let $\alpha=c/(c+1)$, $\alpha'=1/(c+1)$ and
\[
 \mathcal W_w(x;c)=S_w(x)-\alpha S_w(cx)-\alpha'S_w(x/c)+\mu(\alpha-\alpha')\log c\qquad(x\ge c),
\]
which equals $E_w(x)-\alpha E_w(cx)-\alpha'E_w(x/c)$ with $E_w=S_w-\mu\log x$; no constant of the
prime sum is needed. Its Mellin transform carries the factor
$P_c(s)=1-\alpha c^{s-1}-\alpha'c^{1-s}=(1-u)(cu-1)/((c+1)u)$, $u=c^{s-1}$, which vanishes only on
$\Re s\in\set{0,1}$ and on the critical line equals $1-\kappa_c\cos(\gamma\log c)>0$, where
$\kappa_c=2\sqrt c/(c+1)<1$.

\begin{theorem}[Sharp balanced-window law; \stG]\label{zer:thm:balanced}
Put
\[
 L_w(c)=\sum_\chi b_\chi\sum_\gamma m_\chi(\gamma)\frac{1-\kappa_c\cos(\gamma\log c)}{1/4+\gamma^2}>0 .
\]
For every strictly increasing $x_j\to\infty$ with $x_{j+1}/x_j\to1$,
\[
 \limsup_{j\to\infty}\sqrt{x_j}\,\mathcal W_w(x_j;c)=\limsup_{j\to\infty}\sqrt{x_j}\,\abs{\mathcal W_w(x_j;c)}=L_w(c),
\]
and the same value is the limsup over all real $x$.
\end{theorem}
\begin{proof}
By \cref{zer:lem:ramp}, $E_w(x)=\sum_\chi b_\chi c_\chi+x^{-1/2}f_w(\log x)+O((1+\log x)/x)$ with
$f_w=\sum_\chi b_\chi K_\chi$. Since $\alpha+\alpha'=1$ and $\alpha c^{-1/2}=\alpha'c^{1/2}=\kappa_c/2$,
\[
 \sqrt x\,\mathcal W_w(x;c)=f_w(t)-\tfrac{\kappa_c}2\bigl(f_w(t+\log c)+f_w(t-\log c)\bigr)
 +O\bigl((1+t)e^{-t/2}\bigr),\qquad t=\log x,
\]
a series with coefficients $b_\chi m_\chi(\gamma)(1-\kappa_c\cos(\gamma\log c))/(\tfrac14+\gamma^2)\ge0$ and value
$L_w(c)$ at $t=0$. Apply \cref{zer:lem:posspec}(ii); $\log x_j$ has gaps tending to zero.
\end{proof}

The positive weights are essential, and the statement uses neither simplicity nor
independence of ordinates. On a mesh with power-saving gaps, the sparse one-sided criterion of
\cref{sec:ana} shows conversely that a finite limsup forces all active zeros onto the critical
line, and that without GRH the limsup is $+\infty$. For $\zeta$ ($q=1$, $w=1$), \eqref{zer:eq:zetaexact}
evaluates the constant from one prime sum:
\[
 L_1(c)=M_\zeta-\kappa_cK_\zeta(\log c)
 =2+\zergE-\log(4\pi)-\frac{2c}{c+1}\Bigl(S(c)-\log c+1+\zergE-\frac{\log(2\pi)}c+\mathcal T(1/c)\Bigr),
\]
and $S(2)=0$ gives $L_1(2)=0.0625323180445\ldots$ (a numerical evaluation of this closed form).
Consequently, for every fixed integer $k\ge1$,
\[
 \limsup_{n\to\infty}n^{k/2}\Bigl(S(n^k)-\tfrac23S(2n^k)-\tfrac13S(n^k/2)+\tfrac13\log2\Bigr)=L_1(2).
\]
Since $P_c(0)=0$, the term $\log(2\pi)/x$ of \eqref{zer:eq:zetaexact} cancels exactly in
$\mathcal W_1$, so for all $x\ge c$ the corrected quantity
$\sqrt x\bigl(\mathcal W_1(x;c)+\mathcal T(1/x)-\alpha\mathcal T(1/(cx))-\alpha'\mathcal T(c/x)\bigr)$ is exactly the
positive-coefficient series at $\log x$ and has absolute value at most $L_1(c)$; the constant is
optimal on every mesh of the theorem.

\paragraph{Sharp inequalities between compact prime windows.}
Fix an odd prime $q$, $d=q-1$, $G=(\Z/q\Z)^\times$ and $0<h<\log2$. Define
\[
 g_h(u)=\begin{cases}e^{h-\abs u/2}-e^{\abs u/2},&\abs u\le h,\\0,&\abs u>h,\end{cases}\qquad
 \mathsf H=e^h-1,\qquad\mathsf B=h\mathsf H,
\]
and $b_j(u)=e^{-(j+1/2)u}/(1-e^{-2u})$ for $u>0$ and $j=0,1$.
\[
\]
For $a\in G$ and $t>h$ put
\begin{align*}
 P(a,t)&=\sum_{n\equiv a\bmod q}\frac{\Lambda(n)}{\sqrt n}g_h(\log n-t)+\frac1d\sum_{k\ge1}\frac{\log q}{q^{k/2}}g_h(k\log q-t),\\
 A(a,t)&=\frac12\int_{t-h}^{t+h}g_h(u-t)\bigl[\mathbf 1_{a=1}(b_0+b_1)(u)+\mathbf 1_{a=-1}(b_0-b_1)(u)\bigr]\dd u,\\
 C(t)&=\frac{\mathsf B}d\bigl(e^{t/2}+e^{-t/2}\bigr),\qquad \mathsf E(a,t)=C(t)-P(a,t)-A(a,t).
\end{align*}
The prime sum samples $e^{t-h}\le n\le e^{t+h}$; the second term comes from $\zeta$, the
primitive character inducing the principal one. With
$J_j=2\int_0^\infty(\mathsf H-g_h(u))b_j(u)\,\dd u$,
$\mathcal D(f,j,\delta)=\mathsf H(\log(f/\pi)+\psi(\tfrac14+\tfrac j2))+J_j+2\delta\mathsf B$ and
\[
 \mathsf M=\frac1d\Bigl(\mathcal D(1,0,1)+\frac{q-3}2\mathcal D(q,0,0)+\frac{q-1}2\mathcal D(q,1,0)\Bigr),
\]
no prime or zero enters $\mathsf M$, because $h<\log2$.

\begin{theorem}[Sharp multiscale contrast; \stG]\label{zer:thm:sharpwin}
$\mathsf E$ extends to a continuous real positive-definite function $\mathcal K$ on $G\times\R$,
\[
 \mathcal K(a,t)=\frac1d\sum_{\chi\bmod q}\overline{\chi(a)}\sum_\gamma m_\chi(\gamma)w_h(\gamma)e^{i\gamma t},\qquad
 w_h(y)=\frac{e^h+1-2e^{h/2}\cos(hy)}{1/4+y^2}>0,
\]
with $\mathcal K(1,0)=\mathsf M>0$ and $\mathcal K(a^{-1},-t)=\mathcal K(a,t)$. Hence
\[
 \abs{\mathsf E(a,t)}\le\mathsf M\quad(t>h),\qquad\limsup_{t\to\infty}\mathsf E(1,t)=\mathsf M .
\]
For a finite family $(a_i,t_i)$ and $c_i\in\C$ let
$Q=\sum_{i,j}c_i\overline{c_j}\mathcal K(a_ia_j^{-1},t_i-t_j)\ge0$. For
$T>h+\max\abs{t_i-t_j}$,
\[
 \Bigl|\sum_{i,j}c_i\overline{c_j}\,\mathsf E(a_ia_j^{-1},T+t_i-t_j)\Bigr|\le Q,
\]
and the limsup of the left side as $T\to\infty$ is $Q$. If all differences $\abs{t_i-t_j}$, $i\ne j$,
exceed $h$, then $Q$ is computed from the displayed prime sums, integrals and $\mathsf M$.
\end{theorem}
\begin{proof}
With $k(v)=e^{v/2}\mathbf 1_{[0,h]}(v)$ one has $g_h(u)=\int k(v)k(v+u)\,\dd v$, so
$\widehat g_h(y)=\int g_h(u)e^{iyu}\dd u=\abs{(e^{(1/2+iy)h}-1)/(1/2+iy)}^2=w_h(y)$, and
$\int g_h(u)e^{\pm u/2}\dd u=\mathsf B$. For a primitive character of conductor $f$, parity $j$ and
$\delta=1$ for conductor one, and for a real compactly supported Lipschitz $F$ that is piecewise
smooth, the Weil explicit formula reads
\begin{align*}
 \sum_\gamma m(\gamma)\widehat F(\gamma)={}&\delta\int F(u)(e^{u/2}+e^{-u/2})\dd u
 +F(0)\bigl(\log(f/\pi)+\psi(\tfrac14+\tfrac j2)\bigr)\\
 &+\int_0^\infty(2F(0)-F(u)-F(-u))b_j(u)\dd u\\
 &-\sum_{n\ge2}\frac{\Lambda(n)}{\sqrt n}\bigl(\chi(n)F(\log n)+\overline{\chi(n)}F(-\log n)\bigr),
\end{align*}
obtained from the residue theorem for the completed logarithmic derivative, integrated
against the transform $\int F(u)e^{(s-1/2)u}\dd u$, with the digamma integral representation for the Gamma
terms \cite{Koukoulopoulos2019Distribution}; smooth compact mollification and the bound
$\widehat F(y)=O((1+y^2)^{-1})$ justify it for these $F$. Take $F=g_h(\cdot-t)$ with $t>h$, so
$F(0)=F(-\log n)=0$, and average with $\overline{\chi(a)}/d$: the pole terms give $C(t)$;
orthogonality gives the class $a$ for $(n,q)=1$, and $n=q^k$ only through $\zeta$, with weight
$1/d$, which is $P(a,t)$; and writing
\[
 b_j=\tfrac12(b_0+b_1)+\tfrac12\chi(-1)(b_0-b_1)
\]
gives $A(a,t)$.
So $\mathcal K=\mathsf E$ for $t>h$. At $(1,0)$ the prime terms vanish since $h<\log2$, each character
gives $\mathcal D(f,j,\delta)$, and the $d/2$ even characters include the principal one. The
remaining assertions are \cref{zer:lem:posspec}(iii), the weights $m_\chi(\gamma)w_h(\gamma)/d$ being
positive.
\end{proof}

\begin{corollary}[Two and three windows; \stG]\label{zer:cor:sharpwin}
For $u>h$, $a\in G$ and $T>u+h$,
\[
 \abs{2\mathsf E(1,T)-\mathsf E(a,T+u)-\mathsf E(a^{-1},T-u)}\le2\bigl(\mathsf M-\mathsf E(a,u)\bigr),
\]
sharp in the limsup, which improves the triangle bound $4\mathsf M$ whenever $\mathsf E(a,u)>-\mathsf M$.
For $g_1,g_2\in G\times\R$ with $X=\mathcal K(g_1)$, $Y=\mathcal K(g_2)$, $Z=\mathcal K(g_1g_2^{-1})$,
\[
 \mathsf M^3+2XYZ-\mathsf M(X^2+Y^2+Z^2)\ge0,
\]
equivalently, if $\abs Z<\mathsf M$, $\dfrac{(X+Y)^2}{2\mathsf M(\mathsf M+Z)}+\dfrac{(X-Y)^2}{2\mathsf M(\mathsf M-Z)}\le1$.
\end{corollary}
\begin{proof}
Take $(1,0),(a,u)$ with coefficients $(1,-1)$, so $Q=2(\mathsf M-\mathcal K(a,u))$; and use the
nonnegative determinant of the Gram matrix of $(1,0),g_1,g_2$.
\end{proof}

For $q=5$, $h=0.4$, $a=2$, $u=\log10$ and $T=\log1000$, the prime-sum evaluation gives
$\mathsf M\approx0.44998$, a translated contrast $\approx-0.21147$ and the sharp bound $\approx0.63564$,
about $35\%$ of the triangle bound $4\mathsf M$; these are floating-point values from
\begin{center}\artifact{sharp_multiscale_prime_windows/check_windows.py}.\end{center}

\paragraph{Arithmetic tangents of the zero spectrum.}
For a primitive real character $\chi$ of conductor $q$ (conductor one gives $\zeta$) the kernel
$K_\chi$ is real, even, and positive semidefinite:
\[
 \sum_{i,j}c_i\overline{c_j}K_\chi(t_i-t_j)=\sum_\gamma m_\chi(\gamma)\frac{\abs{\sum_ic_ie^{i\gamma t_i}}^2}{1/4+\gamma^2}.
\]
It is
the covariance of a centred stationary Gaussian process $X_\chi$: take, for each distinct
positive ordinate, the coordinates $\sqrt{2m_\chi(\gamma)/(1/4+\gamma^2)}(\cos\gamma t,\sin\gamma t)$, and
$2\sqrt{m_\chi(0)}$ for a central zero, and multiply by independent standard normal variables.
These coordinates also give a deterministic Hilbert-space realization with
$\langle v_\chi(t),v_\chi(u)\rangle=K_\chi(t-u)$. By \cref{zer:lem:ramp}, $K_\chi(\log x)$ is the
normalized prime error $\sqrt x(S_\chi(x)-\delta_\chi\log x-c_\chi)$ up to $O((1+\log x)/\sqrt x)$. Put
\[
 d_q=1-\zergE-\log(2\pi)+\log q,\qquad J_\chi(a)=\begin{cases}\chi(n)\Lambda(n)/\sqrt n,&e^a=n\in\N,\ n\ge2,\\0,&\text{otherwise},\end{cases}
\]
and $B(u,v)=-\tfrac12(u\log u+v\log v-\abs{u-v}\log\abs{u-v})$ with $0\log0=0$.

\begin{lemma}[Cusps of the kernel; \stG]\label{zer:lem:cusp}
As $t\downarrow0$, $K_\chi(t)=M_\chi+\tfrac t2(\log t-d_q)+O_\chi(t^2)$. Near each $a>0$,
$K_\chi(t)=G_a(t)+J_\chi(a)(t-a)_+$ with $G_a$ of class $C^2$. For $q>1$ and parity $b$, for $t\ge0$,
\begin{align*}
 K_\chi(t)={}&K_\zeta(t)+(M_\chi-M_\zeta)\cosh(t/2)-\log q\,\sinh(t/2)
 +(2-b)t\sinh(t/2)\\
 &+2b\cosh(t/2)\log\cosh(t/2)-\sum_{n\le e^t}\frac{2\Lambda(n)(1-\chi(n))}{\sqrt n}\sinh\Bigl(\frac{t-\log n}2\Bigr).
\end{align*}
\end{lemma}
\begin{proof}
For $0<t<\log2$ the prime sum in \eqref{zer:eq:zetaexact} is empty, and expanding
$1-e^{-t}=t(1-t/2+O(t^2))$ in $\mathcal T$ gives the zeta case. The prime power $n$ enters
\eqref{zer:eq:zetaexact} through $2\Lambda(n)n^{-1/2}\sinh((t-\log n)/2)\mathbf 1_{t\ge\log n}$, whose linear
part is $J_\zeta(\log n)(t-\log n)_+$ and whose remainder is $C^2$. For $\chi$, let
$F_\chi(s)=(q/\pi)^{(s+b)/2}\Gamma((s+b)/2)L(s,\chi)$ and $H_\chi(s)=F_\chi'/F_\chi(\tfrac12+s)=\sum_\gamma s/(s^2+\gamma^2)$
for $s>\tfrac12$, pairing $\pm\gamma$ (a central zero contributes $m/s$); the same holds for $\xi$. Pairing
gives $\int_0^\infty e^{-st}K_\chi(t)\dd t=(sM_\chi-H_\chi(s))/(s^2-\tfrac14)$. The Dirichlet series and
\[
 \tfrac12\bigl[\psi(\tfrac{z+1}2)-\psi(\tfrac z2)\bigr]=\int_0^\infty\frac{e^{-zu}}{1+e^{-u}}\dd u,\qquad z=\tfrac12+s,
\]
give $H_\chi-H_\zeta=\tfrac12\log q+\int_0^\infty e^{-st}\bigl(b/(2\cosh(t/2))-2\cosh(t/2)\bigr)\dd t
+\sum_n\Lambda(n)(1-\chi(n))n^{-1/2-s}$, where $2\cosh(t/2)$ inverts the terms $1/z+1/(z-1)$ of
$\xi'/\xi$. Inverting with $(s^2-\tfrac14)^{-1}\leftrightarrow2\sinh(t/2)$ gives the formula, and the
convolution integral evaluates to the two elementary terms. For $t<\log2$ the prime sum is empty
and the local expansion follows from the zeta case; at $\log n$ the two prime contributions
combine to $2\chi(n)\Lambda(n)n^{-1/2}\sinh((t-\log n)/2)\mathbf 1_{t\ge\log n}$.
\end{proof}

Consequently $V_\chi(t)=2(K_\chi(0)-K_\chi(t))=t(\log(1/t)+d_q)+O_\chi(t^2)$, and $X_\chi$ has a
continuous version.

\begin{theorem}[Brownian tangents with an arithmetic correction; \stG]\label{zer:thm:tangent}
Fix $U>0$ and distinct reals $a_1,\dots,a_m$, and let $A_\chi(\boldsymbol a)$ be the symmetric matrix
with zero diagonal and entries $J_\chi(\abs{a_i-a_j})$. As $h\downarrow0$, uniformly for $u,v\in[0,U]$,
\begin{multline*}
 \frac1h\operatorname{Cov}\bigl(X_\chi(a_i+hu)-X_\chi(a_i),X_\chi(a_j+hv)-X_\chi(a_j)\bigr)\\
 =\delta_{ij}\bigl[(\log(1/h)+d_q)\min(u,v)+B(u,v)\bigr]-A_\chi(\boldsymbol a)_{ij}\min(u,v)+O(h).
\end{multline*}
The processes $Y_{i,h}(u)=(X_\chi(a_i+hu)-X_\chi(a_i))/\sqrt{h\log(1/h)}$ converge jointly in
$C([0,U])^m$ to independent standard Brownian motions. Moreover, for $a=\log n$ and $\abs v\le M$,
\[
 \frac{2K_\chi(a+vh)-K_\chi(a+(v+1)h)-K_\chi(a+(v-1)h)}h=-\frac{\chi(n)\Lambda(n)}{\sqrt n}(1-\abs v)_++O(h),
\]
and the numerator is $O(h^2)$ at a fixed positive lag outside the logarithms of prime powers
with $\chi(n)\ne0$.
\end{theorem}
\begin{proof}
By stationarity the covariance is
$K(\Delta+h(u-v))-K(\Delta+hu)-K(\Delta-hv)+K(\Delta)$ with $\Delta=a_i-a_j$. For $i=j$ it is
$\tfrac12(V(hu)+V(hv)-V(h\abs{u-v}))$, and the local expansion with $u+v-\abs{u-v}=2\min(u,v)$
gives the diagonal, with error $O(h^2)$. For $i\ne j$ use \cref{zer:lem:cusp} near $\Delta$ (the
even reflection has the same jump at negative lags): the $C^2$ part contributes $O(h^2uv)$, and the
ramp contributes $A_{ij}h((u-v)_+-u)=-A_{ij}h\min(u,v)$. After division by $h\log(1/h)$ the
covariance tends to $\delta_{ij}\min(u,v)$, which gives finite-dimensional convergence of these
centred Gaussian processes, and $\mathbb E\abs{Y_{i,h}(u)-Y_{i,h}(v)}^4=3(V(h\abs{u-v})/(h\log(1/h)))^2
\ll\abs{u-v}^{3/2}$ gives tightness. The triangular profile follows from
$2v_+-(v+1)_+-(v-1)_+=-(1-\abs v)_+$.
\end{proof}

The Gaussian process is built from the zero spectrum; nothing is asserted about the
distribution of prime errors themselves. Take now $a_i=\log i$, $1\le i\le N$, and write
$A_{\chi,N}$ for the matrix: for $i<j$ its entry is $\chi(p)^k\log p/p^{k/2}$ if $j=ip^k$, and $0$
otherwise. Let $V_\chi(h)=2(K_\chi(0)-K_\chi(h))$,
$R_{\chi,N}(h)_{ij}=\bigl(2K_\chi(\log(i/j))-K_\chi(\log(i/j)+h)-K_\chi(\log(i/j)-h)\bigr)/V_\chi(h)$, the
correlation matrix of $N$ increments of length $h$, and $\ell_q(h)=\log(1/h)+d_q$. Let $A_N^{(q)}$ be
the zeta matrix $A_{\zeta,N}$ with every edge at a power of a prime dividing $q$ removed.

\begin{theorem}[Determinants to every logarithmic order; \stG]\label{zer:thm:det}
For fixed $\chi$ and $N$, as $h\downarrow0$, $\ell_q(h)(I-R_{\chi,N}(h))=A_{\chi,N}+O(h)$ in operator
norm, $R_{\chi,N}(h)$ is positive definite for small $h$, and for $\ell_q(h)>\norm{A_{\chi,N}}$
\[
 \log\det R_{\chi,N}(h)=-\sum_{r\ge2}\frac{\tr(A_{\chi,N}^r)}{r\,\ell_q(h)^r}+O\Bigl(\frac h{\ell_q(h)}\Bigr),
\]
with a convergent series. There is a diagonal matrix $D_\chi$ with entries $\pm1$ such that
$A_{\chi,N}=D_\chi A^{(q)}_ND_\chi$. Consequently every coefficient of the expansion depends only on
the set of primes dividing $q$, and if $\chi,\psi$ have conductors $q,r$ with the same prime
divisors,
\[
 \log\det R_{\chi,N}(q\eps)-\log\det R_{\psi,N}(r\eps)=O\Bigl(\frac\eps{\log(1/\eps)}\Bigr)\qquad(\eps\downarrow0).
\]
Explicitly, with
\[
 S_2(N;q)=\sum_{\substack{p^k\le N\\p\nmid q}}\Bigl\lfloor\frac N{p^k}\Bigr\rfloor\frac{(\log p)^2}{p^k},\qquad
 S_3(N;q)=\sum_{\substack{p^k\le N\\p\nmid q}}(k-1)\Bigl\lfloor\frac N{p^k}\Bigr\rfloor\frac{(\log p)^3}{p^k},
\]
\[
 \log\det R_{\chi,N}(h)=-\frac{S_2(N;q)}{\ell_q(h)^2}-\frac{2S_3(N;q)}{\ell_q(h)^3}
 +O\bigl(\ell_q(h)^{-4}+h/\ell_q(h)\bigr).
\]
\end{theorem}
\begin{proof}
With $u=v=1$, $B(1,1)=0$ and \cref{zer:thm:tangent}, the covariance matrix of the increments is
$h(\ell I-A+E_h)$ with $\norm{E_h}=O(h)$, and $V_\chi(h)=h(\ell+O(h))$; hence
$R=I-A/\ell+F_h$ with $\norm{F_h}=O(h/\ell)$. For large $\ell$ the smallest eigenvalues of $I-A/\ell$
and $R$ exceed $\tfrac12$ and $\tfrac14$; the derivative of $\log\det$ along the segment gives
$\abs{\log\det R-\log\det(I-A/\ell)}\le4N\norm{F_h}$, and
$\log\det(I-A/\ell)=-\sum_{r\ge1}\tr(A^r)/(r\ell^r)$ with vanishing $r=1$ term. For the gauge write
$n=n_{\mathrm{ram}}n_{\mathrm{unr}}$ with the primes of $n_{\mathrm{ram}}$ dividing $q$ and
$(n_{\mathrm{unr}},q)=1$, and let $D_\chi=\operatorname{diag}(\chi(n_{\mathrm{unr}}))_{n\le N}$; if $j=ip^k$
with $p\nmid q$ then $\chi(i_{\mathrm{unr}})\chi(j_{\mathrm{unr}})=\chi(p)^k$, and edges at $p\mid q$ vanish
in both matrices. Similar matrices have equal traces of powers, and
$\ell_q(q\eps)=\ell_r(r\eps)=\log(1/\eps)+1-\zergE-\log(2\pi)$. Finally each edge $\set{i,ip^k}$
contributes twice to $\tr(A^2)$, so $\tr(A^2)=2S_2$; in a triangle $i<j<l$ the ratios $j/i=p^a$,
$l/j=r^b$ and $l/i$ are prime powers only if $p=r$, so the triangles are $i,ip^a,ip^{a+b}$, each
traversed six times, and $\tr(A^3)=6S_3$.
\end{proof}

For $\zeta$ and $N=8$ one finds $\tr(A^2)/2=3.1248179\ldots$ and
$\tr(A^3)/3=0.4995370\ldots$. The convergence in
$1/\log(1/h)$ is slow, as the numerical record
\begin{center}\artifact{grh_prime_resonance_2026-11-26/results/verification.json}\end{center}
shows. At higher orders
$\tr(A^r)$ is the weighted number of closed walks of length $r$ in the graph whose edges join $i$
and $ip^k$.

\begin{theorem}[Coalescing Gram determinants; \stG]\label{zer:thm:gramcoal}
Fix $x_0<x_1<\dots<x_n$, let $\Delta_j=x_j-x_{j-1}$ and $L=\log(1/h)$. As $h\downarrow0$,
\[
 \det\bigl(K_\chi(h(x_i-x_j))\bigr)_{0\le i,j\le n}
 =M_\chi(hL)^n\prod_{j=1}^n\Delta_j\Bigl[1+\frac{nd_q-\sum_j\log\Delta_j}L+O(L^{-2})\Bigr].
\]
\end{theorem}
\begin{proof}
Replace $v_\chi(hx_j)$, $j\ge1$, by the consecutive differences, a change of determinant one. The
top-left entry is $M_\chi$; by \cref{zer:lem:cusp} the lower block is
$h(L\operatorname{diag}(\Delta_j)+Q+O(h))$ with $Q_{jj}=\Delta_j(d_q-\log\Delta_j)$, the leading
off-diagonal terms cancelling because disjoint consecutive intervals are orthogonal for
$\min(u,v)$; the top-right entries are $O(hL)$. The Schur complement changes the block by
$O(h^2L^2)$, and $\det(D+Q/L+O(hL))=\det D\,[1+\tr(D^{-1}Q)/L+O(L^{-2})]$.
\end{proof}

This kernel is the prime error itself; the autocorrelation of the prime error at merging scales
behaves differently and is treated next.

\subsection{Merging scales: a cubic logarithmic law}\label{zer:sec:coalesce}
\sources{prime_scale_coalescence_2026-11-26}

\Cref{zer:thm:coercive} keeps the scales fixed. For the zeta function the rate at which the
covariance degenerates as the scales merge is determined exactly up to a bounded factor. Let
\[
 E(t)=e^{t/2}\bigl(S(e^t)-t+\zergE+1\bigr),\qquad S(x)=\sum_{n\le x}\Lambda(n)\Bigl(\frac1n-\frac1x\Bigr),
\]
let $\gamma>0$ run over the distinct positive ordinates of zeta zeros, $m_\gamma$ the
multiplicities, $N(T)=\sum_{0<\gamma\le T}m_\gamma$ and $M_2(T)=\sum_{0<\gamma\le T}m_\gamma^2$. Since
$\zeta(\tfrac12)\ne0$ (the alternating eta series is positive at $\tfrac12$), \cref{zer:lem:ramp}
gives $E(t)=2\sum_{\gamma>0}m_\gamma\cos(\gamma t)/(\tfrac14+\gamma^2)+O((1+t)e^{-t/2})$, and by
\cref{zer:lem:ap}(i) the long-time covariance of $E(t+h_i)$ and $E(t+h_j)$ is $C(h_i-h_j)$, where
\[
 C(h)=2\sum_{\gamma>0}\frac{m_\gamma^2\cos(\gamma h)}{(1/4+\gamma^2)^2},\qquad C_0=C(0),\qquad
 D_2=2\sum_{\gamma>0}\frac{m_\gamma^2\gamma^2}{(1/4+\gamma^2)^2}.
\]
The multiplicity is squared because coincident frequencies are combined first. Fix $r\ge3$
distinct reals $u_1,\dots,u_r$, put $G_\eps=(C(\eps(u_i-u_j)))_{i,j}$, $\tau_\eps=\eps^3\log(1/\eps)$,
$\bar u$ the mean of the $u_i$, $v=u-\bar u\mathbf 1$, and
$W=\set{c\in\R^r:c\cdot\mathbf 1=c\cdot u=0}$. For $c\in W$ let
\[
 H_c(t)=\sum_ic_i(t-u_i)_+,\qquad \mathcal B(c)=\int_\R H_c(t)^2\dd t=\frac1{12}\sum_{i,j}c_ic_j\abs{u_i-u_j}^3,
\]
and let $b_1\ge\dots\ge b_{r-2}>0$ be the eigenvalues of $\mathcal B$ on $W$.

\begin{theorem}[Arithmetic coalescence law; \stG]\label{zer:thm:coalesce}
The long-time limit is taken first, then $\eps\downarrow0$.
\begin{enumerate}[label=\textup{(\roman*)}]
\item Uniformly for $c$ in the unit sphere of $W$,
$\mathcal B(c)-o(1)\le c^{\mathsf T}G_\eps c/\tau_\eps\le\tfrac43\mathcal B(c)+o(1)$.
\item $\lambda_1(G_\eps)\sim rC_0$, $\lambda_2(G_\eps)\sim D_2\norm v^2\eps^2$, and for $1\le j\le r-2$
\[
 b_j\le\liminf_{\eps\downarrow0}\frac{\lambda_{j+2}(G_\eps)}{\tau_\eps}\le\limsup_{\eps\downarrow0}\frac{\lambda_{j+2}(G_\eps)}{\tau_\eps}\le\frac43b_j .
\]
\item $\det G_\eps\asymp_u\eps^{3r-4}\log(1/\eps)^{r-2}$ and the condition number is
$\asymp_u\eps^{-3}/\log(1/\eps)$.
\end{enumerate}
All conclusions persist after deleting any fixed finite set of positive ordinates with their
negatives; $C_0$ and $D_2$ change, while $\mathcal B$ and the factor $\tfrac43$ do not.
\end{theorem}

With two offsets only the first two statements of (ii) hold. No simplicity, independence or
pair-correlation conjecture is used; the factor $\tfrac43$ is an explicit comparison, not a
claim of optimality. The arithmetic input is Montgomery's theorem, proved under RH
\cite{Montgomery1973Pair}, which gives the classical bound $M_2(T)\le(\tfrac43+o(1))N(T)$
\cite{ConreyGhoshGonek1998Simple}. A cumulative bound does not compare nonmonotone test
functions; the following local form does.

\begin{lemma}[Local squared multiplicities; \stG]\label{zer:lem:localM2}
For fixed $0<a<b$,
\[
 b-a\le\liminf_{T\to\infty}\frac{M_2(bT)-M_2(aT)}{T\log T/(2\pi)}\le
 \limsup_{T\to\infty}\frac{M_2(bT)-M_2(aT)}{T\log T/(2\pi)}\le\frac43(b-a),
\]
and $M_2(T)=O(T\log(2+T))$.
\end{lemma}
\begin{proof}
Let $w(y)=4/(4+y^2)$. Montgomery's equation (6) gives, for $0<\alpha<1$ and uniformly for $\alpha$ in
compact subintervals,
\[
 \sum_{0<\gamma,\gamma'\le X}\zersinc^2\Bigl(\frac\alpha2(\gamma-\gamma')\log X\Bigr)w(\gamma-\gamma')
 =\Bigl(\frac1\alpha+\frac\alpha3+o(1)\Bigr)\frac{X\log X}{2\pi},
\]
the double sum repeating zeros with multiplicity: the left side divided by $X\log X/(2\pi)$ is
$\alpha^{-2}\int_{-\alpha}^\alpha(\alpha-\abs\beta)F_X(\beta)\dd\beta$, and
$F_X(\beta)=(1+o(1))X^{-2\abs\beta}\log X+\abs\beta+o(1)$ uniformly on $\abs\beta\le1-\delta$. Fix $\alpha$ and use
the single kernel $k_T(y)=\zersinc^2(\tfrac\alpha2y\log T)w(y)\ge0$ for all $X=cT$, with
$\alpha_X=\alpha\log T/\log(cT)\to\alpha$; uniformity gives
$S_T(cT)=(1/\alpha+\alpha/3)cT\log T/(2\pi)+o(T\log T)$. Raising $X$ from $aT$ to $bT$ adds the
diagonal terms $m_\gamma^2$ of the new ordinates and nonnegative off-diagonal terms, so
$M_2(bT)-M_2(aT)\le S_T(bT)-S_T(aT)$. Subtracting the two asymptotic equalities and letting
$\alpha\uparrow1$ after the height limit gives the upper bound. The lower bound follows from
$m_\gamma^2\ge m_\gamma$ and the Riemann--von Mangoldt formula, and the global bound from one fixed
$\alpha$.
\end{proof}

Subtracting two upper bounds for $M_2$ would be invalid; the proof subtracts asymptotic
equalities for one kernel, and stays in the proved range of Montgomery's theorem.

\begin{proof}[Proof of \cref{zer:thm:coalesce}]
\emph{Spline identity.} For $c\in W$, $H_c$ vanishes outside $[\min u_i,\max u_i]$ by the two
moment conditions. With $J_c(t)=\sum_ic_i\abs{t-u_i}^3$ one has $J_c''=12H_c$ and
$H_c''=\sum_ic_i\delta_{u_i}$, so $12\int H_c^2=\sum_ic_iJ_c(u_i)$, which is the displayed cubic form;
it is positive for $c\ne0$, since $H_c=0$ would make all slope jumps vanish. By Parseval,
$\mathcal B(c)=\pi^{-1}\int_0^\infty\abs{p_c(x)}^2x^{-4}\dd x$ with $p_c(x)=\sum_ic_ie^{iu_ix}$.

\emph{Scaled measures.} Let $\nu_T=\frac{2\pi}{T\log T}\sum_{\gamma>0}m_\gamma^2\delta_{\gamma/T}$. By
\cref{zer:lem:localM2} and approximation by step functions,
$\int f\le\liminf\int f\dd\nu_T\le\limsup\int f\dd\nu_T\le\tfrac43\int f$ for continuous nonnegative
compactly supported $f$ on $(0,\infty)$. Since $\abs{p_c(x)}\le A_u\norm c\min(x^2,1)$ by the two
cancellations, and $\limsup\nu_T((0,\delta])=O(\delta)$,
$\limsup\int_R^\infty x^{-4}\dd\nu_T=O((1+\log R)/R^3)$, the comparison extends to
$f_c=\abs{p_c}^2/x^4$, also with $x^4$ replaced by $(x^2+1/(4T^2))^2$, uniformly over unit $c\in W$
by a finite net. Since
\[
 \frac{c^{\mathsf T}G_\eps c}{\tau_\eps}=\frac1\pi\int_0^\infty\frac{\abs{p_c(x)}^2}{(x^2+1/(4T^2))^2}\dd\nu_T(x),\qquad T=1/\eps,
\]
this proves (i).

\emph{Separating the scales.} With $K(x)=\cos x-1+x^2/2\in[0,\min(x^4/24,x^2/2)]$ and
$M_2(T)=O(T\log T)$, $\sum_{\gamma>T}m_\gamma^2/\gamma^2=O(\log T/T)$, so
$C(h)=C_0-\tfrac12D_2h^2+O(\abs h^3\log(1/\abs h))$. Centre the nodes, so $v=u$. Then
$G_\eps=C_0\mathbf 1\mathbf 1^{\mathsf T}+\eps^2A+O(\tau_\eps)$ with
$A=-\tfrac12D_2(u^{\circ2}\mathbf 1^{\mathsf T}+\mathbf 1(u^{\circ2})^{\mathsf T}-2uu^{\mathsf T})$, whose
compression to $\mathbf 1^\perp$ is $D_2uu^{\mathsf T}$ and to $W$ is $0$. In an orthonormal basis
$\mathbf 1/\sqrt r$, $u/\norm u$, $W$, the matrix has blocks $rC_0+O(\eps^2)$,
$D_2\norm u^2\eps^2+O(\tau_\eps)$ and $\tau_\eps Q_\eps$ with $\mathcal B-o(1)\preceq Q_\eps\preceq\tfrac43\mathcal B+o(1)$,
off-diagonal blocks $O(\eps^2)$ against the first and $O(\tau_\eps)$ between the last two (this
uses both cancellations). Eliminating the first block changes the $W$-block by $O(\eps^4)$, and the
second by $O(\tau_\eps^2/\eps^2)=o(\tau_\eps)$, through congruences tending to the identity. By
min--max, congruence by matrices with singular values tending to $1$ changes positive
eigenvalues by factors $1+o(1)$, which gives (ii), and (iii) follows by multiplying and dividing
eigenvalues. Deleting finitely many frequencies changes $N$ and $M_2$ by bounded amounts, and
each deleted frequency contributes $O(\eps^4)$ to $c^{\mathsf T}G_\eps c$ for $c\in W$.
\end{proof}

For fixed $c\ne0$ the mean square of $\sum_ic_iE(t+\eps u_i)$ thus has exactly three possible
orders: $C_0(\sum c_i)^2$ if $\sum c_i\ne0$; $\sim D_2(\sum u_ic_i)^2\eps^2$ if $\sum c_i=0\ne\sum u_ic_i$;
and between $\mathcal B(c)\tau_\eps$ and $\tfrac43\mathcal B(c)\tau_\eps$ asymptotically if $c\in W$. Further
moment cancellations do not lower the order, since $\mathcal B>0$ on $W$.

\begin{corollary}[Raw finite differences; \stG]\label{zer:cor:rawdiff}
For a fixed integer $d\ge2$, the long-time mean square $V_d(\eps)$ of
$\sum_{j=0}^d(-1)^{d-j}\binom djE(t+j\eps)$ satisfies
\[
 \mathcal B_d\le\liminf\frac{V_d(\eps)}{\tau_\eps}\le\limsup\frac{V_d(\eps)}{\tau_\eps}\le\frac43\mathcal B_d,\qquad
 \mathcal B_d=\frac1{12}\sum_{i,j=0}^d(-1)^{i+j}\binom di\binom dj\abs{i-j}^3>0,
\]
with $\mathcal B_2,\mathcal B_3,\mathcal B_4,\mathcal B_5=\tfrac23,1,\tfrac83,\tfrac{25}3$.
\end{corollary}
\begin{proof}
The difference coefficients annihilate constants and linear functions, so they lie in $W$.
\end{proof}

These are raw differences; a finite truncation of the zero series would give $O(\eps^{2d})$, so
no fixed table of zeros can confirm the cubic law.

\begin{theorem}[Saturation of linear prediction; \stG]\label{zer:thm:predict}
For a target node $k$ let $V_k^*(\eps)$ be the infimum over real $a_i$ $(i\ne k)$, which may depend
on $\eps$, of the long-time mean square of $E(t+\eps u_k)-\sum_{i\ne k}a_iE(t+\eps u_i)$, and
$B_k^*=\min\set{\mathcal B(c):c\in W,\ c_k=1}>0$. Then
\[
 B_k^*\le\liminf\frac{V_k^*(\eps)}{\tau_\eps}\le\limsup\frac{V_k^*(\eps)}{\tau_\eps}\le\frac43B_k^*,
\]
and an intercept does not help, since all observations have mean zero.
\end{theorem}
\begin{proof}
The projection $w_k$ of $e_k$ to $W$ is nonzero, since otherwise a nonzero affine function would
vanish at the other $r-1\ge2$ nodes. Then $V_k^*=1/(e_k^{\mathsf T}G_\eps^{-1}e_k)$ and
$B_k^*=1/(w_k^{\mathsf T}\mathcal B^{-1}w_k)$. The block elimination gives
$\tau_\eps G_\eps^{-1}=P_WQ_\eps^{-1}P_W+o(1)$, since $\tau_\eps/\eps^2\to0$; invert the bounds on $Q_\eps$.
\end{proof}

For $u=(-1,0,1)$, $W$ is spanned by $(1,-2,1)$, $\mathcal B(1,-2,1)=\tfrac23$, $b_1=\tfrac19$ and
$B^*_{\mathrm{middle}}=\tfrac16$. So $\lambda_3(G_\eps)/\tau_\eps$ lies asymptotically in
$[\tfrac19,\tfrac4{27}]$, and the optimal prediction error of the middle value in
$[\tfrac16,\tfrac29]$ after division by $\tau_\eps$; the midpoint predictor $\tfrac12(E(t-\eps)+E(t+\eps))$
is optimal up to a factor $1+o(1)$, its excess error being
$[C_0+C(2\eps)-2C(\eps)]^2/(2[C_0+C(2\eps)])=O(\eps^4)$. These finite constants are checked in exact
rational arithmetic in
\begin{center}\artifact{prime_scale_coalescence_2026-11-26/verification/verify_geometry.py}.\end{center}

\begin{proposition}[Exact constants; \stG, with the asymptotic for $M_2$ as a hypothesis]\label{zer:prop:kappa}
If moreover $M_2(T)\sim\kappa T\log T/(2\pi)$, then $1\le\kappa\le\tfrac43$,
$G_\eps|_W/\tau_\eps\to\kappa\mathcal B$,
$\lambda_{j+2}(G_\eps)/\tau_\eps\to\kappa b_j$, $V_k^*(\eps)/\tau_\eps\to\kappa B_k^*$, the optimal predictor
converges to the minimizer defining $B_k^*$, and
$\det G_\eps/(\eps^{3r-4}\log(1/\eps)^{r-2})\to rC_0D_2\norm v^2\kappa^{r-2}\det\mathcal B$. Also
$C(h)=C_0-\tfrac12D_2h^2+\tfrac\kappa{12}\abs h^3\log(1/\abs h)+o(\abs h^3\log(1/\abs h))$.
\end{proposition}
\begin{proof}
Now $\nu_T$ converges vaguely to $\kappa\,\dd x$, and the tail bounds upgrade this to the quadratic
forms; the block elimination and inversion give the limits, the optimal contrast being
\[
 \frac{G_\eps^{-1}e_k}{e_k^{\mathsf T}G_\eps^{-1}e_k}\to\frac{\mathcal B^{-1}w_k}{w_k^{\mathsf T}\mathcal B^{-1}w_k}.
\]
For $C(h)$ scale the remainder $K$, using
$\int_0^\infty K(x)x^{-4}\dd x=\tfrac16\int_0^\infty(1-\cos x)x^{-2}\dd x=\pi/12$.
\end{proof}

If all but finitely many zeros are simple then $\kappa=1$; that almost all zeros are simple would
not suffice, because the covariance involves squared multiplicities. The mean-square derivative
of the limiting signal exists, its squared coefficients summing to $D_2$, but no second
derivative does, since $\sum m_\gamma^2\gamma^4/(\tfrac14+\gamma^2)^2=\infty$; this roughness at the scale
$\gamma\asymp1/\eps$, together with the density of zeros, produces the factor $\eps^3\log(1/\eps)$. The
corresponding matrix problem for residue classes, where coincident zeros of different
$L$-functions add structure, is not settled by this argument.

\subsection{A prime race in a fixed window}\label{zer:sec:window}
\sources{grh_transfers; grh-transfer}

Let $b(p)$ be a real weight on primes, $B(y)=\sum_{p\le y}b(p)\log p$, and for a fixed
dilation $\tau>1$ let
\begin{equation}\label{zer:eq:Jh}
 J_\tau(x)=-\int_x^{\tau x}B(y)\,\frac{\dd y}{y^2}
 =-\sum_{p\le\tau x}b(p)\log p\Bigl(\frac1{\max(p,x)}-\frac1{\tau x}\Bigr).
\end{equation}
Every value of $J_\tau$ involves finitely many primes. Two weights are considered:
$b(p)=\chi_4(p)$ for odd $p$, attached to the Dirichlet beta function $\beta(s)$, for which
a positive value of $J_\tau$ favors primes $3\bmod4$; and
$b(p)=\#\set{y\in\F_p:y^3-y-1=0}-1$ for $p\ne23$, attached to the standard two-dimensional
Artin $L$-function $L(s,\rho)$ of the cubic field of discriminant $-23$, for which a
positive value favors irreducible primes over twice the completely split primes. The
factor $1/n-1/x$ is the smoothing that appears in the source proof.

\begin{theorem}[The modulo-four race; \stG]\label{zer:thm:mod4}
For $b=\chi_4$ one has $J_4(x)>0$ for every real $x>3/4$, $J_4(x)=0$ for $0<x\le3/4$, and
$\liminf_{x\to\infty}\sqrt x\,J_4(x)>3/8$.
\end{theorem}

The starting point $3/4$ is sharp. The proof has an asymptotic part common to both weights
and an explicit part special to $\chi_4$.

\paragraph{The square terms have the decisive sign.}
For the cubic, the polynomial is irreducible by the rational root test and has nonsquare
discriminant $-23$, so its splitting field has group $S_3$; the standard representation
$\rho$ has character equal to the number of fixed letters minus one. For unramified $p$ let
$g_p$ denote Frobenius. The traces are as follows.
\begin{center}\small
\begin{tabular}{lrrrr}\toprule
Cycle type & Number in $S_3$ & Roots & $\tr\rho(g_p)$ & $\tr\rho(g_p^2)$\\\midrule
Identity &1&3&2&2\\
Transposition &3&1&0&2\\
Three-cycle &2&0&$-1$&$-1$\\\bottomrule
\end{tabular}\end{center}
Let $A(y)=\sum_{p^k\le y}\tr\rho(g_p^k)\log p$, using $\chi_4(p)^k$ in the quadratic case.
The mean of the last column is $(2+6-2)/6=1$. Chebotarev's theorem and partial summation
give a square contribution $\sqrt y+o(\sqrt y)$ for the fixed cubic field, without any
hypothesis on its zeros. For $\chi_4$ the squares have weight one at every odd prime, so
the ordinary prime number theorem suffices. Powers with $k\ge3$ contribute
$O(y^{1/3}\log^2y)$ in both cases. With $H_\tau(x)=-\int_x^{\tau x}A(y)y^{-2}\,\dd y$ this
gives
\begin{equation}\label{zer:eq:squares}
 \sqrt x\,\bigl[J_\tau(x)-H_\tau(x)\bigr]\longrightarrow2(1-\tau^{-1/2}).
\end{equation}
Passing from primes to prime powers changes the main term at precisely the scale of the
desired bias.

\paragraph{The oscillation of the prime powers.}
Use $L_*$ for the $L$-function with the omitted local factor removed, and
$D_*(s)=-L_*'(s)/L_*(s)$; in the cubic case $L_*=(1-23^{-s})L$ \cite{Serre2003Jordan}.
Write $a_{p^k}=\tr\rho(g_p^k)\log p$, or $a_{p^k}=\chi_4(p)^k\log p$, and set
\[
 S(x)=\sum_{p^k\le x}a_{p^k}\bigl(p^{-k}-x^{-1}\bigr),\qquad H_\tau(x)=S(x)-S(\tau x).
\]
Suppose that RH holds for this individual $L$. Denote its nontrivial zeros, repeated with
multiplicity, by $\rho'=\tfrac12+i\gamma$, and put
\begin{equation}\label{zer:eq:Z}
 Z=\sum_{\rho'}\frac1{1/4+\gamma^2}.
\end{equation}
The smoothed explicit formula is
\begin{equation}\label{zer:eq:smoothexplicit}
 S(x)=D_*(1)+\sum_{\rho'}\frac{x^{\rho'-1}}{\rho'(1-\rho')}+O(x^{-3/4}).
\end{equation}
Indeed the inverse Mellin kernel for the sum is $D_*(s)x^{s-1}/[s(s-1)]$ on $\Re s=2$. Move
the line to $1/4$. The pole at $1$ comes from the kernel, and the residue of $D_*$ at a
zero of multiplicity $m$ is $-m$. The omitted local factor is nonzero and has bounded
logarithmic derivative on this new line. The remaining vertical integral is $O(x^{-3/4})$:
the standard partial fraction formula for the logarithmic derivative of a fixed completed
$L$-function, Stirling's formula and the $O(\log(2+T))$ count of zeros in a unit interval
give $D_*(1/4+it)=O(\log^2(3+\abs t))$. For the horizontal segments take heights separated
by $\gg1/\log T$ from nearby zero ordinates; the same bound there and the squared
denominator let the segments tend to zero. These fixed-function estimates and their Artin
version are described in \cite{Ng2000Limiting}. The series in
\eqref{zer:eq:smoothexplicit} converges absolutely, since the number of ordinates up to $T$
is $O(T\log T)$.

The constant $D_*(1)$ disappears from $H_\tau$. Equations \eqref{zer:eq:squares} and
\eqref{zer:eq:smoothexplicit} yield
\begin{equation}\label{zer:eq:liminf}
 \liminf_{x\to\infty}\sqrt x\,J_\tau(x)\ \ge\ 2(1-\tau^{-1/2})-(1+\tau^{-1/2})Z .
\end{equation}
For the completed primitive $L$-function $\zerLL$, both examples have sign $+1$ and
$\zerLL(1/2)\ne0$. The Hadamard product for the even entire function
$z\mapsto\zerLL(1/2+z)$ pairs zeros at $\pm i\gamma$. Its order is one, and evenness
removes the linear exponential. Differentiating the paired product at $z=1/2$ gives
\begin{equation}\label{zer:eq:Zformula}
 Z=2\,\zerLL'(1)/\zerLL(1).
\end{equation}

\begin{proof}[Proof of \cref{zer:thm:mod4}]
\emph{The asymptotic bound.} Here
$\zerLL(s)=(4/\pi)^{(s+1)/2}\Gamma((s+1)/2)\beta(s)$, and the positive alternating beta
series gives $\beta(s)>0$ for all real $s>0$. Formula \eqref{zer:eq:Zformula} gives
\[
 Z=\log(4/\pi)-\zergE+2\beta'(1)/\beta(1)<5/12 .
\]
Indeed $\beta(1)=\pi/4>3/4$. The alternating upper bound
$\beta'(1)\le\sum_{j=1}^{15}(-1)^{j+1}\log(2j+1)/(2j+1)<1/4$ uses terms decreasing from $3$
onwards; the finite inequality is verified by rational logarithm enclosures. Use also
$\log(4/\pi)<1/4$ and $\zergE>1/2$. \Cref{src:thm:dgrh} supplies RH for $\beta$. Taking
$\tau=4$ in \eqref{zer:eq:liminf} proves the asymptotic lower bound $1-\tfrac32Z>3/8$.

\emph{An explicit estimate for large $x$.} For $x>1$ shift the contour defining $S(x)$
farther to the left. In addition to the constant $D_*(1)$ and the zero series from
\eqref{zer:eq:smoothexplicit}, the exact formula has the terms
\[
 \frac{\beta'(0)}{x\beta(0)}-\sum_{j=0}^\infty\frac{x^{-2j-2}}{(2j+1)(2j+2)},
\]
the residues at $s=0$ and at the negative odd integers, respectively. To justify the full
shift, move to $\Re s=-2M$ along good horizontal heights and then let $M\to\infty$. On that
line the functional equation, the digamma reflection formula, Stirling's bound and the
convergent logarithmic derivative at $1-s$ give $O(\log(2+M+\abs{\Im s}))$ for the
logarithmic derivative, uniformly for integers $M\ge1$. The resulting integral is
$O(x^{-2M-1}\log(2+M)/M)$. The functional equation also gives
\[
 \beta'(0)/\beta(0)=\zergE+\log(\pi/2)-\beta'(1)/\beta(1)>0 .
\]
The trivial-zero coefficients sum to $\log2<1$. Consequently
\[
 H_4(x)=S(x)-S(4x)\ge-\frac{3Z}{2\sqrt x}-x^{-2}.
\]

The higher prime powers supply an explicit gain. For $p\equiv3\bmod4$, exponents $k\ge2$
contribute alternately $+\log p,-\log p,\dots$, so their sum is nonnegative; for
$p\equiv1\bmod4$ each term is positive. Every odd prime in $t^{1/3}<p\le\sqrt t$
contributes $\log p$. Therefore for $t\ge65{,}536$,
\[
 A(t)-B(t)\ \ge\ \vartheta(\sqrt t)-\vartheta(t^{1/3})\ >\ 0.84\sqrt t-1.02\,t^{1/3}
 \ >\ 0.67\sqrt t ,
\]
where $\vartheta(y)=\sum_{p\le y}\log p$. Theorems 9--10 of \cite{RosserSchoenfeld1962Approximate}
give the upper bound $1.01624\,y$ for all positive $y$ and the lower bound $0.84\,y$ for
$y\ge101$; we used $t^{1/6}>6$. Integration yields
\[
 \sqrt x\,J_4(x)>0.67-5/8-x^{-3/2}>0\qquad(x\ge65{,}536).
\]

\emph{The finite range.} On $[1,65{,}536]$, \eqref{zer:eq:Jh} is continuous and affine in
$1/x$ between successive breakpoints $p$ and $p/4$ for odd primes $p$. It suffices to check
these breakpoints and the two endpoints. The certificate
\artifact{grh_transfers/checks/mod4_all_x.py} encloses each $\log p$ using integer arithmetic:
reduce the argument to $1\le a/b\le2$, put $w=(a-b)/(a+b)\le1/3$, and round
outward in
\[
 \log(a/b)=2\sum_{j=0}^{15}\frac{w^{2j+1}}{2j+1}+R,\qquad
 0\le R\le\frac{2(1/3)^{33}}{33(1-1/9)} .
\]
At scale $2^{48}$ the tail costs less than one integer unit. The certificate covers all
$29{,}541$ test points and gives $J_4(x)\ge3681/10^6$ on the whole interval
(\artifact{grh_transfers/checks/mod4_all_x.txt}). Finally, for $3/4<x<1$ the only
contributing prime is $3$, with weight $\log3\,(1/3-1/(4x))>0$, and for $0<x\le3/4$ no
term has positive support.
\end{proof}

\paragraph{Controlling a degree-two zero sum.}
For $K=\Q[\alpha]$, $\alpha^3-\alpha-1=0$, the standard Artin factor is the Hecke
$L$-function of the weight-one form
\[
 f(z)=\eta(z)\eta(23z),\qquad\zeta_K(s)=\zeta(s)L(s,\rho),
\]
with completed function $\zerLL(s)=(\sqrt{23}/(2\pi))^s\Gamma(s)L(s,\rho)$ and
$\zerLL(s)=\zerLL(1-s)$ \cite{Serre2003Jordan,Zagier2008Elliptic}. The eta product is
positive on the positive imaginary axis, and its Mellin integral makes $\zerLL(s)>0$ for
all real $s$, including $s=1/2$. The zero sum $Z$ of \eqref{zer:eq:Z} needs quantitative
control of the logarithmic derivative as well.

\begin{proposition}[\stU]\label{zer:prop:eta}
Let $F(y)=\prod_j\eta(id_jy)^{r_j}$ with finitely many $d_j,r_j>0$. For the entire Mellin
transform $M(s)=\int_0^\infty F(y)y^s\,\dd y/y$ let
\[
 c=\frac\pi{12}\inf_{y>0}\sum_jr_j\max\set{d_jy,(d_jy)^{-1}}.
\]
For every real $\sigma$, $(\log M)''(\sigma)\le1/c$. The transform has no zero at any $s$
with $\abs{\Im s}\le(\pi/2)\sqrt c$.
\end{proposition}
\begin{proof}
Set $v=\log y$. The product formula for $\eta(ie^v)$ starts with $\exp(-\pi e^v/12)$. For
$z=ae^v$, the derivative of $\log(1-e^{-z})$ in $v$ is $z/(e^z-1)$, decreasing in $v$. Thus
$-(\log\eta(ie^v))''\ge(\pi/12)e^v$. Termwise differentiation is justified locally by
exponential convergence. Eta inversion gives $\eta(ie^{-v})=e^{v/2}\eta(ie^v)$, making this
second logarithmic derivative even. Apply the two lower bounds to each factor separately.
The potential $V(v)=-\log F(e^v)$ satisfies $V''\ge c$. The eta product and its inversion
also give double exponential decay in $v$ at both ends, up to an exponential linear factor.
The integral is entire and all differentiations below are justified.

Give $U$ density proportional to $\exp(\sigma v-V(v))$, and write $\mu=\mathbb E\,U$.
Logarithmic differentiation gives $(\log M)''(\sigma)=\operatorname{Var}U$. Integration by
parts gives $\mathbb E[(U-\mu)(V'(U)-\sigma)]=1$. Subtracting $V'(\mu)$ in place of $\sigma$
leaves this expectation unchanged. Strong convexity then gives $1\ge c\operatorname{Var}U$.
Finally, $\cos w\ge\pi/4-w^2/\pi$ on $\R$, strictly except at $w=\pm\pi/2$: on $[0,\pi/2]$
this follows by differentiating and using $\sin w\ge2w/\pi$, and past $\pi/2$ one has
$2w/\pi\ge1$. The centered characteristic function consequently has strictly positive real
part whenever $t^2\operatorname{Var}U\le\pi^2/4$. It equals
$e^{-it\mu}M(\sigma+it)/M(\sigma)$.
\end{proof}

\begin{corollary}\label{zer:cor:cubicZ}
For $f=\eta(z)\eta(23z)$ one has $c=\pi\sqrt{23}/6$, and the completed function
$\zerLL(s)$ has no zero with $\abs{\Im s}\le(\pi/2)\sqrt c=2.48915\ldots$ \textup{(\stU)}.
If all zeros of $L(s,\rho)$ in the critical strip have real part $1/2$, then
$Z\le1/c=6/(\pi\sqrt{23})$ and
\[
 \liminf_{x\to\infty}\sqrt x\,J_4(x)\ \ge\ 1-\frac9{\pi\sqrt{23}}>0.4026
\]
for the cubic weight $b(p)=\#\set{y\in\F_p:y^3-y-1=0}-1$.
\end{corollary}
\begin{proof}
For $d_1=1$, $d_2=23$ the infimum is $2\sqrt{23}$, attained at $y=1/\sqrt{23}$, by
comparison with $1/y+23y$. Multiplying $M(s)$ by $23^{s/2}$ gives $\zerLL(s)$ and leaves
the second logarithmic derivative unchanged. The symmetry forces $(\log\zerLL)'(1/2)=0$.
Integrating the variance bound up to $1$ gives $(\log\zerLL)'(1)\le1/(2c)$, and
\eqref{zer:eq:Zformula} gives $Z\le1/c$ under the Riemann hypothesis for $L(s,\rho)$. Take
$\tau=4$ in \eqref{zer:eq:liminf}.
\end{proof}

\begin{theorem}[The cubic race; \stR]\label{zer:thm:cubicrace}
For the cubic weight $b(p)=\#\set{y\in\F_p:y^3-y-1=0}-1$, $p\ne23$, the following are
equivalent:
\begin{enumerate}[label=\textup{(\roman*)}]
\item every zero of $L(s,\rho)$ in the critical strip has real part $1/2$;
\item $J_4(x)\ge0$ for every sufficiently large real $x$;
\item $\displaystyle\liminf_{x\to\infty}\sqrt x\,J_4(x)\ge1-\frac9{\pi\sqrt{23}}>0.4026$.
\end{enumerate}
The same equivalence holds with strict positivity in \textup{(ii)}. Without any hypothesis,
the cubic $J_4$ takes positive values at arbitrarily large $x$.
\end{theorem}
\begin{proof}
The implication (i)$\Rightarrow$(iii) is \cref{zer:cor:cubicZ}, and (iii)$\Rightarrow$(ii) is
immediate, with strict positivity. For (ii)$\Rightarrow$(i), the Mellin transform of
$J_\tau$ is $(1-\tau^{s-1})/(s(s-1))$ times $\sum_{p\ne23}b(p)(\log p)p^{-s}$ for
$\Re s>1$; it has a pole at every zero of $L(s,\rho)$ with $1/2<\Re s<1$, it is regular on
the real interval $(1/2,\infty)$ because $\zerLL$ is positive there, and the Landau step
given in \cref{sec:hec} (an instance of \cref{ana:thm:onesided}) then shows that an
eventually nonnegative $J_\tau$ forces (i). An eventually nonpositive $J_\tau$ forces (i)
in the same way, by applying the argument to $-J_\tau$; for $\tau=4$ this would contradict
(iii), which proves the last assertion.
\end{proof}

The Riemann hypothesis for $L(s,\rho)$ is not a consequence of \cref{src:thm:dgrh}, and
establishing the sign in (ii) is the open step. The obstruction to replacing the cubic
weight by congruence weights is treated in
\cref{sec:hec}. \Cref{zer:prop:eta} gives one formula for all positive eta products; earlier
computations reach much farther for this particular field \cite{Ng2000Limiting}. The
quantity $2\zerLL'(1)/\zerLL(1)$ has also been enclosed numerically: with
$G(t)=f(it/\sqrt{23})$ and $G(1/t)=tG(t)$,
\[
 \zerLL(s)=\int_1^\infty G(t)\bigl(t^{s-1}+t^{-s}\bigr)\,\dd t,
\]
so $2\zerLL'(1)/\zerLL(1)$ is twice the ratio of the integrals of $G(t)\log t\,(1-1/t)$
and $G(t)(1+1/t)$, and \artifact{grh-transfer/checks/cubic23.py} encloses it in
$(0.32770,0.35161)$, consistently with the bound $1/c=0.39823\ldots$. The enclosure
evaluates the positive eta product on intervals of width $1/64$ up to $16$, using the first
$12$ factors of each infinite product; for omitted factors $1-u_i$ their product lies
between $1-\sum u_i$ and $1$, and for the integral tails one uses
$0<G(t)\le e^{-2\pi t/\sqrt{23}}$ and $\log t\le t$. These integrals and enclosures hold
independently of any hypothesis on zeros.

\begin{remark}[Sample values]\label{zer:rem:samples}
Rounded values of \eqref{zer:eq:Jh} at $\tau=4$ illustrate the scale and the sign; the
proofs do not extrapolate them. The primes $p=2,5,59$ respectively exhibit the irreducible,
mixed and completely split cases of the cubic. The root counts are exact and the weighted
sums use floating point; the output is in \artifact{grh_transfers/checks/cubic_check.txt}.
\begin{center}\small
\begin{tabular}{rrrrrr}\toprule
$x$ & Three roots & One root & No roots & $\sqrt x\,J_4$ cubic & $\sqrt x\,J_4$ mod $4$\\\midrule
100&1&14&9&1.343855&0.674377\\
1000&25&87&55&0.9648&0.747457\\
10000&193&627&408&1.150540&0.961868\\
100000&1567&4823&3201&1.1356&1.032079\\\bottomrule
\end{tabular}\end{center}
\end{remark}

\subsection{Zhao's Mertens-error question}\label{zer:sec:zhao}
\sources{grh-transfer}

Zhao asks for a classification of number fields related to the sign of the integrated
prime-ideal Mertens error \cite[\S8(3)]{Zhao2025Mertens}. His sufficient condition
compares $2$ with a sum over the nontrivial zeros of $\zeta_K$. With the Riemann hypothesis
for $\zeta_K$ in force, write
\begin{equation}\label{zer:eq:BK}
 B_K=\sideset{}{^*}\sum_{\zeta_K(1/2+i\gamma)=0}\frac1{1/4+\gamma^2},
\end{equation}
where the star means that zeros are counted with multiplicity, with all ordinates of both
signs and any central zero included. For every finite abelian extension $K/\Q$,
\cref{src:thm:dgrh} supplies the Riemann hypothesis for $\zeta_K$ through the factorization
of $\zeta_K$ into Dirichlet $L$-functions. The numerical condition can be resolved in this
entire class.

\begin{theorem}[\stG]\label{zer:thm:zhao}
Up to isomorphism over $\Q$, exactly $229$ finite abelian extensions $K/\Q$ satisfy
$B_K<2$. Their distribution is given in the following table. For all $229$ fields the
integral of the reciprocal prime-ideal Mertens error is strictly positive for every
sufficiently large upper endpoint.
\end{theorem}
\begin{center}\small\begin{tabular}{@{}lrl@{}}\toprule
Galois group & Number & Identification or conductors\\\midrule
Trivial & 1 & $\Q$\\
$C_2$ & 168 & 116 real, 52 imaginary\\
$C_3$ & 7 & $7,9,13,19,31,43,63$\\
$C_2\times C_2$ & 38 & 10 totally real; generator pairs in \cref{zer:tab:biquadratic}\\
$C_4$ & 9 & $5,13,15,16,16,17,20,35,40$\\
$C_5$ & 1 & $\Q(\zeta_{11})^+$\\
$C_6$ & 4 & $7,9,13,35$\\
$C_2\times C_4$ & 1 & $\Q(\zeta_{15})$\\\bottomrule
\end{tabular}\end{center}
Here $C_m$ denotes a cyclic group of order $m$, and the plus sign denotes the maximal real
subfield. Conductors can identify several fields: the file
\artifact{grh-transfer/checks/abelian_fields.json} specifies each of the $229$ fields
uniquely, through generators of its character group, with an enclosing interval for $B_K$;
\artifact{grh-transfer/checks/fields.json} lists every quadratic discriminant and all three
quadratic subfield discriminants of each biquadratic field. Among the cyclic cubic fields
of conductor $63$, the one selected has a generator with local angle $1/3$ at both $3$ and
$7$ in the convention of the proof. Both cyclic quartic fields of conductor $16$ pass. The
distinct full cyclotomic fields passing the test have minimal conductors
$1,3,4,5,7,8,9,12,15$; in particular $\Q(\zeta_{15})$ has $1.9967391<B_K<1.9967467$. All
abelian fields of degree greater than eight fail. The classification establishes the stated
sufficient condition; the sign of the prime-ideal average for fields with $B_K\ge2$ is a
separate question.

\begin{table}[ht]\centering\small
\begin{tabular}{r|l}
$d_1$ & choices of $d_2$\\\hline
$-3$ & $-23,-20,-19,-11,-8,-7,-4,5,8,13$\\
$-4$ & $-24,-15,-11,-8,-7,5,17$\\
$5$ & $-11,-8,-7,8,12,13,17,21$\\
$-7$ & $-15,-8,8,17$\\
$-8$ & $-20,-11,12$\\
$8$ & $12,17,28$\\
$12$ & $-15,21,33$
\end{tabular}
\caption{The $38$ biquadratic fields $\Q(\sqrt{d_1},\sqrt{d_2})$ with $B_K<2$. Each entry
in the second column gives a separate field.}\label{zer:tab:biquadratic}
\end{table}

\begin{proof}
\emph{Reduction to characters.} By the Kronecker--Weber theorem, every abelian $K/\Q$ is a
subfield of a cyclotomic field. Galois duality identifies it uniquely with the finite group
$V$ of Dirichlet characters trivial on the subgroup fixing $K$. Take every character at its
primitive conductor when forming its $L$-function; products of characters mean products
followed by passage to the inducing primitive character. Factoring the Dedekind zeta
function and adding zero multiplicities gives
\begin{equation}\label{zer:eq:BKsum}
 B_K=B_\Q+\sum_{1\ne\chi\in V}B_\chi,\qquad B_\Q=2+\zergE-\log(4\pi),
\end{equation}
where $B_\chi$ is the same zero sum for $L(s,\chi)$; the formula for $B_\Q$ follows by
logarithmically differentiating the completed zeta function at $1$. Every summand is
positive. A real nontrivial character of a qualifying field must satisfy
$B_\chi<2-B_\Q$. A nonreal one comes paired with its conjugate of equal weight, since the
ordinates for $\bar\chi$ are the negatives of those for $\chi$, and must satisfy
$B_\chi<(2-B_\Q)/2$. The next two steps find every character satisfying the appropriate
inequality; this permits a finite exhaustive group calculation.

\emph{Quadratic characters: an exterior cutoff.} Label a quadratic character by the
fundamental discriminant $d$ of its field, put $q=\abs d$, and write $a=0$ for an even
character and $a=1$ for an odd one. The real part of the logarithmic derivative of the
completed primitive $L$-function at $2$ is the sum of $(3/2)/((3/2)^2+\gamma^2)$ over the
ordinates of its zeros; this follows by subtracting the Hadamard logarithmic derivatives at
$2$ and $-1$ and using the functional equation. Each term is at most
$(3/2)/(1/4+\gamma^2)$. Consequently
\[
 B_{\chi_d}\ge\frac13\Bigl[\log(q/\pi)+\psi\bigl((2+a)/2\bigr)+2\frac{L'}L(2,\chi_d)\Bigr],
\]
where $\psi=\Gamma'/\Gamma$. The digamma term is at least $-\zergE$. The Euler product
gives $L'/L(2,\chi_d)\ge\zeta'/\zeta(2)$, and $-\zeta'/\zeta(2)<0.577$: sum
$\Lambda(k)/k^2$ up to $1000$ and bound the tail by
$\int_{1000}^\infty(\log t)/t^2\,\dd t$.
The displayed bound exceeds $2.0371$ for $q\ge8000$, including both parities.

\emph{Quadratic characters: the finite range.} This leaves $4{,}861$ nontrivial
fundamental discriminants with $\abs d<8000$. They are divided at the extrema $-551$ and
$1201$ reported by Zhao. The $534$ interior values recover the $178$ characters with
$B_{\chi_d}<2$ in his study. A separate computation treats all $4{,}327$ exterior values
in the finite range: every lower endpoint exceeds $2.0442$
(\artifact{grh-transfer/checks/exterior-results.json}). Thus completeness is independent
of his numerical search. The numerical comparison is made without locating zeros: Zhao's
equation (32) is
\begin{equation}\label{zer:eq:zhao32}
 B_{\chi_d}=\log(q/\pi)-\zergE-2\,\mathbf 1_{d>0}\log2+2\,\frac{L'(1,\chi_d)}{L(1,\chi_d)} .
\end{equation}
Evaluate the sums with weights $f(t)=1/t$ and $g(t)=\log(t)/t$, respectively, for $L$ and
$-L'$. The computation sums $m=24$ periods and then uses Euler--Maclaurin separately on
each residue class, starting at $t=mq+r$, $1\le r\le q$. For either weight $w$ with
primitive $\mathcal W$, the three retained terms for the tail are
\[
 -\mathcal W(t)/q+w(t)/2-q\,w'(t)/12 .
\]
The common integration constant cancels because $\sum_r\chi_d(r)=0$. With the $B_4$
endpoint term absorbed in the error, the omitted terms have total absolute value per
progression at most $q^3(\abs{w'''(t)}+\int_t^\infty\abs{w^{(4)}(y)}\,\dd y)/720$, by
$\abs{B_4(\set y)}\le1/30$. In the range used here $w^{(4)}>0$ and $w'''<0$, so the
integral equals $\abs{w'''(t)}$. Summing these bounds over $r$ gives errors at most
$1/(60m^4)$ for $f$ and $\log(mq+q)/(60m^4)$ for $g$. The interior computation uses
$30$-decimal-place interval arithmetic and compares an interval with $2$ only when they are
disjoint. The exterior one uses twelve periods, fixed-point integer sums with rounding
allowances, and interval arithmetic for the final enclosure: with $P=12$, $S=2^{42}$ and
$n=Pq$, the integer weights approximate
$1/t$, $(\log t)/t$, $\log t$, $(\log t)^2$, $1/t^2$, $(\log t-1)/t^2$ to within $2/S$ for
$1\le t\le104000$, so that for either sum the rounded coefficients and the finite prefix
contribute at most $(2/S)(n+1+q/2+q^2/12)$ in addition to the stated tail error; all dot
products are signed $64$-bit integer operations, since the sum of absolute values for any
correction row is at most $144Sq<2^{63}$ and for a full prefix row at most $nS<2^{63}$.
Interval substitution in \eqref{zer:eq:zhao32} closes each test, including positivity of
the denominator.

\emph{Multiquadratic fields.} Start with the $168$ singleton characters for which
$B_\Q+B_{\chi_d}<2$. Given a qualifying subgroup $V$ of $\Q^\times/(\Q^\times)^2$, adjoin
a character $a$ by taking $V\cup aV$. Retain the larger subgroup only if each character
lies in the allowed set and the total is below $2$. If a subgroup fails, positivity
excludes every overgroup. Every finite elementary abelian $2$-group has a flag of
subgroups with successive index two, so this enumeration is exhaustive. It produces the
$38$ fields of \cref{zer:tab:biquadratic}, then no rank-three group. Thus exactly $207$
multiquadratic fields qualify: the rational field, $168$ quadratic and $38$ biquadratic
fields. One near miss illustrates why a certified comparison matters: for the biquadratic
field with subfield discriminants $-15,-24,40$ the computed interval lies strictly inside
$(2.0000952,\,2.0000983)$.

\emph{Nonreal characters.} For a primitive character $\chi$ of conductor $q$ and parity
$a\in\set{0,1}$, the real part of the Hadamard logarithmic derivative of the completed
$L$-function at $s=2$ gives
\[
 \frac32B_\chi\ \ge\ \sum_\rho\frac{3/2}{9/4+\gamma^2}
 =\frac12\log(q/\pi)+\frac12\psi(1+a/2)+\Re\frac{L'}L(2,\chi)
 \ \ge\ \frac12\log(q/\pi)-\frac{\zergE}2+\frac{\zeta'}\zeta(2).
\]
Here $\psi$ is increasing on the positive real axis and the last inequality follows from
the absolutely convergent Euler series. The bound $-\zeta'(2)/\zeta(2)<0.6$ follows by
summing $\sum_{n\le1000}\Lambda(n)/n^2$ and bounding the tail by $(1+\log1000)/1000$. For
$q\ge400$ the displayed lower bound implies $B_\chi>1$, which excludes this character.

For $q<400$, primitive local characters are generated by the Chinese remainder theorem. At
odd $p$ take the least primitive root modulo $p^2$, which also generates the units modulo
every $p^e$, and write its character value as $\exp(2\pi ij/\varphi(p^e))$. A nontrivial
character at $e=1$ is primitive locally, and for $e\ge2$ primitivity requires $p\nmid j$.
At $2^e$, $e\ge3$, use $-1$ and $5$ and require the exponent numerator on $5$ to be odd
when expressed with denominator $2^{e-2}$. Modulo $4$ there is one primitive local
character and modulo $2$ there are none. This generates each primitive character exactly
once. Independently the Dirichlet convolution $\mu*\varphi$ counts $29405$ primitive
characters of conductor below $400$, including the trivial conductor-one character;
applying the same inversion to the size of the two-torsion of the unit group counts $243$
real characters, leaving $29162$ nonreal characters. For the finite comparisons use the
real part of the same logarithmic derivative at $s=1$, or equivalently the functional
equation, to obtain
\[
 B_\chi=\log(q/\pi)-\zergE-2(1-a)\log2+2\Re\frac{L'}L(1,\chi).
\]
The real and imaginary parts of $L(1,\chi)$ and $-L'(1,\chi)$ follow from the same
Euler--Maclaurin progression sums as in the quadratic calculation, now using $32$ complete
periods. For a transparent rounding bound choose $S=2^{34}$. For each unit residue $r$ let
$w_r/S$ approximate the Euler--Maclaurin weight and $h_r/S$ the coordinate of its root of
unity, with absolute errors at most $1/S$ verified in $32$-decimal-place interval
arithmetic, and accumulate the dot product $\sum_rw_rh_r$ in exact integer arithmetic. The
absolute rounding error after dividing by $S^2$ is at most
$\bigl(\varphi(q)(S+1)+\sum_r\abs{w_r}\bigr)/S^2$, since $\abs{h_r}\le S$. Add the analytic
tail error to each component and use complex interval division to enclose $B_\chi$. All
comparisons with the budget separate it from the computed interval. The file
\artifact{grh-transfer/checks/nonreal.json} records the $394$ qualifying characters ($197$
conjugate pairs) by their local angles with rigorous decimal bounds. Their largest
conductor is $109$.

\emph{Exhaustion of the character groups.} Encode a character by its nonzero local values
on the fixed generators, as angles in $\Q/\Z$. At each odd prime these generator choices
are compatible through the prime-power tower. At $2$ use one coordinate for $-1$ and
another for $5$. Thus adding angle vectors modulo one gives the primitive product,
including products of different conductors. There are $168$ admissible quadratic
characters and $394$ admissible nonreal characters, with the identity adjoined. For each
single generator form its cyclic subgroup. Discard it if any power is absent or its total
lower weight reaches $2$. Keep one generator of each remaining cyclic subgroup. Start at
the identity subgroup, and repeatedly adjoin each retained generator, forming the union of
the distinct cosets until it closes. A group with an inadmissible member or total lower
weight at least $2$ can be discarded together with its overgroups. Each qualifying
character group has a generating sequence whose every intermediate subgroup qualifies, so
this procedure reaches all of them. Canonical sets of angle vectors prevent multiple
counting. The sums of interval endpoints are compared with $2$ by exact rational
arithmetic: every tested full group either has lower endpoint at least $2$ or upper
endpoint strictly below $2$. The resulting group counts are those of the table. The
conductor assigned to a field is the least common multiple of its members' conductors.
Two Galois subfields of $\overline\Q$ that are isomorphic over $\Q$ coincide, which
justifies counting the distinct character subgroups themselves.

\emph{From the zero sum to the sign of the mean.} The sums run over prime ideals of $K$,
each ideal separately. Let $E_{1,K}(x)$ be the sum of $\log\Nm\mathfrak p/\Nm\mathfrak p$
for $\Nm\mathfrak p\le x$, less $\log x$ and its limiting constant, and define $E_{2,K}$
in the same way using $1/\Nm\mathfrak p$ and $\log\log x$. Under RH for $K$, the integrated
explicit formula used in the proof of \cref{zer:thm:endpoint} below gives
\[
 \frac{F(x)}{\sqrt x}=2+\sum_{\rho=1/2+i\gamma}\frac{e^{i\gamma\log x}}{1/4+\gamma^2}+o_K(1),
 \qquad F(x)=\int_2^xE_{1,K}(y)\,\dd y .
\]
Indeed integration of the first Mertens error gives
$-x\int_x^\infty(\theta_K(t)-t)t^{-2}\,\dd t+O_K(1)$, where $\theta_K$ counts prime ideals
with weight $\log\Nm\mathfrak p$. Replacing $\theta_K$ by $\psi_K$ gives the displayed zero
series, as in formula (1.3a) of \cite{GrenieMolteni2016Smoothed}. Removing prime-ideal
squares restores the term $2$, since $\theta_K(\sqrt t)\sim\sqrt t$ by the prime ideal
theorem; higher powers contribute $o(\sqrt t)$ before integration. This square term counts
all prime ideals and so differs from the completely split rational-prime term of
\cref{zer:thm:endpoint}.

Thus $F(x)=O_K(\sqrt x)$ and $F(x)/\sqrt x\ge2-B_K+o_K(1)$. To pass to the reciprocal sum,
use the exact identity of Stieltjes differentials
$\dd E_{2,K}(t)=\dd E_{1,K}(t)/\log t$. Integration to infinity gives
\[
 E_{2,K}(x)=\frac{E_{1,K}(x)}{\log x}-\int_x^\infty\frac{E_{1,K}(t)}{t\log^2t}\,\dd t .
\]
The last integral is $O_K(x^{-1/2}/\log^2x)$ by integration by parts against $\dd F(t)$.
Integrating in $x$ and integrating the first term by parts once more yields
\[
 \int_2^xE_{2,K}(t)\,\dd t=\frac{F(x)}{\log x}+O_K\Bigl(\frac{\sqrt x}{\log^2x}\Bigr).
\]
It follows that the reciprocal prime-ideal Mertens error has strictly positive integral for
all sufficiently large $x$ whenever $B_K<2$.
\end{proof}

For $B_K\ge2$ the zero phases can restrict the lower endpoint more strongly than the
triangle inequality, so excluding such fields from the numerical test leaves their eventual
sign undecided.

\begin{remark}[Which step needs the input theorem]\label{zer:rem:needsGRH}
The finite calculation has an interpretation before any zero-location theorem is assumed.
Set $A_\chi=2\Re(\Lambda'/\Lambda)(1,\chi)$ for the completed function $\Lambda(s,\chi)$ of
a nonprincipal primitive character, with the harmless factors at $0$ and $1$ included for
the zeta factor. In the Hadamard formula the functional equation pairs $\beta+i\gamma$ with
$1-\beta+i\gamma$. At $s=1$ the two real logarithmic-derivative terms together are at
least $1/(1+\gamma^2)$; at $s=2$ they are at most $3/(1+\gamma^2)$, by bounding the two
denominators and adding their numerators, with fixed points counted with half weight when
making pairs. Thus $A_\chi\ge(2/3)\Re(\Lambda'/\Lambda)(2,\chi)$ even before localization,
and $A_\chi>0$. The conductor cutoffs and the finite comparisons therefore classify the
corresponding completed logarithmic-derivative constant on their own.
\Cref{src:thm:dgrh} makes $A_\chi$ equal to the absolute zero sum $B_\chi$ and makes the
statement about the Mertens integral possible.
\end{remark}

\begin{proposition}[A finite frontier for arbitrary number fields; \stU, with RH for
$\zeta_K$ as a hypothesis]\label{zer:prop:frontier}
Let $K$ be a number field of signature $(r_1,r_2)$ and degree $n=r_1+2r_2$ whose Dedekind
zeta function satisfies the Riemann hypothesis. If $B_K<2$, then
\[
 n\le8,\qquad\abs{D_K}<193\,000\,000 .
\]
In degree $8$ only the signatures $(0,4)$, $(2,3)$ and $(4,2)$ survive.
\end{proposition}

The degree bound is sharp, since $\Q(\zeta_{15})$ attains it. By Hermite's theorem the
proposition leaves a finite search for the fields satisfying Zhao's sufficient condition,
subject to their own zeta-function hypotheses.

\begin{proof}
Write $u=s-1/2$ and let $\xi_K(s)$ be the completed Dedekind zeta function multiplied by
$s(s-1)$. The paired Hadamard formula on real $s>1$ gives
$(\log\xi_K)'(s)=\sum_\rho u/(u^2+\gamma^2)$. Its derivative converges locally uniformly.
Since $(\gamma^2-u^2)/(u^2+\gamma^2)^2\le(1/4+\gamma^2)^{-1}$, we have
$B_K\ge(\log\xi_K)''(2)$. The Euler product gives $(\log\zeta_K)''(2)\ge0$. The
pole-removal and gamma factors therefore give
\[
 B_K\ \ge\ \frac{\pi^2}{24}r_1+\Bigl(\frac{\pi^2}6-1\Bigr)r_2-\frac54 .
\]
For fixed degree the right side increases with $r_1$, since two real places contribute
more than one complex place. Direct substitution gives the degree bound $n\le10$. To
sharpen it, set $s=5/2$, $u=2$, $c=47/48$ and $\Phi(u,t)=u/(u^2+t)$. For every $t\ge0$,
\[
 (t+1/4)^{-1}\ge\Phi_u(u,t)-c\,\Phi_{uu}(u,t).
\]
Indeed after multiplication by $(t+1/4)(t+4)^3$ the difference has numerator
$(64+13c)t+68+4c$; the coefficient of $t^2$ vanishes for this choice of $c$. Apply the
inequality to $t=\gamma^2$. Each prime-power term in
$(\log\zeta_K)''-c(\log\zeta_K)'''$ is nonnegative, giving
\[
 B_K\ge A\,r_1+H\,r_2-J>0.4618\,r_1+0.7216\,r_2-1.3101 ,
\]
where $A=\psi'(s/2)/4-c\psi''(s/2)/8$, $H=\psi'(s)-c\psi''(s)$ and
$J=s^{-2}+(s-1)^{-2}+2c(s^{-3}+(s-1)^{-3})$. For the numerical lower bounds on $A,H$ one
may sum the first $100{,}000$ positive terms of $\psi'(x)=\sum_{j\ge0}(x+j)^{-2}$ and
$-\psi''(x)=2\sum_{j\ge0}(x+j)^{-3}$. At degree $9$ even the smallest possible signature
contribution, for $(1,4)$, exceeds $2$ after subtraction of $J$. At degree $10$ the pair
$(0,5)$ already exceeds $2$. At degree $8$ at least six real places would give a
contribution above $2$.

For the discriminant, use $(\log\xi_K)'(4)\le(7/2)B_K$. The local Euler factors imply
$\zeta_K'/\zeta_K(4)\ge n\,\zeta'/\zeta(4)$, since the coefficient at $p^k$ in
$-\zeta_K'/\zeta_K$ is at most $n\log p$. With $c_4=-(\zeta'/\zeta)(4)$ this gives
\[
 \log\abs{D_K}<\frac{77}6+r_1\bigl(\log\pi-\psi(2)+2c_4\bigr)
 +r_2\bigl(2\log(2\pi)-2\psi(4)+4c_4\bigr)<\frac{77}6+0.850\,r_1+1.420\,r_2 .
\]
Here $c_4<0.0638$, by summing $\Lambda(m)/m^4$ through $100$ and bounding the tail by
$\int_{100}^\infty(\log t)t^{-4}\,\dd t$, and $\psi(2)=1-\zergE$,
$\psi(4)=11/6-\zergE$. The largest resulting right side over the allowed signatures is
$77/6+4(0.850)+2(1.420)<\log(193\,000\,000)$. Both bounds use only the functional equation,
the Euler product and RH for $K$.
\end{proof}

\begin{example}[An explicit nonabelian case on its own hypothesis]\label{zer:ex:cubic23}
Let $F$ be the cubic field defined by $x^3-x-1$, of discriminant $-23$, and let $E$ be
its splitting field. With $L_f$ the $L$-series of $f(z)=\eta(z)\eta(23z)$, the permutation
representation on the three roots and the regular representation give
$\zeta_F=\zeta L_f$ and $\zeta_E=\zeta L_{-23}L_f^2$
\cite{Zagier1991Polylogarithms}. The enclosure $(0.32770,0.35161)$ of
$2\zerLL'(1)/\zerLL(1)$ recorded at the end of \cref{zer:sec:window}, together with
$B_{\chi_{-23}}\in(0.56930,0.56931)$ from the quadratic calculation, yields
\[
 0.3738<B_F<0.3979,\qquad1.2708<B_E<1.3188
\]
under the Riemann hypothesis for the respective zeta functions. That hypothesis identifies
the logarithmic derivatives with Zhao's positive spectral sums; in this $S_3$ case the
remaining analytic input is the localization of the zeros of the single degree-two Artin
factor $L_f$.
\end{example}

\subsubsection{The exact upper limiting value of the split-prime error}
Let $M/\Q$ be finite Galois with group $G$ and degree $n$, and let $P_M$ be the set of
unramified rational primes that split completely in $M$. Put
\[
 R_M(x)=\sum_{p\in P_M,\ p\le x}\log p\Bigl(\frac1p-\frac1x\Bigr)-\frac{\log x}n-c_M,
\]
where $c_M$ is the unique constant for which $R_M(x)\to0$; it exists by the prime ideal
theorem with its classical error term (\cref{sec:grp}). Write
\[
 b_M=\frac{\#\set{g\in G:g^2=1}}n,
\]
the identity being included in the count.

\begin{theorem}[Endpoint; \stU]\label{zer:thm:endpoint}
If $\zeta_M$ satisfies the Riemann hypothesis, with nontrivial zeros
$\rho=\tfrac12+i\gamma$ listed with multiplicity, then
\[
 \limsup_{x\to\infty}\sqrt x\,R_M(x)=2b_M+\frac1n\sum_\rho\frac1{1/4+\gamma^2}.
\]
If the Riemann hypothesis fails for $\zeta_M$, the same limsup is $+\infty$.
\end{theorem}

In particular the one-sided criterion for completely split primes (\cref{sec:grp}) can use
the exponent $-1/2$ exactly. For abelian $M$, \cref{src:thm:dgrh} and the abelian
factorization supply the first case.

\begin{proof}
Put $\theta(t)=\sum_{p\le t,\ p\in P_M}\log p$. Integrating this step function gives the
ramp sum as $\int_1^x\theta(t)t^{-2}\,\dd t$. Thus
\[
 R_M(x)=-\int_x^\infty\bigl(\theta(t)-t/n\bigr)t^{-2}\,\dd t .
\]
Under RH for $\zeta_M$ this integral converges absolutely by the prime ideal theorem. For
the moment replace $\theta$ by $\psi_M/n$, with $\psi_M$ the sum of $\log\Nm\mathfrak p$
over prime ideal powers of norm at most $t$. The integrated explicit formula of
\cite[(1.3a) and Lemma~3.3]{GrenieMolteni2016Smoothed} gives for the fixed field
\[
 A(t):=\int_0^t\bigl(\psi_M(y)-y\bigr)\,\dd y
 =-\sum_\rho\frac{t^{\rho+1}}{\rho(\rho+1)}+O_M(t\log t)\qquad(t\ge3).
\]
The completed Dedekind zeta function has order one, so its zeros have
$\sum_\rho(1+\abs\rho^2)^{-1}<\infty$. Under RH the series for $A(t)$ is bounded absolutely
by a constant times $t^{3/2}$. Integrating by parts and then integrating that series term
by term gives
\begin{align*}
 -\int_x^\infty\bigl(\psi_M(t)-t\bigr)t^{-2}\,\dd t
 &=\frac{A(x)}{x^2}-2\int_x^\infty\frac{A(t)}{t^3}\,\dd t\\
 &=\sum_\rho\frac{x^{\rho-1}}{\rho(1-\rho)}+O_M\Bigl(\frac{\log x}x\Bigr),
\end{align*}
the coefficient of $x^{\rho-1}$ on the first line being
$-\bigl(1-2/(1-\rho)\bigr)/(\rho(\rho+1))$.

For unramified $p$ with Frobenius $g$, the coefficient of $p^j$ in $\psi_M/n$ is $\log p$
when $g^j=1$, and zero otherwise; this follows by counting the $n/f$ primes above $p$,
each of norm $p^f$, for $f=\ord(g)$. The $j=1$ terms give $\theta$. The $j=2$ terms sum
to $b_M\sqrt t+o_M(\sqrt t)$ by the Chebotarev density theorem and partial summation. The
terms $j\ge3$ contribute $O(t^{1/3}\log t)$, since $\sum_{p\le y}\log p=O(y)$ and
$j\le\log_2t$. The finitely many ramified primes contribute $O_M(\log t)$. Consequently
\begin{equation}\label{zer:eq:endpointexpansion}
 \sqrt x\,R_M(x)=2b_M+\frac1n\sum_\rho\frac{e^{i\gamma\log x}}{1/4+\gamma^2}+o_M(1).
\end{equation}
Here $2b_M$ comes from integrating $b_M\sqrt t/t^2$ from $x$ to infinity, and the sum is
real by conjugation. The triangle inequality gives the claimed upper endpoint. To attain it
in the limit, truncate the absolutely convergent series and apply simultaneous Dirichlet
approximation to the finitely many numbers $\gamma/(2\pi)$: there are arbitrarily large
integers $k$ with every $e^{i\gamma k}$ as close to $1$ as desired. Put $x=e^k$ and let the
tail bound tend to zero. This requires no independence assumption on the ordinates.

Suppose finally that the limsup is finite while RH fails for $\zeta_M$. Then
$R_M(x)\le Cx^{-1/2}$ for all large $x$. For $\Re s>1$,
\[
 \int_1^\infty R_M(x)x^{-s}\,\dd x=\frac{D_M(s)}{s(s-1)}-\frac{1/n}{(s-1)^2}
 -\frac{c_M}{s-1},\qquad
 D_M(s)=\sum_{p\in P_M}\frac{\log p}{p^s}=-\frac1n\frac{\zeta_M'}{\zeta_M}(s)+H_M(s),
\]
with $H_M$ holomorphic in $\Re s>1/2$, by the count of coefficients above. The right side
is regular at $s=1$, and each zero $\rho$ of $\zeta_M$ of multiplicity $m_\rho$ with
$1/2<\Re\rho<1$ gives a simple pole with residue $m_\rho/[n\rho(1-\rho)]$, which is
positive at a real zero. The one-sided Mellin--Landau test (\cref{ana:thm:onesided};
see \cref{sec:grp} for this weight) then excludes every such zero, and the functional
equation supplies the other half of the critical strip, a contradiction.
\end{proof}

\begin{corollary}[A cyclotomic specialization; \stG]\label{zer:cor:cyclotomic}
Let $q\ge1$, write $C(q,1)=\#\set{a\in(\Z/q\Z)^\times:a^2=1}$, and let
\[
 E_1(t;q,1)=\sum_{p\le t,\ p\equiv1\bmod q}\frac{\log p}p-\frac{\log t}{\varphi(q)}-e_{q,1},
 \qquad\lim_{t\to\infty}E_1(t;q,1)=0,
\]
be the first Mertens error in the residue class of $1$. Then
\[
 \limsup_{x\to\infty}\frac1{\sqrt x}\int_2^xE_1(t;q,1)\,\dd t
 =\frac1{\varphi(q)}\Bigl(2C(q,1)+\sum_\rho\frac1{1/4+\gamma^2}\Bigr),
\]
the sum being over the nontrivial zeros of $\zeta_{\Q(\zeta_q)}$ with multiplicity.
\end{corollary}
\begin{proof}
For $M=\Q(\zeta_q)$ the counted primes, up to finitely many exceptions, are
$p\equiv1\pmod q$; centering separately absorbs the constant from any omitted primes.
Integration gives
\[
 \frac1x\int_2^xE_1(t;q,1)\,\dd t=R_M(x)+O_q(1/x).
\]
Apply
\cref{zer:thm:endpoint}, whose hypothesis holds by \cref{src:thm:dgrh}.
\end{proof}

Zhao treats residue-class versions of these mean errors and gives their explicit formula
\cite[Theorem~4 and \S7]{Zhao2025Mertens}, and also the exact rational-prime case
$M=\Q$. For the identity residue class the upper endpoint follows with no
linear-independence hypothesis, since every character equals one there.
\Cref{zer:thm:endpoint} extends the calculation to arbitrary finite Galois extensions
through the count of the elements whose square is the identity.

\begin{example}\label{zer:ex:Qi}
For $M=\Q(i)$ the limit superior for the centered sum over $p\equiv1\pmod4$ is
\[
 3+\frac{\beta'(1)}{\beta(1)}-\log\pi=2.10087969891685\ldots .
\]
Indeed $n=2$, $b_M=1$, and the completed entire function is proportional to
$\xi_M(s)=s(s-1)4^{s/2}(2\pi)^{-s}\Gamma(s)\zeta(s)\beta(s)$. The Hadamard product and the
symmetry around $1/2$ give
$\sum_\rho[\rho(1-\rho)]^{-1}=\xi_M'(1)/\xi_M(1)-\xi_M'(0)/\xi_M(0)=2\xi_M'(1)/\xi_M(1)$.
Evaluating the logarithmic derivative at $1$ cancels the Euler constants from $\Gamma(s)$
and $(s-1)\zeta(s)$.
\end{example}

\subsection{Open questions}\label{zer:sec:open}
{\emergencystretch=6em\sources{grh-transfer; grh_transfers; dirichlet-grh-research; exponent_seven_and_grh_transfers; grh-torsion-and-zero-data; grh-methods-polynomial-density; grh_cm_four_terms_and_prime_bias; grh_method_transfers_2026-11-26; grh_multiscale_covariance; multiscale_prime_fluctuations; grh_spectral_rigidity; grh_filtered_prime_races; grh_prime_error_balancing; prime_peak_rigidity; sparse_grh_balanced_windows; sharp_multiscale_prime_windows; prime_scale_coalescence_2026-11-26}}

\begin{enumerate}
\item \emph{Fixed-height Euler products.} Can $F_0(v)/v$ in \eqref{zer:eq:single} be
controlled for every $v$, that is, is $\Psi(x)=o(\sqrt x\log x)$ for a fixed primitive
character? This is the term that must tend to zero to reach the fixed-height prediction.
Repeating the endpoint cancellation with more shifts gives additional linked identities
but supplies no upper bound on it.
\item \emph{The boundary window scale.} In windows $[U,U+\exp(-a\sqrt U)]$,
\cref{zer:thm:exceptional} gives a vanishing exceptional fraction for tolerances
$\eps>a^2/(2\pi)$. Does the fraction vanish for every fixed tolerance, and is the constant
$2\pi$ optimal? An improvement would have to retain some phase information of the zeros
below the cutoff or use different arithmetic input.
\item \emph{Short intervals.} Does a bounded-length interval suffice for each directional
multiplicity excess for a fixed pair of characters, and is the coefficient $2/\pi$ in
\cref{zer:thm:short} optimal? Can one force a genuinely unshared zero in every
triple-logarithmic interval? The first asks for a stronger localization argument, the last
for information on individual orders at common zeros.
\item \emph{The boundary value in Suzuki's question.} Does the ordinary logarithmic
density of $\set{x:E(x)\le0}$ exist, and what is its value? An atom-free limiting
distribution at the threshold $c$ would give it; a quantitative small-interval estimate at
$c$ surviving the infinite-tail limit would address the missing point.
\item \emph{Suzuki's averages over moduli.} What is the sign of $A_Q(x)$ in the
intermediate regime $Q\to\infty$ with $Q=o(x)$? Can the exceptional set of finite
logarithmic measure in \cref{zer:prop:UQ} be replaced by an eventual pointwise statement
at fixed $Q$?
\item \emph{Effective constants.} Can the radius $r_q$ and ratio $C_q$ of
\cref{zer:thm:balancing}, the window $H$ and radius of \cref{zer:cor:smoothbalance}, the
constants $c,C,K$ of \cref{zer:thm:coercive}, the least eigenvalues in
\cref{zer:thm:jets,zer:thm:filtered} and a certified window length $L_0$, and the return lengths
near the peak in \cref{zer:thm:peak}, be made explicit, ideally uniformly in $q$, in the number
of scales, or in the cut $H$? Near-coincident ordinates must then be separated rigorously, and
efficient probe primes chosen.
\item \emph{The eigenvalue gap.} The two bounds of \eqref{zer:eq:coercive} differ by three powers
of $\log T$. Can information on pair correlation and multiplicities close the gap, and what are
the sharp asymptotics for arbitrary shifts under stronger hypotheses?
\item \emph{Merging scales.} \Cref{zer:thm:coalesce} settles the rate for $\zeta$. What is the
corresponding law for the full residue-class matrix of \cref{zer:cor:multiscale}, where coincident
zeros of different $L$-functions add structure? Is the factor $\tfrac43$ optimal, and does
$M_2(T)\sim\kappa T\log T/(2\pi)$ hold, as \cref{zer:prop:kappa} requires?
\item \emph{Beyond covariance.} Can the convex-hull statements be upgraded to the interior of the
actual limiting support, to every ordering of the unsmoothed counts, or to inclusivity of
ordinary prime races, without a hypothesis on the ordinates? Which nonlinear constraints hold on
the cluster sets? Positive covariance excludes affine hyperplanes only.
\item \emph{Less smoothing and other families.} Does strict positivity of the multiscale or jet
covariance survive for unsmoothed prime counts, where only a mean-square expansion is available,
after the prime-square biases are removed? Does \cref{zer:thm:polyrigid} extend to families with
other Gamma factors and distinct Euler coefficients, and do the probe arguments extend to
Frobenius-class functions of Artin $L$-functions under their own Riemann hypothesis (which
\cref{src:thm:dgrh} does not provide)?
\item \emph{Kernels and filters.} Which compact filters retain every noncentral frequency with the
shortest memory or the smallest error, and which other Mellin kernels have nonnegative spectral
weights on the critical line, with sharper laws or shorter recurrence? Can the compact window
kernel of \cref{zer:thm:sharpwin} be optimized for given contrasts, or constrained so that the pole
terms cancel, and do its finite Gram conditions on separated windows have a converse forcing GRH
without the full positivity criterion of the explicit formula?
\item \emph{Prime-square bias against the zero sum.} Which other self-dual Artin or
modular $L$-functions have a prime-square bias large enough to dominate their zero sum?
Inequality \eqref{zer:eq:liminf} and the curvature bound of \cref{zer:prop:eta} give
verifiable sufficient conditions.
\item \emph{Zhao's classification beyond abelian fields.} Which nonabelian number fields
satisfy $B_K<2$? By \cref{zer:prop:frontier} the search is finite, subject to the Riemann
hypothesis for each candidate. What is the eventual sign of the integrated Mertens error
for fields with $B_K\ge2$?
\end{enumerate}
